\documentclass[a4paper,10pt]{article}
\usepackage[utf8]{inputenc}
\usepackage[T1]{fontenc}
\usepackage[ngerman]{babel}
\usepackage{lmodern}
\usepackage[left=14mm,right=12mm,top=13mm,bottom=11mm,headsep=4mm]{geometry}
\usepackage{graphicx}
\usepackage{tikz}
\usepackage{fancyhdr}
\usepackage{xcolor}
\usepackage{tabularx}
\usepackage{array}
\usepackage{ragged2e}
\usepackage{needspace}
\usepackage[most]{tcolorbox}
\usepackage{hyperref}
\hypersetup{colorlinks=true, linkcolor=black, pdfborder={0 0 0}}

\definecolor{A1Farbe}{HTML}{1F5C8B}
\definecolor{A2Farbe}{HTML}{116B54}
\definecolor{A3Farbe}{HTML}{8B4A0F}
\definecolor{A4Farbe}{HTML}{6A2C70}
\definecolor{A5Farbe}{HTML}{9B1B30}
\definecolor{Grau}{HTML}{555555}
\definecolor{Marker}{HTML}{FFD83D}
\definecolor{MarkerG}{HTML}{4FD08A}

% Kopfzeile ueber Marken: sonst zeigt die Seite den zuletzt gesetzten Wert
\newcommand{\marke}[1]{\markboth{#1}{#1}}
\pagestyle{fancy}
\fancyhf{}
\fancyhead[L]{\small\sffamily\bfseries\leftmark}
\fancyhead[R]{\small\sffamily\color{Grau}\thepage}
\renewcommand{\headrulewidth}{0.6pt}
\setlength{\parindent}{0pt}
\setlength{\fboxsep}{0pt}

% Kopf einer Altklausur-Aufgabe
\newcommand{\aufgabenkopf}[3]{%
  \par\vspace{1mm}
  \noindent\colorbox{#3}{\parbox{\dimexpr\textwidth-2\fboxsep}{%
    \color{white}\sffamily\bfseries\strut\hspace{1mm}#1\hfill#2\hspace{1mm}}}%
  \par\vspace{2.5mm}}

% Loesungsblock
\newtcolorbox{loesung}[1]{colback=#1!4, colframe=#1!45, boxrule=0pt,
  leftrule=1.8pt, arc=0.6mm, left=2.5mm, right=1mm, top=1.2mm, bottom=1.2mm,
  breakable, enhanced,
  overlay unbroken and first={\node[anchor=north west, font=\sffamily\bfseries
    \tiny, color=#1, inner sep=0] at ([shift={(2.5mm,-1.0mm)}]frame.north west)
    {LÖSUNG};}}

% Gesuchte Groesse im Register -- gleiches Gruen wie auf den Aufgabenseiten.
% \fboxsep ist global 0pt (fuer die Bild-Ausschnitte), hier lokal wieder Luft.
\newcommand{\gs}[1]{{\setlength{\fboxsep}{0.5mm}%
  \colorbox{MarkerG!45}{\strut #1}}}

% Teilaufgabe, die anderswo schon ausgeschrieben steht.
% #1 Farbe  #2 Nummer  #3 Kurztext der Frage  #4 Fundstelle (+ ggf. Zusatz)
\newcommand{\verweis}[4]{%
  \par\vspace{1.2mm}%
  \noindent\colorbox{#1!7}{\parbox{\dimexpr\textwidth-2\fboxsep}{%
    \vspace{0.6mm}\leftskip=2mm\rightskip=2mm
    {\small\sffamily\bfseries\color{#1}#2}\hspace{1.5mm}{\small #3}\par
    \vspace{0.5mm}%
    {\footnotesize\color{Grau}$\hookrightarrow$~#4}%
    \vspace{0.8mm}}}%
  \par\vspace{1.2mm}}

\begin{document}

\thispagestyle{empty}
\vspace*{14mm}
{\sffamily\bfseries\Huge Elektrodynamik}\par\vspace{3mm}
{\sffamily\LARGE\color{Grau} Altklausuren nach Aufgabentyp \textendash{} Aufgabe und Lösung ineinander}\par\vspace{4mm}
{\color{black}\rule{\textwidth}{1.2pt}}\par\vspace{7mm}
{\large\textbf{Übungs-Version} \textendash{} alle Altklausuren
außer den drei für Probeklausuren reservierten Jahrgängen (SoSe 2024,
WiSe 2024/25, WiSe 2025/26). Spoilerfrei zum Üben.}\par\vspace{6mm}

\textbf{Aufbau.} Jede Teilaufgabe steht direkt über ihrer Lösung (grau
hinterlegt) \textendash{} kein Blättern zwischen Aufgabe und Lösungsvorschlag.
Der Leerraum zum Rechnen ist herausgeschnitten, Skizzen sind verkleinert,
In der Aufgabenstellung ist \colorbox{Marker!45}{gegeben} gelb und
\colorbox{MarkerG!45}{gesucht} grün markiert.
Alle Aufgaben- und Lösungstexte sind unveränderte Ausschnitte aus den
Original-PDFs \textendash{} es wurde nichts abgetippt.\par\vspace{3mm}

\textbf{Keine Klausur fehlt, keine Frage steht doppelt.} Aussortiert wird
nicht klausur-, sondern \emph{teilaufgabenweise} \textendash{} und nur dort,
wo die Rechnung wirklich dieselbe ist und bloß die Zahlenwerte andere sind.
Dann steht statt Aufgabe und Lösung ein farbig hinterlegter \textbf{Verweis}
mit der Frage in einer Zeile und der Seite, auf der sie ausgerechnet ist.
\textbf{Sobald sich sonst etwas ändert} \textendash{} anderes Bauteil,
Serie statt parallel, gefüllt statt luftgefüllt, andere gesuchte Größe
\textendash{} ist es eine \emph{eigene} Aufgabe und steht vollständig da.
So sieht man an jeder Aufgabe, was neu daran ist.
\par\vspace{3mm}

Pro Aufgabentyp zuerst ein \textbf{Register} (\glqq wenn gefragt wird \ldots\
$\rightarrow$ Seite\grqq{}; fett die ausgeschriebene Fundstelle, dahinter die
Klausuren, die dasselbe fragen), dann ein Kasten mit möglichen \textbf{neuen}
Teilaufgaben samt Rechenweg-Skizze (eigene Prognose, keine offizielle
Ankündigung), dann die Aufgaben in chronologischer Reihenfolge.
\par\vspace{7mm}

{\sffamily\bfseries\large Inhalt}\par\vspace{2mm}
\begin{tabular}{@{}p{75mm}l@{}}
\textbf{A1 -- Koaxialleitung} & Seite \pageref{typ:1}\\[0.8mm]
\textbf{A2 -- Rechteckhohlleiter} & Seite \pageref{typ:2}\\[0.8mm]
\textbf{A3 -- Grenzflächen / Optik} & Seite \pageref{typ:3}\\[0.8mm]
\textbf{A4 -- Leitungstransformation} & Seite \pageref{typ:4}\\[0.8mm]
\textbf{A5 -- Antennen} & Seite \pageref{typ:5}\\[0.8mm]
\end{tabular}\par\vspace{12mm}
{\color{Grau}\small Alle 11 Klausuren stammen von Prof.\ Janker (90 min,
Open Book). Der nächste Anlauf im WiSe 2026/27 hat einen \emph{neuen} Prüfer
\textendash{} das Fünf-Typen-Schema ist damit nicht mehr garantiert, der Stoff
schon. Quelle: \texttt{ED/Klausur/Aufgaben} und \texttt{ED/Klausur/Lösung}.
Erzeugt mit \texttt{build\_kompakt.py}.}
\clearpage
\clearpage
\marke{\color{A1Farbe}A1 \textendash{} Koaxialleitung}
\phantomsection\label{typ:1}
\par\vspace*{-3mm}{\color{A1Farbe}\rule{\textwidth}{2.5pt}}\par\vspace{1.5mm}
{\sffamily\bfseries\huge\color{A1Farbe} A1 \textendash{} Koaxialleitung}\par\vspace{1mm}
{\sffamily\color{Grau} alle 8 Klausuren \textendash{} 23 verschiedene Teilaufgaben ausgeschrieben, 6 Wiederholungen als Verweis}\par\vspace{2mm}
{\color{A1Farbe}\rule{\textwidth}{0.8pt}}\par\vspace{4mm}
{\sffamily\bfseries Register: Wenn gefragt wird \ldots}\hfill{\footnotesize\color{Grau}\colorbox{MarkerG!45}{grün} = das ist gesucht \textbf{·} rechts steht, wo die Aufgabe ausgeschrieben ist}\par\vspace{1.5mm}
\renewcommand{\arraystretch}{1.12}\footnotesize
\begin{tabularx}{\textwidth}{@{}>{\RaggedRight}X>{\RaggedRight\scriptsize}p{27mm}@{}}
\hline\rule{0pt}{3.5mm}
Eindeutigkeit: welches \gs{$\varepsilon_r$} darf das Dielektrikum maximal haben? & \textbf{S.\,\pageref{tf:SoSe21:1:1}} \textcolor{Grau}{SoSe 21, 1.} \\
Eindeutigkeit: welchen \gs{Innendurchmesser $D$} darf der Außenleiter maximal haben? & \textbf{S.\,\pageref{tf:SoSe22:1:1}} \textcolor{Grau}{SoSe 22, 1.} \\
Eindeutigkeit: bis zu welcher \gs{Frequenz} ist nur ein Wellentyp ausbreitungsfähig? \newline \hspace*{3mm}\textcolor{Grau}{\itshape \textendash{} Grundfall} & \textbf{S.\,\pageref{tf:WiSe22x23:1:2}} \textcolor{Grau}{WiSe 22/23, 2.} \\
Eindeutigkeit: bis zu welcher \gs{Frequenz} ist nur ein Wellentyp ausbreitungsfähig? \newline \hspace*{3mm}\textcolor{Grau}{\itshape \textendash{} hier für \emph{beide} Leitungen des Triaxkabels einzeln, die kleinere Grenzfrequenz gewinnt} & \textbf{S.\,\pageref{tf:SoSe25:1:2}} \textcolor{Grau}{SoSe 25, 2.} \\
\gs{Innenleiter-Durchmesser $d$} für \emph{minimale Dämpfung} ($D/d=3{,}591$), dazu \gs{$Z_L$} \newline \hspace*{3mm}\textcolor{Grau}{\itshape \textendash{} Grundfall} & \textbf{S.\,\pageref{tf:SoSe21:1:2}} \textcolor{Grau}{SoSe 21, 2.} \\
\gs{Innenleiter-Durchmesser $d$} für \emph{minimale Dämpfung} ($D/d=3{,}591$), dazu \gs{$Z_L$} \newline \hspace*{3mm}\textcolor{Grau}{\itshape \textendash{} hier ohne $Z_L$} & \textbf{S.\,\pageref{tf:WiSe23x24:1:2}} \textcolor{Grau}{WiSe 23/24, 2.} \\
\gs{$\varepsilon_r$} aus vorgegebenem $Z_L$ bei bekannten $D$ und $d$ & \textbf{S.\,\pageref{tf:WiSe20x21:1:1}} \textcolor{Grau}{WiSe 20/21, 1.} \\
\gs{$d$} aus vorgegebenem $Z_L$ (bei bekanntem $D$) & \textbf{S.\,\pageref{tf:SoSe22:1:2}} \textcolor{Grau}{SoSe 22, 2.} \\
\gs{$Z_L$} vorwärts aus $D$, $d$ und $\varepsilon_r$ & \textbf{S.\,\pageref{tf:WiSe23x24:1:3}} \textcolor{Grau}{WiSe 23/24, 3.} \\
Triaxkabel: \gs{beide Durchmesser} so wählen, dass \emph{beide} Leitungen dasselbe $Z_L$ haben & \textbf{S.\,\pageref{tf:SoSe25:1:1}} \textcolor{Grau}{SoSe 25, 1.} \\
\gs{Dämpfung in dB} allein aus den Leiterverlusten (Skineffekt) & \textbf{S.\,\pageref{tf:WiSe20x21:1:2}} \textcolor{Grau}{WiSe 20/21, 2.} \\
\gs{Gesamtdämpfung in dB}: Leiterverluste \emph{plus} dielektrische Verluste & \textbf{S.\,\pageref{tf:WiSe22x23:1:3}} \textcolor{Grau}{WiSe 22/23, 3.} \\
\gs{Welche von zwei Leitungen} dämpft stärker, und \gs{um wieviel dB}? & \textbf{S.\,\pageref{tf:SoSe25:1:3}} \textcolor{Grau}{SoSe 25, 3.} \\
\gs{$\tan\delta$} aus gemessener Gesamt- und Leiterdämpfung & \textbf{S.\,\pageref{tf:WiSe20x21:1:3}} \textcolor{Grau}{WiSe 20/21, 3.} \\
Maximaler \gs{$\tan\delta$}, damit die Gesamtdämpfung ein Budget einhält & \textbf{S.\,\pageref{tf:WiSe21x22:1:2}} \textcolor{Grau}{WiSe 21/22, 2.} \\
Bei welcher \gs{Frequenz} sind die dielektrischen Verluste ein bestimmter Anteil der Leiterverluste? & \textbf{S.\,\pageref{tf:SoSe22:1:3}} \textcolor{Grau}{SoSe 22, 3.} \\
Schirmdämpfung: in welchem \gs{Frequenzbereich} wird eine Vorgabe eingehalten? \newline \hspace*{3mm}\textcolor{Grau}{\itshape \textendash{} Grundfall} & \textbf{S.\,\pageref{tf:SoSe21:1:3}} \textcolor{Grau}{SoSe 21, 3.} \\
Schirmdämpfung: in welchem \gs{Frequenzbereich} wird eine Vorgabe eingehalten? \newline \hspace*{3mm}\textcolor{Grau}{\itshape \textendash{} hier als Übersprechen zwischen den beiden Triax-Leitungen verpackt} & \textbf{S.\,\pageref{tf:SoSe25:1:4}} \textcolor{Grau}{SoSe 25, 4.} \\
\gs{Schirmdämpfung in dB} bei gegebener Wandstärke und Frequenz & \textbf{S.\,\pageref{tf:SoSe22:1:4}} \textcolor{Grau}{SoSe 22, 4.} \\
Schirmdämpfung: welche \gs{Wandstärke $w$} ist nötig? & \textbf{S.\,\pageref{tf:WiSe22x23:1:4}} \textcolor{Grau}{WiSe 22/23, 4.} \\
Welche \gs{Leistung} wird im Kabel in Wärme umgesetzt? & \textbf{S.\,\pageref{tf:SoSe23:1:3}} \textcolor{Grau}{SoSe 23, 3.} \\
Koaxialresonator: \gs{niedrigste Resonanzfrequenz} & \textbf{S.\,\pageref{tf:SoSe23:1:4}} \textcolor{Grau}{SoSe 23, 4.} \\
\gs{Zusatzverluste} durch Fehlanpassung (Sender $\neq Z_L$) & \textbf{S.\,\pageref{tf:WiSe23x24:1:5}} \textcolor{Grau}{WiSe 23/24, 5.} \\
\hline\end{tabularx}\normalsize
\par\vspace{5mm}
\begin{tcolorbox}[colback=A1Farbe!4, colframe=A1Farbe!55, boxrule=0.5pt, arc=1mm, left=3mm, right=2mm, top=2mm, bottom=2mm, breakable, title={\sffamily\bfseries\small Was diesmal neu drankommen könnte \normalfont\itshape\footnotesize -- eigene Prognose (keine offizielle Ankündigung): Umkehrungen, Kombinationen und Transfers zwischen den Aufgabentypen}, coltitle=white, colbacktitle=A1Farbe!70]
\sffamily\small
\begin{itemize}\setlength{\itemsep}{2.2mm}\setlength{\leftmargini}{4mm}
\item Welche Leistung kann das Koaxialkabel maximal übertragen, wenn die Feldstärke die Durchbruchfeldstärke nicht überschreiten darf? ($E$ ist am Innenleiter am größten) \textcolor{Grau}{\itshape(Diese Frage gab es bisher nur beim Hohlleiter)}\\[0.6mm] {\footnotesize\color{A1Farbe}$\rightarrow$ $E$ ist am Innenleiter am größten: $E_{max}=\frac{U}{(d/2)\ln(D/d)}$ $\to$ daraus $\hat{U}$, dann $P=\hat{U}^2/(2Z_L)$. FS 7.1}
\item Wie lange braucht ein Signal / Impuls durch das $l$ lange Kabel? ($v_{ph}=c_0/\sqrt{\varepsilon_r}$)\\[0.6mm] {\footnotesize\color{A1Farbe}$\rightarrow$ $t=l/v_{ph}$ mit $v_{ph}=c_0/\sqrt{\varepsilon_r}$ -- beim Koax (TEM-Welle) keine Dispersion, also \emph{keine} Gruppengeschwindigkeit nötig}
\item Koaxialresonator: Resonanzfrequenz der \emph{n-ten} Ordnung oder bei dielektrisch gefülltem Kabel \textcolor{Grau}{\itshape(bisher nur die niedrigste: SoSe 23 A1.4)}\\[0.6mm] {\footnotesize\color{A1Farbe}$\rightarrow$ $l=n\cdot\lambda/2 \Rightarrow f_n=n\,c_0/(2l\sqrt{\varepsilon_r})$ -- dieselbe Rechnung wie SoSe 23, nur mit $n>1$}
\item Umkehrung: Welche Leistung darf eingespeist werden, damit die Verlustwärme einen vorgegebenen Wert nicht überschreitet? \textcolor{Grau}{\itshape(Hinrichtung: SoSe 23 A1.3)}\\[0.6mm] {\footnotesize\color{A1Farbe}$\rightarrow$ $a_{dB}=8{,}686\,\alpha l$ und $P_{Verlust}=P_{ein}(1-10^{-a_{dB}/10})$ nach $P_{ein}$ auflösen. K1}
\item Ab welcher Frequenz ist der Außenleiter \glqq ausreichend dick\grqq, d.\,h.\ Wandstärke $w \ge 3\delta$ (Skintiefe)?\\[0.6mm] {\footnotesize\color{A1Farbe}$\rightarrow$ $\delta=1/\sqrt{\pi f\mu\kappa}$, Bedingung $w\ge 3\delta$ nach $f$ auflösen: $f\ge 9/(\pi\mu\kappa w^2)$. FS 5.2}
\item Um wieviel dB ändert sich die Dämpfung, wenn statt Kupfer Silber oder Messing verwendet wird? ($\alpha \sim 1/\sqrt{\kappa}$)\\[0.6mm] {\footnotesize\color{A1Farbe}$\rightarrow$ $\alpha_L\sim 1/\sqrt{\kappa}$, also $a_{neu}=a_{alt}\sqrt{\kappa_{alt}/\kappa_{neu}}$; $\kappa$-Werte aus K0b}
\item Stehwellen anders gemessen: Abstand Minimum$\to$Minimum ($=\lambda/2$) oder Verhältnis $U_{max}/U_{min}$ direkt gegeben $\rightarrow$ $|r|$, $Z_V$\\[0.6mm] {\footnotesize\color{A1Farbe}$\rightarrow$ Min$\to$Min $=\lambda/2$ (Max$\to$Min $=\lambda/4$); $|r|=\frac{U_{max}-U_{min}}{U_{max}+U_{min}}$, $Z_V=Z_L\frac{1+r}{1-r}$. K1 rechte Spalte}
\end{itemize}\end{tcolorbox}
\clearpage
\Needspace*{237.9mm}
\marke{\color{A1Farbe}A1 \textendash{} WiSe 2020/21}
\phantomsection\label{auf:WiSe20x21:1}
\aufgabenkopf{WiSe 2020/21 \textendash{} Aufgabe 1}{13 Punkte}{A1Farbe}
\begin{tikzpicture}[inner sep=0]
\node[anchor=north west,inner sep=0] (bi) at (0,0) {\includegraphics[page=2,trim={52.00bp 670.50bp 39.00bp 106.50bp},clip,scale=0.840]{A_WiSe20x21.pdf}};
\fill[Marker,opacity=0.30,rounded corners=0.6bp] (99.12bp,-1.26bp) rectangle (140.62bp,-11.34bp);
\fill[Marker,opacity=0.30,rounded corners=0.6bp] (288.96bp,-1.26bp) rectangle (317.15bp,-11.34bp);
\fill[Marker,opacity=0.30,rounded corners=0.6bp] (84.84bp,-22.26bp) rectangle (120.30bp,-32.34bp);
\fill[Marker,opacity=0.30,rounded corners=0.6bp] (270.48bp,-22.26bp) rectangle (309.07bp,-32.34bp);
\fill[Marker,opacity=0.30,rounded corners=0.6bp] (12.60bp,-44.94bp) rectangle (43.68bp,-55.02bp);
\end{tikzpicture}
\par\vspace{1.2mm}
\includegraphics[page=2,trim={52.00bp 472.50bp 39.00bp 205.50bp},clip,scale=0.620]{A_WiSe20x21.pdf}
\par\vspace{1.2mm}
\phantomsection\label{tf:WiSe20x21:1:1}
\par\vspace{1.6mm}
\Needspace*{39mm}
\begin{tikzpicture}[inner sep=0]
\node[anchor=north west,inner sep=0] (bi) at (0,0) {\includegraphics[page=2,trim={52.00bp 377.50bp 39.00bp 437.50bp},clip,scale=0.840]{A_WiSe20x21.pdf}};
\fill[Marker,opacity=0.30,rounded corners=0.6bp] (237.72bp,-1.26bp) rectangle (262.92bp,-11.34bp);
\fill[MarkerG,opacity=0.30,rounded corners=0.6bp] (23.52bp,-12.18bp) rectangle (111.48bp,-22.26bp);
\end{tikzpicture}
\par\nopagebreak
\vspace{0.8mm}\nopagebreak
\begin{loesung}{A1Farbe}
\vspace{1.2mm}
\includegraphics[page=1,trim={52.00bp 609.50bp 39.00bp 144.50bp},clip,scale=0.700]{L_WiSe20x21.pdf}
\par
\end{loesung}
\par\vspace{1.6mm}
\begin{tikzpicture}[inner sep=0]
\node[anchor=north west,inner sep=0] (bi) at (0,0) {\includegraphics[page=2,trim={52.00bp 314.50bp 39.00bp 514.50bp},clip,scale=0.840]{A_WiSe20x21.pdf}};
\fill[Marker,opacity=0.30,rounded corners=0.6bp] (68.04bp,-1.26bp) rectangle (96.81bp,-11.34bp);
\end{tikzpicture}
\par
\phantomsection\label{tf:WiSe20x21:1:2}
\par\vspace{1.6mm}
\Needspace*{43mm}
\begin{tikzpicture}[inner sep=0]
\node[anchor=north west,inner sep=0] (bi) at (0,0) {\includegraphics[page=2,trim={52.00bp 275.50bp 39.00bp 552.50bp},clip,scale=0.840]{A_WiSe20x21.pdf}};
\fill[MarkerG,opacity=0.30,rounded corners=0.6bp] (87.95bp,-1.26bp) rectangle (119.94bp,-11.34bp);
\end{tikzpicture}
\par\nopagebreak
\vspace{0.8mm}\nopagebreak
\begin{loesung}{A1Farbe}
\vspace{1.2mm}
\includegraphics[page=1,trim={52.00bp 451.50bp 39.00bp 271.50bp},clip,scale=0.700]{L_WiSe20x21.pdf}
\par
\end{loesung}
\phantomsection\label{tf:WiSe20x21:1:3}
\par\vspace{1.6mm}
\Needspace*{71mm}
\begin{tikzpicture}[inner sep=0]
\node[anchor=north west,inner sep=0] (bi) at (0,0) {\includegraphics[page=2,trim={52.00bp 198.50bp 39.00bp 590.50bp},clip,scale=0.840]{A_WiSe20x21.pdf}};
\fill[Marker,opacity=0.30,rounded corners=0.6bp] (231.00bp,-1.26bp) rectangle (267.05bp,-11.34bp);
\fill[Marker,opacity=0.30,rounded corners=0.6bp] (370.44bp,-13.02bp) rectangle (403.11bp,-23.10bp);
\fill[MarkerG,opacity=0.30,rounded corners=0.6bp] (387.45bp,-1.26bp) rectangle (417.61bp,-11.34bp);
\end{tikzpicture}
\par\nopagebreak
\vspace{0.8mm}\nopagebreak
\begin{loesung}{A1Farbe}
\vspace{1.2mm}
\includegraphics[page=1,trim={52.00bp 203.50bp 39.00bp 450.50bp},clip,scale=0.700]{L_WiSe20x21.pdf}
\par
\end{loesung}
\par\vspace{6mm}
\Needspace*{209.7mm}
\marke{\color{A1Farbe}A1 \textendash{} SoSe 2021}
\phantomsection\label{auf:SoSe21:1}
\aufgabenkopf{SoSe 2021 \textendash{} Aufgabe 1}{14 Punkte}{A1Farbe}
\begin{tikzpicture}[inner sep=0]
\node[anchor=north west,inner sep=0] (bi) at (0,0) {\includegraphics[page=2,trim={52.00bp 682.50bp 39.00bp 106.50bp},clip,scale=0.840]{A_SoSe21.pdf}};
\fill[Marker,opacity=0.30,rounded corners=0.6bp] (15.85bp,-12.18bp) rectangle (42.07bp,-22.26bp);
\fill[Marker,opacity=0.30,rounded corners=0.6bp] (105.00bp,-12.18bp) rectangle (133.06bp,-22.26bp);
\fill[Marker,opacity=0.30,rounded corners=0.6bp] (286.87bp,-12.18bp) rectangle (318.88bp,-22.26bp);
\fill[Marker,opacity=0.30,rounded corners=0.6bp] (193.20bp,-22.26bp) rectangle (238.71bp,-32.34bp);
\fill[Marker,opacity=0.30,rounded corners=0.6bp] (24.50bp,-34.02bp) rectangle (43.68bp,-44.10bp);
\end{tikzpicture}
\par\vspace{1.2mm}
\includegraphics[page=2,trim={52.00bp 484.50bp 39.00bp 192.50bp},clip,scale=0.620]{A_SoSe21.pdf}
\par\vspace{1.2mm}
\phantomsection\label{tf:SoSe21:1:1}
\par\vspace{1.6mm}
\Needspace*{39mm}
\begin{tikzpicture}[inner sep=0]
\node[anchor=north west,inner sep=0] (bi) at (0,0) {\includegraphics[page=2,trim={52.00bp 391.50bp 39.00bp 424.50bp},clip,scale=0.840]{A_SoSe21.pdf}};
\fill[MarkerG,opacity=0.30,rounded corners=0.6bp] (49.65bp,-1.26bp) rectangle (122.27bp,-11.34bp);
\end{tikzpicture}
\par\nopagebreak
\vspace{0.8mm}\nopagebreak
\begin{loesung}{A1Farbe}
\vspace{1.2mm}
\includegraphics[page=1,trim={52.00bp 606.50bp 39.00bp 144.50bp},clip,scale=0.700]{L_SoSe21.pdf}
\par
\end{loesung}
\par\vspace{1.6mm}
\begin{tikzpicture}[inner sep=0]
\node[anchor=north west,inner sep=0] (bi) at (0,0) {\includegraphics[page=2,trim={52.00bp 340.50bp 39.00bp 488.50bp},clip,scale=0.840]{A_SoSe21.pdf}};
\fill[Marker,opacity=0.30,rounded corners=0.6bp] (37.80bp,-1.26bp) rectangle (66.57bp,-11.34bp);
\end{tikzpicture}
\par
\phantomsection\label{tf:SoSe21:1:2}
\par\vspace{1.6mm}
\Needspace*{37mm}
\begin{tikzpicture}[inner sep=0]
\node[anchor=north west,inner sep=0] (bi) at (0,0) {\includegraphics[page=2,trim={52.00bp 301.50bp 39.00bp 514.50bp},clip,scale=0.840]{A_SoSe21.pdf}};
\fill[MarkerG,opacity=0.30,rounded corners=0.6bp] (68.76bp,-0.42bp) rectangle (111.90bp,-10.50bp);
\end{tikzpicture}
\par\nopagebreak
\vspace{0.8mm}\nopagebreak
\begin{loesung}{A1Farbe}
\vspace{1.2mm}
\includegraphics[page=1,trim={52.00bp 464.50bp 39.00bp 296.50bp},clip,scale=0.700]{L_SoSe21.pdf}
\par
\end{loesung}
\phantomsection\label{tf:SoSe21:1:3}
\par\vspace{1.6mm}
\Needspace*{50mm}
\begin{tikzpicture}[inner sep=0]
\node[anchor=north west,inner sep=0] (bi) at (0,0) {\includegraphics[page=2,trim={52.00bp 238.50bp 39.00bp 589.50bp},clip,scale=0.840]{A_SoSe21.pdf}};
\fill[MarkerG,opacity=0.30,rounded corners=0.6bp] (67.48bp,-1.26bp) rectangle (129.11bp,-11.34bp);
\end{tikzpicture}
\par\nopagebreak
\vspace{0.8mm}\nopagebreak
\begin{loesung}{A1Farbe}
\vspace{1.2mm}
\includegraphics[page=1,trim={52.00bp 275.50bp 39.00bp 432.50bp},clip,scale=0.700]{L_SoSe21.pdf}
\par
\par\vspace{1.0mm}
\includegraphics[page=1,trim={52.00bp 243.50bp 39.00bp 584.50bp},clip,scale=0.700]{L_SoSe21.pdf}
\par
\end{loesung}
\par\vspace{6mm}
\Needspace*{176.2mm}
\marke{\color{A1Farbe}A1 \textendash{} WiSe 2021/22}
\phantomsection\label{auf:WiSe21x22:1}
\aufgabenkopf{WiSe 2021/22 \textendash{} Aufgabe 1}{14 Punkte}{A1Farbe}
\begin{tikzpicture}[inner sep=0]
\node[anchor=north west,inner sep=0] (bi) at (0,0) {\includegraphics[page=2,trim={52.00bp 669.50bp 39.00bp 106.50bp},clip,scale=0.840]{A_WiSe21x22.pdf}};
\fill[Marker,opacity=0.30,rounded corners=0.6bp] (308.28bp,-12.18bp) rectangle (348.24bp,-22.26bp);
\fill[Marker,opacity=0.30,rounded corners=0.6bp] (5.04bp,-23.10bp) rectangle (37.25bp,-33.18bp);
\fill[Marker,opacity=0.30,rounded corners=0.6bp] (232.68bp,-23.10bp) rectangle (267.40bp,-33.18bp);
\fill[Marker,opacity=0.30,rounded corners=0.6bp] (54.09bp,-34.02bp) rectangle (85.28bp,-44.10bp);
\fill[Marker,opacity=0.30,rounded corners=0.6bp] (24.50bp,-44.94bp) rectangle (43.68bp,-55.02bp);
\end{tikzpicture}
\par\vspace{1.2mm}
\includegraphics[page=2,trim={52.00bp 471.50bp 39.00bp 205.50bp},clip,scale=0.620]{A_WiSe21x22.pdf}
\par\vspace{1.2mm}
\phantomsection\label{tf:WiSe21x22:1:1}
\Needspace*{18mm}
\verweis{A1Farbe}{1.}{\gs{Innenleiter-Durchmesser $d$} für \emph{minimale Dämpfung} ($D/d=3{,}591$), dazu \gs{$Z_L$}}{gleiche Frage, nur andere Zahlenwerte \textendash{} siehe \textbf{SoSe 2021, Aufgabe 1.2}, \textbf{Seite \pageref{tf:SoSe21:1:2}}}
\phantomsection\label{tf:WiSe21x22:1:2}
\par\vspace{1.6mm}
\Needspace*{92mm}
\begin{tikzpicture}[inner sep=0]
\node[anchor=north west,inner sep=0] (bi) at (0,0) {\includegraphics[page=2,trim={52.00bp 299.50bp 39.00bp 501.50bp},clip,scale=0.840]{A_WiSe21x22.pdf}};
\fill[Marker,opacity=0.30,rounded corners=0.6bp] (193.20bp,-12.18bp) rectangle (226.91bp,-22.26bp);
\fill[MarkerG,opacity=0.30,rounded corners=0.6bp] (96.10bp,-1.26bp) rectangle (150.19bp,-11.34bp);
\end{tikzpicture}
\par\nopagebreak
\vspace{0.8mm}\nopagebreak
\begin{loesung}{A1Farbe}
\vspace{1.2mm}
\includegraphics[page=1,trim={52.00bp 463.50bp 39.00bp 270.50bp},clip,scale=0.700]{L_WiSe21x22.pdf}
\par
\par\vspace{1.0mm}
\includegraphics[page=1,trim={52.00bp 261.50bp 39.00bp 403.50bp},clip,scale=0.700]{L_WiSe21x22.pdf}
\par
\end{loesung}
\par\vspace{6mm}
\Needspace*{264.8mm}
\marke{\color{A1Farbe}A1 \textendash{} SoSe 2022}
\phantomsection\label{auf:SoSe22:1}
\aufgabenkopf{SoSe 2022 \textendash{} Aufgabe 1}{18 Punkte}{A1Farbe}
\begin{tikzpicture}[inner sep=0]
\node[anchor=north west,inner sep=0] (bi) at (0,0) {\includegraphics[page=2,trim={52.00bp 694.50bp 39.00bp 106.50bp},clip,scale=0.840]{A_SoSe22.pdf}};
\fill[Marker,opacity=0.30,rounded corners=0.6bp] (101.64bp,-2.10bp) rectangle (129.78bp,-12.18bp);
\fill[Marker,opacity=0.30,rounded corners=0.6bp] (147.84bp,-2.10bp) rectangle (173.88bp,-12.18bp);
\fill[Marker,opacity=0.30,rounded corners=0.6bp] (5.04bp,-13.02bp) rectangle (43.88bp,-23.10bp);
\fill[Marker,opacity=0.30,rounded corners=0.6bp] (269.64bp,-13.02bp) rectangle (313.85bp,-23.10bp);
\fill[Marker,opacity=0.30,rounded corners=0.6bp] (24.50bp,-23.94bp) rectangle (43.68bp,-34.02bp);
\end{tikzpicture}
\par\vspace{1.2mm}
\includegraphics[page=2,trim={52.00bp 496.50bp 39.00bp 180.50bp},clip,scale=0.620]{A_SoSe22.pdf}
\par\vspace{1.2mm}
\phantomsection\label{tf:SoSe22:1:1}
\par\vspace{1.6mm}
\Needspace*{31mm}
\begin{tikzpicture}[inner sep=0]
\node[anchor=north west,inner sep=0] (bi) at (0,0) {\includegraphics[page=2,trim={52.00bp 416.50bp 39.00bp 412.50bp},clip,scale=0.840]{A_SoSe22.pdf}};
\fill[MarkerG,opacity=0.30,rounded corners=0.6bp] (94.92bp,-0.42bp) rectangle (168.00bp,-10.50bp);
\fill[MarkerG,opacity=0.30,rounded corners=0.6bp] (166.32bp,-0.42bp) rectangle (174.72bp,-10.50bp);
\end{tikzpicture}
\par\nopagebreak
\vspace{0.8mm}\nopagebreak
\begin{loesung}{A1Farbe}
\vspace{1.2mm}
\includegraphics[page=1,trim={52.00bp 657.50bp 39.00bp 144.50bp},clip,scale=0.700]{L_SoSe22.pdf}
\par
\par\vspace{1.0mm}
\includegraphics[page=1,trim={52.00bp 614.50bp 39.00bp 194.50bp},clip,scale=0.700]{L_SoSe22.pdf}
\par
\end{loesung}
\par\vspace{1.6mm}
\begin{tikzpicture}[inner sep=0]
\node[anchor=north west,inner sep=0] (bi) at (0,0) {\includegraphics[page=2,trim={52.00bp 367.50bp 39.00bp 462.50bp},clip,scale=0.840]{A_SoSe22.pdf}};
\fill[Marker,opacity=0.30,rounded corners=0.6bp] (37.80bp,-1.26bp) rectangle (76.44bp,-11.34bp);
\end{tikzpicture}
\par
\phantomsection\label{tf:SoSe22:1:2}
\par\vspace{1.6mm}
\Needspace*{36mm}
\begin{tikzpicture}[inner sep=0]
\node[anchor=north west,inner sep=0] (bi) at (0,0) {\includegraphics[page=2,trim={52.00bp 326.50bp 39.00bp 488.50bp},clip,scale=0.840]{A_SoSe22.pdf}};
\fill[Marker,opacity=0.30,rounded corners=0.6bp] (386.06bp,-1.26bp) rectangle (403.20bp,-11.34bp);
\fill[MarkerG,opacity=0.30,rounded corners=0.6bp] (68.76bp,-1.26bp) rectangle (111.90bp,-11.34bp);
\end{tikzpicture}
\par\nopagebreak
\vspace{0.8mm}\nopagebreak
\begin{loesung}{A1Farbe}
\vspace{1.2mm}
\includegraphics[page=1,trim={52.00bp 494.50bp 39.00bp 271.50bp},clip,scale=0.700]{L_SoSe22.pdf}
\par
\end{loesung}
\phantomsection\label{tf:SoSe22:1:3}
\par\vspace{1.6mm}
\Needspace*{72mm}
\begin{tikzpicture}[inner sep=0]
\node[anchor=north west,inner sep=0] (bi) at (0,0) {\includegraphics[page=2,trim={52.00bp 237.50bp 39.00bp 564.50bp},clip,scale=0.840]{A_SoSe22.pdf}};
\fill[Marker,opacity=0.30,rounded corners=0.6bp] (155.40bp,-1.26bp) rectangle (193.91bp,-11.34bp);
\fill[MarkerG,opacity=0.30,rounded corners=0.6bp] (69.02bp,-1.26bp) rectangle (101.04bp,-11.34bp);
\end{tikzpicture}
\par\nopagebreak
\vspace{0.8mm}\nopagebreak
\begin{loesung}{A1Farbe}
\vspace{1.2mm}
\includegraphics[page=1,trim={52.00bp 253.50bp 39.00bp 382.50bp},clip,scale=0.700]{L_SoSe22.pdf}
\par
\end{loesung}
\phantomsection\label{tf:SoSe22:1:4}
\par\vspace{1.6mm}
\Needspace*{40mm}
\begin{tikzpicture}[inner sep=0]
\node[anchor=north west,inner sep=0] (bi) at (0,0) {\includegraphics[page=2,trim={52.00bp 187.50bp 39.00bp 641.50bp},clip,scale=0.840]{A_SoSe22.pdf}};
\fill[Marker,opacity=0.30,rounded corners=0.6bp] (290.98bp,-1.26bp) rectangle (316.98bp,-11.34bp);
\fill[MarkerG,opacity=0.30,rounded corners=0.6bp] (51.89bp,-1.26bp) rectangle (110.30bp,-11.34bp);
\end{tikzpicture}
\par\nopagebreak
\vspace{0.8mm}\nopagebreak
\begin{loesung}{A1Farbe}
\vspace{1.2mm}
\includegraphics[page=1,trim={52.00bp 220.50bp 39.00bp 610.50bp},clip,scale=0.700]{L_SoSe22.pdf}
\par
\par\vspace{1.0mm}
\includegraphics[page=1,trim={52.00bp 108.50bp 39.00bp 634.50bp},clip,scale=0.700]{L_SoSe22.pdf}
\par
\end{loesung}
\par\vspace{6mm}
\Needspace*{230.6mm}
\marke{\color{A1Farbe}A1 \textendash{} WiSe 2022/23}
\phantomsection\label{auf:WiSe22x23:1}
\aufgabenkopf{WiSe 2022/23 \textendash{} Aufgabe 1}{18 Punkte}{A1Farbe}
\begin{tikzpicture}[inner sep=0]
\node[anchor=north west,inner sep=0] (bi) at (0,0) {\includegraphics[page=2,trim={52.00bp 682.50bp 39.00bp 106.50bp},clip,scale=0.840]{A_WiSe22x23.pdf}};
\fill[Marker,opacity=0.30,rounded corners=0.6bp] (113.40bp,-2.10bp) rectangle (144.97bp,-12.18bp);
\fill[Marker,opacity=0.30,rounded corners=0.6bp] (272.16bp,-2.10bp) rectangle (301.14bp,-12.18bp);
\fill[Marker,opacity=0.30,rounded corners=0.6bp] (318.36bp,-2.10bp) rectangle (346.92bp,-12.18bp);
\fill[Marker,opacity=0.30,rounded corners=0.6bp] (191.52bp,-13.02bp) rectangle (236.08bp,-23.10bp);
\fill[Marker,opacity=0.30,rounded corners=0.6bp] (84.84bp,-23.10bp) rectangle (126.34bp,-33.18bp);
\fill[Marker,opacity=0.30,rounded corners=0.6bp] (12.60bp,-34.86bp) rectangle (43.68bp,-44.94bp);
\end{tikzpicture}
\par\vspace{1.2mm}
\includegraphics[page=2,trim={52.00bp 484.50bp 39.00bp 193.50bp},clip,scale=0.620]{A_WiSe22x23.pdf}
\par\vspace{1.2mm}
\phantomsection\label{tf:WiSe22x23:1:1}
\Needspace*{18mm}
\verweis{A1Farbe}{1.}{\gs{Innenleiter-Durchmesser $d$} für \emph{minimale Dämpfung} ($D/d=3{,}591$), dazu \gs{$Z_L$}}{gleiche Frage, nur andere Zahlenwerte \textendash{} siehe \textbf{SoSe 2021, Aufgabe 1.2}, \textbf{Seite \pageref{tf:SoSe21:1:2}}}
\phantomsection\label{tf:WiSe22x23:1:2}
\par\vspace{1.6mm}
\Needspace*{26mm}
\begin{tikzpicture}[inner sep=0]
\node[anchor=north west,inner sep=0] (bi) at (0,0) {\includegraphics[page=2,trim={52.00bp 327.50bp 39.00bp 488.50bp},clip,scale=0.840]{A_WiSe22x23.pdf}};
\fill[MarkerG,opacity=0.30,rounded corners=0.6bp] (122.46bp,-0.42bp) rectangle (155.80bp,-10.50bp);
\end{tikzpicture}
\par\nopagebreak
\vspace{0.8mm}\nopagebreak
\begin{loesung}{A1Farbe}
\vspace{1.2mm}
\includegraphics[page=1,trim={52.00bp 532.50bp 39.00bp 271.50bp},clip,scale=0.700]{L_WiSe22x23.pdf}
\par
\end{loesung}
\par\vspace{1.6mm}
\begin{tikzpicture}[inner sep=0]
\node[anchor=north west,inner sep=0] (bi) at (0,0) {\includegraphics[page=2,trim={52.00bp 290.50bp 39.00bp 538.50bp},clip,scale=0.840]{A_WiSe22x23.pdf}};
\fill[Marker,opacity=0.30,rounded corners=0.6bp] (172.87bp,-1.26bp) rectangle (202.78bp,-11.34bp);
\end{tikzpicture}
\par
\phantomsection\label{tf:WiSe22x23:1:3}
\par\vspace{1.6mm}
\Needspace*{64mm}
\begin{tikzpicture}[inner sep=0]
\node[anchor=north west,inner sep=0] (bi) at (0,0) {\includegraphics[page=2,trim={52.00bp 251.50bp 39.00bp 576.50bp},clip,scale=0.840]{A_WiSe22x23.pdf}};
\fill[MarkerG,opacity=0.30,rounded corners=0.6bp] (50.96bp,-1.26bp) rectangle (115.36bp,-11.34bp);
\end{tikzpicture}
\par\nopagebreak
\vspace{0.8mm}\nopagebreak
\begin{loesung}{A1Farbe}
\vspace{1.2mm}
\includegraphics[page=1,trim={52.00bp 268.50bp 39.00bp 369.50bp},clip,scale=0.700]{L_WiSe22x23.pdf}
\par
\end{loesung}
\phantomsection\label{tf:WiSe22x23:1:4}
\par\vspace{1.6mm}
\Needspace*{45mm}
\begin{tikzpicture}[inner sep=0]
\node[anchor=north west,inner sep=0] (bi) at (0,0) {\includegraphics[page=2,trim={52.00bp 188.50bp 39.00bp 627.50bp},clip,scale=0.840]{A_WiSe22x23.pdf}};
\fill[Marker,opacity=0.30,rounded corners=0.6bp] (307.36bp,-0.42bp) rectangle (324.60bp,-10.50bp);
\fill[Marker,opacity=0.30,rounded corners=0.6bp] (22.68bp,-11.34bp) rectangle (63.97bp,-21.42bp);
\fill[MarkerG,opacity=0.30,rounded corners=0.6bp] (89.62bp,-0.42bp) rectangle (130.18bp,-10.50bp);
\end{tikzpicture}
\par\nopagebreak
\vspace{0.8mm}\nopagebreak
\begin{loesung}{A1Farbe}
\vspace{1.2mm}
\includegraphics[page=1,trim={52.00bp 232.50bp 39.00bp 598.50bp},clip,scale=0.700]{L_WiSe22x23.pdf}
\par
\par\vspace{1.0mm}
\includegraphics[page=1,trim={52.00bp 115.50bp 39.00bp 622.50bp},clip,scale=0.700]{L_WiSe22x23.pdf}
\par
\end{loesung}
\par\vspace{6mm}
\Needspace*{190.7mm}
\marke{\color{A1Farbe}A1 \textendash{} SoSe 2023}
\phantomsection\label{auf:SoSe23:1}
\aufgabenkopf{SoSe 2023 \textendash{} Aufgabe 1}{16 Punkte}{A1Farbe}
\begin{tikzpicture}[inner sep=0]
\node[anchor=north west,inner sep=0] (bi) at (0,0) {\includegraphics[page=2,trim={52.00bp 695.50bp 39.00bp 106.50bp},clip,scale=0.840]{A_SoSe23.pdf}};
\fill[Marker,opacity=0.30,rounded corners=0.6bp] (160.44bp,-1.26bp) rectangle (192.80bp,-11.34bp);
\fill[Marker,opacity=0.30,rounded corners=0.6bp] (54.04bp,-12.18bp) rectangle (88.48bp,-22.26bp);
\fill[Marker,opacity=0.30,rounded corners=0.6bp] (24.50bp,-23.10bp) rectangle (43.68bp,-33.18bp);
\end{tikzpicture}
\par\vspace{1.2mm}
\includegraphics[page=2,trim={52.00bp 497.50bp 39.00bp 179.50bp},clip,scale=0.620]{A_SoSe23.pdf}
\par\vspace{1.2mm}
\phantomsection\label{tf:SoSe23:1:1}
\Needspace*{18mm}
\verweis{A1Farbe}{1.}{Eindeutigkeit: welchen \gs{Innendurchmesser $D$} darf der Außenleiter maximal haben?}{gleiche Frage, nur andere Zahlenwerte \textendash{} siehe \textbf{SoSe 2022, Aufgabe 1.1}, \textbf{Seite \pageref{tf:SoSe22:1:1}}}
\par\vspace{1.6mm}
\begin{tikzpicture}[inner sep=0]
\node[anchor=north west,inner sep=0] (bi) at (0,0) {\includegraphics[page=2,trim={52.00bp 353.50bp 39.00bp 474.50bp},clip,scale=0.840]{A_SoSe23.pdf}};
\fill[Marker,opacity=0.30,rounded corners=0.6bp] (209.16bp,-1.26bp) rectangle (247.11bp,-11.34bp);
\end{tikzpicture}
\par
\phantomsection\label{tf:SoSe23:1:2}
\Needspace*{18mm}
\verweis{A1Farbe}{2.}{\gs{$d$} aus vorgegebenem $Z_L$ (bei bekanntem $D$)}{gleiche Frage, nur andere Zahlenwerte \textendash{} siehe \textbf{SoSe 2022, Aufgabe 1.2}, \textbf{Seite \pageref{tf:SoSe22:1:2}}}
\par\vspace{1.6mm}
\begin{tikzpicture}[inner sep=0]
\node[anchor=north west,inner sep=0] (bi) at (0,0) {\includegraphics[page=2,trim={52.00bp 264.50bp 39.00bp 563.50bp},clip,scale=0.840]{A_SoSe23.pdf}};
\fill[Marker,opacity=0.30,rounded corners=0.6bp] (199.08bp,-1.26bp) rectangle (236.04bp,-11.34bp);
\end{tikzpicture}
\par
\phantomsection\label{tf:SoSe23:1:3}
\par\vspace{1.6mm}
\Needspace*{51mm}
\begin{tikzpicture}[inner sep=0]
\node[anchor=north west,inner sep=0] (bi) at (0,0) {\includegraphics[page=2,trim={52.00bp 214.50bp 39.00bp 589.50bp},clip,scale=0.840]{A_SoSe23.pdf}};
\fill[Marker,opacity=0.30,rounded corners=0.6bp] (141.12bp,-0.42bp) rectangle (177.38bp,-10.50bp);
\fill[MarkerG,opacity=0.30,rounded corners=0.6bp] (263.25bp,-0.42bp) rectangle (296.77bp,-10.50bp);
\end{tikzpicture}
\par\nopagebreak
\vspace{0.8mm}\nopagebreak
\begin{loesung}{A1Farbe}
\vspace{1.2mm}
\includegraphics[page=1,trim={52.00bp 460.50bp 39.00bp 344.50bp},clip,scale=0.700]{L_SoSe23.pdf}
\par
\par\vspace{1.0mm}
\includegraphics[page=1,trim={52.00bp 362.50bp 39.00bp 391.50bp},clip,scale=0.700]{L_SoSe23.pdf}
\par
\end{loesung}
\phantomsection\label{tf:SoSe23:1:4}
\par\vspace{1.6mm}
\Needspace*{33mm}
\begin{tikzpicture}[inner sep=0]
\node[anchor=north west,inner sep=0] (bi) at (0,0) {\includegraphics[page=2,trim={52.00bp 150.50bp 39.00bp 665.50bp},clip,scale=0.840]{A_SoSe23.pdf}};
\fill[MarkerG,opacity=0.30,rounded corners=0.6bp] (128.35bp,-11.34bp) rectangle (192.15bp,-21.42bp);
\end{tikzpicture}
\par\nopagebreak
\vspace{0.8mm}\nopagebreak
\begin{loesung}{A1Farbe}
\vspace{1.2mm}
\includegraphics[page=1,trim={52.00bp 251.50bp 39.00bp 524.50bp},clip,scale=0.700]{L_SoSe23.pdf}
\par
\end{loesung}
\par\vspace{6mm}
\Needspace*{192.7mm}
\marke{\color{A1Farbe}A1 \textendash{} WiSe 2023/24}
\phantomsection\label{auf:WiSe23x24:1}
\aufgabenkopf{WiSe 2023/24 \textendash{} Aufgabe 1}{19 Punkte}{A1Farbe}
\begin{tikzpicture}[inner sep=0]
\node[anchor=north west,inner sep=0] (bi) at (0,0) {\includegraphics[page=2,trim={52.00bp 682.50bp 39.00bp 119.50bp},clip,scale=0.840]{A_WiSe23x24.pdf}};
\fill[Marker,opacity=0.30,rounded corners=0.6bp] (339.36bp,-1.26bp) rectangle (367.92bp,-11.34bp);
\fill[Marker,opacity=0.30,rounded corners=0.6bp] (212.52bp,-12.18bp) rectangle (240.66bp,-22.26bp);
\fill[Marker,opacity=0.30,rounded corners=0.6bp] (269.92bp,-12.18bp) rectangle (289.80bp,-22.26bp);
\fill[Marker,opacity=0.30,rounded corners=0.6bp] (23.80bp,-23.10bp) rectangle (43.68bp,-33.18bp);
\end{tikzpicture}
\par\vspace{1.2mm}
\includegraphics[page=2,trim={52.00bp 484.50bp 39.00bp 193.50bp},clip,scale=0.620]{A_WiSe23x24.pdf}
\par\vspace{1.2mm}
\phantomsection\label{tf:WiSe23x24:1:1}
\Needspace*{18mm}
\verweis{A1Farbe}{1.}{Eindeutigkeit: welchen \gs{Innendurchmesser $D$} darf der Außenleiter maximal haben?}{gleiche Frage, nur andere Zahlenwerte \textendash{} siehe \textbf{SoSe 2022, Aufgabe 1.1}, \textbf{Seite \pageref{tf:SoSe22:1:1}}}
\par\vspace{1.6mm}
\begin{tikzpicture}[inner sep=0]
\node[anchor=north west,inner sep=0] (bi) at (0,0) {\includegraphics[page=2,trim={52.00bp 353.50bp 39.00bp 475.50bp},clip,scale=0.840]{A_WiSe23x24.pdf}};
\fill[Marker,opacity=0.30,rounded corners=0.6bp] (214.85bp,-0.42bp) rectangle (240.61bp,-10.50bp);
\end{tikzpicture}
\par
\phantomsection\label{tf:WiSe23x24:1:2}
\par\vspace{1.6mm}
\Needspace*{21mm}
\begin{tikzpicture}[inner sep=0]
\node[anchor=north west,inner sep=0] (bi) at (0,0) {\includegraphics[page=2,trim={52.00bp 327.50bp 39.00bp 500.50bp},clip,scale=0.840]{A_WiSe23x24.pdf}};
\fill[MarkerG,opacity=0.30,rounded corners=0.6bp] (68.76bp,-1.26bp) rectangle (111.90bp,-11.34bp);
\end{tikzpicture}
\par\nopagebreak
\vspace{0.8mm}\nopagebreak
\begin{loesung}{A1Farbe}
\vspace{1.2mm}
\includegraphics[page=1,trim={52.00bp 550.50bp 39.00bp 258.50bp},clip,scale=0.700]{L_WiSe23x24.pdf}
\par
\end{loesung}
\par\vspace{1.6mm}
\begin{tikzpicture}[inner sep=0]
\node[anchor=north west,inner sep=0] (bi) at (0,0) {\includegraphics[page=2,trim={52.00bp 289.50bp 39.00bp 538.50bp},clip,scale=0.840]{A_WiSe23x24.pdf}};
\fill[Marker,opacity=0.30,rounded corners=0.6bp] (75.60bp,-1.26bp) rectangle (112.98bp,-11.34bp);
\fill[Marker,opacity=0.30,rounded corners=0.6bp] (138.60bp,-1.26bp) rectangle (178.42bp,-11.34bp);
\end{tikzpicture}
\par
\phantomsection\label{tf:WiSe23x24:1:3}
\par\vspace{1.6mm}
\Needspace*{22mm}
\begin{tikzpicture}[inner sep=0]
\node[anchor=north west,inner sep=0] (bi) at (0,0) {\includegraphics[page=2,trim={52.00bp 264.50bp 39.00bp 563.50bp},clip,scale=0.840]{A_WiSe23x24.pdf}};
\fill[MarkerG,opacity=0.30,rounded corners=0.6bp] (57.12bp,-1.26bp) rectangle (130.20bp,-11.34bp);
\fill[MarkerG,opacity=0.30,rounded corners=0.6bp] (128.52bp,-1.26bp) rectangle (135.24bp,-11.34bp);
\end{tikzpicture}
\par\nopagebreak
\vspace{0.8mm}\nopagebreak
\begin{loesung}{A1Farbe}
\vspace{1.2mm}
\includegraphics[page=1,trim={52.00bp 480.50bp 39.00bp 325.50bp},clip,scale=0.700]{L_WiSe23x24.pdf}
\par
\end{loesung}
\phantomsection\label{tf:WiSe23x24:1:4}
\Needspace*{18mm}
\verweis{A1Farbe}{4.}{\gs{Gesamtdämpfung in dB}: Leiterverluste \emph{plus} dielektrische Verluste}{gleiche Frage, nur andere Zahlenwerte \textendash{} siehe \textbf{WiSe 2022/23, Aufgabe 1.3}, \textbf{Seite \pageref{tf:WiSe22x23:1:3}}}
\phantomsection\label{tf:WiSe23x24:1:5}
\par\vspace{1.6mm}
\Needspace*{39mm}
\begin{tikzpicture}[inner sep=0]
\node[anchor=north west,inner sep=0] (bi) at (0,0) {\includegraphics[page=2,trim={52.00bp 174.50bp 39.00bp 640.50bp},clip,scale=0.840]{A_WiSe23x24.pdf}};
\fill[MarkerG,opacity=0.30,rounded corners=0.6bp] (100.95bp,-1.26bp) rectangle (164.57bp,-11.34bp);
\end{tikzpicture}
\par\nopagebreak
\vspace{0.8mm}\nopagebreak
\begin{loesung}{A1Farbe}
\vspace{1.2mm}
\includegraphics[page=1,trim={52.00bp 172.50bp 39.00bp 638.50bp},clip,scale=0.700]{L_WiSe23x24.pdf}
\par
\par\vspace{1.0mm}
\includegraphics[page=1,trim={52.00bp 104.50bp 39.00bp 679.50bp},clip,scale=0.700]{L_WiSe23x24.pdf}
\par
\end{loesung}
\par\vspace{6mm}
\Needspace*{273.0mm}
\marke{\color{A1Farbe}A1 \textendash{} SoSe 2025}
\phantomsection\label{auf:SoSe25:1}
\aufgabenkopf{SoSe 2025 \textendash{} Aufgabe 1}{23 Punkte}{A1Farbe}
\begin{tikzpicture}[inner sep=0]
\node[anchor=north west,inner sep=0] (bi) at (0,0) {\includegraphics[page=2,trim={52.00bp 644.50bp 39.00bp 106.50bp},clip,scale=0.840]{A_SoSe25.pdf}};
\fill[Marker,opacity=0.30,rounded corners=0.6bp] (168.84bp,-1.26bp) rectangle (214.58bp,-11.34bp);
\fill[Marker,opacity=0.30,rounded corners=0.6bp] (327.20bp,-1.26bp) rectangle (351.67bp,-11.34bp);
\fill[Marker,opacity=0.30,rounded corners=0.6bp] (307.44bp,-34.02bp) rectangle (340.20bp,-44.10bp);
\fill[Marker,opacity=0.30,rounded corners=0.6bp] (24.50bp,-66.78bp) rectangle (43.68bp,-76.86bp);
\fill[Marker,opacity=0.30,rounded corners=0.6bp] (78.12bp,-66.78bp) rectangle (102.48bp,-76.86bp);
\end{tikzpicture}
\par\vspace{1.2mm}
\includegraphics[page=2,trim={52.00bp 457.50bp 39.00bp 220.50bp},clip,scale=0.620]{A_SoSe25.pdf}
\par\vspace{1.2mm}
\phantomsection\label{tf:SoSe25:1:1}
\par\vspace{1.6mm}
\Needspace*{45mm}
\begin{tikzpicture}[inner sep=0]
\node[anchor=north west,inner sep=0] (bi) at (0,0) {\includegraphics[page=2,trim={52.00bp 353.50bp 39.00bp 462.50bp},clip,scale=0.840]{A_SoSe25.pdf}};
\fill[Marker,opacity=0.30,rounded corners=0.6bp] (256.20bp,-12.18bp) rectangle (284.51bp,-22.26bp);
\fill[MarkerG,opacity=0.30,rounded corners=0.6bp] (73.44bp,-1.26bp) rectangle (117.36bp,-11.34bp);
\end{tikzpicture}
\par\nopagebreak
\vspace{0.8mm}\nopagebreak
\begin{loesung}{A1Farbe}
\vspace{1.2mm}
\includegraphics[page=1,trim={52.00bp 599.50bp 39.00bp 127.50bp},clip,scale=0.700]{L_SoSe25.pdf}
\par
\end{loesung}
\par\vspace{1.6mm}
\begin{tikzpicture}[inner sep=0]
\node[anchor=north west,inner sep=0] (bi) at (0,0) {\includegraphics[page=2,trim={52.00bp 289.50bp 39.00bp 526.50bp},clip,scale=0.840]{A_SoSe25.pdf}};
\fill[Marker,opacity=0.30,rounded corners=0.6bp] (186.48bp,-0.42bp) rectangle (226.53bp,-10.50bp);
\fill[Marker,opacity=0.30,rounded corners=0.6bp] (5.04bp,-11.34bp) rectangle (44.89bp,-21.42bp);
\end{tikzpicture}
\par
\phantomsection\label{tf:SoSe25:1:2}
\par\vspace{1.6mm}
\Needspace*{30mm}
\begin{tikzpicture}[inner sep=0]
\node[anchor=north west,inner sep=0] (bi) at (0,0) {\includegraphics[page=2,trim={52.00bp 251.50bp 39.00bp 563.50bp},clip,scale=0.840]{A_SoSe25.pdf}};
\fill[MarkerG,opacity=0.30,rounded corners=0.6bp] (80.22bp,-1.26bp) rectangle (112.14bp,-11.34bp);
\end{tikzpicture}
\par\nopagebreak
\vspace{0.8mm}\nopagebreak
\begin{loesung}{A1Farbe}
\vspace{1.2mm}
\includegraphics[page=1,trim={52.00bp 532.50bp 39.00bp 257.50bp},clip,scale=0.700]{L_SoSe25.pdf}
\par
\end{loesung}
\phantomsection\label{tf:SoSe25:1:3}
\par\vspace{1.6mm}
\Needspace*{98mm}
\begin{tikzpicture}[inner sep=0]
\node[anchor=north west,inner sep=0] (bi) at (0,0) {\includegraphics[page=2,trim={52.00bp 176.50bp 39.00bp 614.50bp},clip,scale=0.840]{A_SoSe25.pdf}};
\fill[Marker,opacity=0.30,rounded corners=0.6bp] (246.96bp,-1.26bp) rectangle (290.28bp,-11.34bp);
\fill[Marker,opacity=0.30,rounded corners=0.6bp] (304.92bp,-11.34bp) rectangle (336.32bp,-21.42bp);
\fill[MarkerG,opacity=0.30,rounded corners=0.6bp] (93.12bp,-1.26bp) rectangle (160.54bp,-11.34bp);
\end{tikzpicture}
\par\nopagebreak
\vspace{0.8mm}\nopagebreak
\begin{loesung}{A1Farbe}
\vspace{1.2mm}
\includegraphics[page=1,trim={52.00bp 207.50bp 39.00bp 333.50bp},clip,scale=0.700]{L_SoSe25.pdf}
\par
\end{loesung}
\phantomsection\label{tf:SoSe25:1:4}
\par\vspace{1.6mm}
\Needspace*{45mm}
\begin{tikzpicture}[inner sep=0]
\node[anchor=north west,inner sep=0] (bi) at (0,0) {\includegraphics[page=2,trim={52.00bp 125.50bp 39.00bp 690.50bp},clip,scale=0.840]{A_SoSe25.pdf}};
\fill[Marker,opacity=0.30,rounded corners=0.6bp] (214.32bp,-11.34bp) rectangle (231.90bp,-21.42bp);
\fill[MarkerG,opacity=0.30,rounded corners=0.6bp] (66.80bp,-1.26bp) rectangle (127.49bp,-11.34bp);
\end{tikzpicture}
\par\nopagebreak
\vspace{0.8mm}\nopagebreak
\begin{loesung}{A1Farbe}
\vspace{1.2mm}
\includegraphics[page=1,trim={52.00bp 51.50bp 39.00bp 674.50bp},clip,scale=0.700]{L_SoSe25.pdf}
\par
\end{loesung}
\par\vspace{6mm}
\clearpage
\marke{\color{A2Farbe}A2 \textendash{} Rechteckhohlleiter}
\phantomsection\label{typ:2}
\par\vspace*{-3mm}{\color{A2Farbe}\rule{\textwidth}{2.5pt}}\par\vspace{1.5mm}
{\sffamily\bfseries\huge\color{A2Farbe} A2 \textendash{} Rechteckhohlleiter}\par\vspace{1mm}
{\sffamily\color{Grau} alle 8 Klausuren \textendash{} 18 verschiedene Teilaufgaben ausgeschrieben, 7 Wiederholungen als Verweis}\par\vspace{2mm}
{\color{A2Farbe}\rule{\textwidth}{0.8pt}}\par\vspace{4mm}
{\sffamily\bfseries Register: Wenn gefragt wird \ldots}\hfill{\footnotesize\color{Grau}\colorbox{MarkerG!45}{grün} = das ist gesucht \textbf{·} rechts steht, wo die Aufgabe ausgeschrieben ist}\par\vspace{1.5mm}
\renewcommand{\arraystretch}{1.12}\footnotesize
\begin{tabularx}{\textwidth}{@{}>{\RaggedRight}X>{\RaggedRight\scriptsize}p{27mm}@{}}
\hline\rule{0pt}{3.5mm}
\gs{Frequenzbereich}, in dem genau ein Wellentyp ausbreitungsfähig ist \newline \hspace*{3mm}\textcolor{Grau}{\itshape \textendash{} Grundfall} & \textbf{S.\,\pageref{tf:WiSe20x21:2:1}} \textcolor{Grau}{WiSe 20/21, 1.} \\
\gs{Frequenzbereich}, in dem genau ein Wellentyp ausbreitungsfähig ist \newline \hspace*{3mm}\textcolor{Grau}{\itshape \textendash{} hier mit Dielektrikum gefüllt, also alle $f_c$ durch $\sqrt{\varepsilon_r}$} & \textbf{S.\,\pageref{tf:SoSe21:2:1}} \textcolor{Grau}{SoSe 21, 1.} \\
\gs{Frequenzbereich}, in dem genau ein Wellentyp ausbreitungsfähig ist \newline \hspace*{3mm}\textcolor{Grau}{\itshape \textendash{} hier mit Dielektrikum gefüllt ($\varepsilon_r=3{,}0$) und als Leiterplatten-Hohlleiter (SIW) verpackt} & \textbf{S.\,\pageref{tf:WiSe21x22:2:1}} \textcolor{Grau}{WiSe 21/22, 1.} \\
Maximale Amplitude \gs{$\hat{E}$} der elektrischen Feldstärke bei gegebener Leistung & \textbf{S.\,\pageref{tf:WiSe20x21:2:2}} \textcolor{Grau}{WiSe 20/21, 2.} \\
Maximale transversale magnetische Feldstärke \gs{$\hat{H}_t$} & \textbf{S.\,\pageref{tf:WiSe20x21:2:3}} \textcolor{Grau}{WiSe 20/21, 3.} \\
\gs{Dämpfung in dB} allein durch die \emph{dielektrischen} Verluste & \textbf{S.\,\pageref{tf:SoSe21:2:2}} \textcolor{Grau}{SoSe 21, 2.} \\
\gs{Dämpfung in dB} durch die \emph{Wandverluste} beim luftgefüllten Hohlleiter & \textbf{S.\,\pageref{tf:WiSe22x23:2:2}} \textcolor{Grau}{WiSe 22/23, 2.} \\
\gs{Dämpfung in dB} durch die Wandverluste beim \emph{dielektrisch gefüllten} Hohlleiter (andere Wellenlänge in der Formel!) & \textbf{S.\,\pageref{tf:WiSe21x22:2:2}} \textcolor{Grau}{WiSe 21/22, 2.} \\
Maximal übertragbare \gs{Leistung} bei vorgegebener Feldstärke & \textbf{S.\,\pageref{tf:SoSe22:2:2}} \textcolor{Grau}{SoSe 22, 2.} \\
Umkehrung: welche \gs{Länge} darf der Hohlleiter bei vorgegebener Dämpfung haben? & \textbf{S.\,\pageref{tf:SoSe22:2:4}} \textcolor{Grau}{SoSe 22, 4.} \\
Dielektrikum eingefüllt: ab welcher \gs{Frequenz} ausbreitungsfähig bzw. welches \gs{$\varepsilon_r$} ist dafür nötig? & \textbf{S.\,\pageref{tf:WiSe22x23:2:3}} \textcolor{Grau}{WiSe 22/23, 3.} \\
\gs{$\tan\delta$} aus der dielektrischen Dämpfung bzw. für ein Dämpfungs-Budget \newline \hspace*{3mm}\textcolor{Grau}{\itshape \textendash{} Grundfall} & \textbf{S.\,\pageref{tf:WiSe22x23:2:4}} \textcolor{Grau}{WiSe 22/23, 4.} \\
\gs{$\tan\delta$} aus der dielektrischen Dämpfung bzw. für ein Dämpfungs-Budget \newline \hspace*{3mm}\textcolor{Grau}{\itshape \textendash{} hier aus einer \emph{gemessenen} Dämpfung} & \textbf{S.\,\pageref{tf:SoSe23:2:3}} \textcolor{Grau}{SoSe 23, 3.} \\
\gs{$\varepsilon_r$-Bereich}, in dem der gefüllte Hohlleiter eindeutig bleibt & \textbf{S.\,\pageref{tf:SoSe23:2:1}} \textcolor{Grau}{SoSe 23, 1.} \\
\gs{$\varepsilon_r$} aus der gemessenen Hohlleiterwellenlänge $\lambda_H$ & \textbf{S.\,\pageref{tf:SoSe23:2:2}} \textcolor{Grau}{SoSe 23, 2.} \\
Prüfen, \gs{ob} der gefüllte Hohlleiter noch im Eindeutigkeitsbereich liegt (ja/nein) & \textbf{S.\,\pageref{tf:WiSe23x24:2:4}} \textcolor{Grau}{WiSe 23/24, 4.} \\
\gs{$a$ und $b$} rückwärts aus zwei beobachteten Grenzfrequenzen & \textbf{S.\,\pageref{tf:SoSe25:2:1}} \textcolor{Grau}{SoSe 25, 1.} \\
Hohlraumresonator: \gs{Mindestlänge} für eine Resonanzfrequenz & \textbf{S.\,\pageref{tf:SoSe25:2:3}} \textcolor{Grau}{SoSe 25, 3.} \\
\hline\end{tabularx}\normalsize
\par\vspace{5mm}
\begin{tcolorbox}[colback=A2Farbe!4, colframe=A2Farbe!55, boxrule=0.5pt, arc=1mm, left=3mm, right=2mm, top=2mm, bottom=2mm, breakable, title={\sffamily\bfseries\small Was diesmal neu drankommen könnte \normalfont\itshape\footnotesize -- eigene Prognose (keine offizielle Ankündigung): Umkehrungen, Kombinationen und Transfers zwischen den Aufgabentypen}, coltitle=white, colbacktitle=A2Farbe!70]
\sffamily\small
\begin{itemize}\setlength{\itemsep}{2.2mm}\setlength{\leftmargini}{4mm}
\item Rundhohlleiter: Mindestdurchmesser für den Grundwellentyp bzw. Eindeutigkeitsbereich des Rundhohlleiters\\[0.6mm] {\footnotesize\color{A2Farbe}$\rightarrow$ gleiche Bauart wie die $H_{01}$-Rechnung, nur mit der Konstante des jeweiligen Wellentyps -- FS 7.5 (Rundhohlleiter)}
\item Hohlraumresonator: Länge für eine \emph{höhere} Ordnung ($l=n\cdot\lambda_H/2$) oder mit Dielektrikum gefüllt\\[0.6mm] {\footnotesize\color{A2Farbe}$\rightarrow$ $l=n\cdot\lambda_H/2$ mit $\lambda_H=\lambda/\sqrt{1-(f_c/f)^2}$; bei Füllung überall $\lambda\to\lambda/\sqrt{\varepsilon_r}$}
\item Wandmaterial: um wieviel dB sinkt die Dämpfung durch Versilbern? Oder: $\kappa$ des Materials aus der gemessenen Dämpfung bestimmen\\[0.6mm] {\footnotesize\color{A2Farbe}$\rightarrow$ $\alpha\sim1/\sqrt{\kappa}$: Verhältnis der Leitfähigkeiten bilden bzw.\ die Dämpfungsformel nach $\kappa$ auflösen}
\item Phasen- und Gruppengeschwindigkeit bzw.\ $\lambda_H$ direkt gefragt (statt über den Umweg der $\varepsilon_r$-Bestimmung)\\[0.6mm] {\footnotesize\color{A2Farbe}$\rightarrow$ $v_{ph}=c/\sqrt{1-(f_c/f)^2}$ (grösser als $c$!), $v_{gr}=c\sqrt{1-(f_c/f)^2}$, $\lambda_H$ analog. FS 7.3}
\item Leiterverluste \emph{und} dielektrische Verluste in \emph{einer} Teilaufgabe -- bisher wurden sie fast immer getrennt abgefragt\\[0.6mm] {\footnotesize\color{A2Farbe}$\rightarrow$ $\alpha_{ges}=\alpha_L+\alpha_D$ addieren, \emph{dann} erst $a_{dB}=8{,}686\,\alpha_{ges}\,l$}
\item Vergleich Koaxialkabel $\leftrightarrow$ Hohlleiter: welcher hat bei dieser Frequenz die geringere Dämpfung?\\[0.6mm] {\footnotesize\color{A2Farbe}$\rightarrow$ beide $\alpha$ bei derselben Frequenz ausrechnen und in dB/m vergleichen (K1 gegen K2)}
\end{itemize}\end{tcolorbox}
\clearpage
\Needspace*{203.6mm}
\marke{\color{A2Farbe}A2 \textendash{} WiSe 2020/21}
\phantomsection\label{auf:WiSe20x21:2}
\aufgabenkopf{WiSe 2020/21 \textendash{} Aufgabe 2}{13 Punkte}{A2Farbe}
\begin{tikzpicture}[inner sep=0]
\node[anchor=north west,inner sep=0] (bi) at (0,0) {\includegraphics[page=3,trim={52.00bp 708.50bp 39.00bp 106.50bp},clip,scale=0.840]{A_WiSe20x21.pdf}};
\fill[Marker,opacity=0.30,rounded corners=0.6bp] (286.44bp,-1.26bp) rectangle (320.88bp,-11.34bp);
\fill[Marker,opacity=0.30,rounded corners=0.6bp] (356.65bp,-1.26bp) rectangle (381.46bp,-11.34bp);
\fill[Marker,opacity=0.30,rounded corners=0.6bp] (32.76bp,-12.18bp) rectangle (59.64bp,-22.26bp);
\end{tikzpicture}
\par\vspace{1.2mm}
\includegraphics[page=3,trim={52.00bp 493.50bp 39.00bp 167.50bp},clip,scale=0.620]{A_WiSe20x21.pdf}
\par\vspace{1.2mm}
\phantomsection\label{tf:WiSe20x21:2:1}
\par\vspace{1.6mm}
\Needspace*{39mm}
\begin{tikzpicture}[inner sep=0]
\node[anchor=north west,inner sep=0] (bi) at (0,0) {\includegraphics[page=3,trim={52.00bp 392.50bp 39.00bp 422.50bp},clip,scale=0.840]{A_WiSe20x21.pdf}};
\fill[MarkerG,opacity=0.30,rounded corners=0.6bp] (65.20bp,-1.26bp) rectangle (123.72bp,-11.34bp);
\end{tikzpicture}
\par\nopagebreak
\vspace{0.8mm}\nopagebreak
\begin{loesung}{A2Farbe}
\vspace{1.2mm}
\includegraphics[page=2,trim={52.00bp 719.50bp 39.00bp 88.50bp},clip,scale=0.700]{L_WiSe20x21.pdf}
\par
\par\vspace{1.0mm}
\includegraphics[page=2,trim={52.00bp 651.50bp 39.00bp 135.50bp},clip,scale=0.700]{L_WiSe20x21.pdf}
\par
\end{loesung}
\phantomsection\label{tf:WiSe20x21:2:2}
\par\vspace{1.6mm}
\Needspace*{45mm}
\begin{tikzpicture}[inner sep=0]
\node[anchor=north west,inner sep=0] (bi) at (0,0) {\includegraphics[page=3,trim={52.00bp 329.50bp 39.00bp 484.50bp},clip,scale=0.840]{A_WiSe20x21.pdf}};
\fill[Marker,opacity=0.30,rounded corners=0.6bp] (112.13bp,-13.02bp) rectangle (143.93bp,-23.10bp);
\fill[Marker,opacity=0.30,rounded corners=0.6bp] (226.45bp,-13.02bp) rectangle (251.25bp,-23.10bp);
\fill[MarkerG,opacity=0.30,rounded corners=0.6bp] (86.74bp,-2.94bp) rectangle (123.98bp,-13.02bp);
\end{tikzpicture}
\par\nopagebreak
\vspace{0.8mm}\nopagebreak
\begin{loesung}{A2Farbe}
\vspace{1.2mm}
\includegraphics[page=2,trim={52.00bp 469.50bp 39.00bp 259.50bp},clip,scale=0.700]{L_WiSe20x21.pdf}
\par
\end{loesung}
\phantomsection\label{tf:WiSe20x21:2:3}
\par\vspace{1.6mm}
\Needspace*{44mm}
\begin{tikzpicture}[inner sep=0]
\node[anchor=north west,inner sep=0] (bi) at (0,0) {\includegraphics[page=3,trim={52.00bp 253.50bp 39.00bp 562.50bp},clip,scale=0.840]{A_WiSe20x21.pdf}};
\fill[Marker,opacity=0.30,rounded corners=0.6bp] (208.32bp,-11.34bp) rectangle (251.16bp,-21.42bp);
\fill[MarkerG,opacity=0.30,rounded corners=0.6bp] (123.55bp,-0.42bp) rectangle (159.86bp,-10.50bp);
\end{tikzpicture}
\par\nopagebreak
\vspace{0.8mm}\nopagebreak
\begin{loesung}{A2Farbe}
\vspace{1.2mm}
\includegraphics[page=2,trim={52.00bp 346.50bp 39.00bp 429.50bp},clip,scale=0.700]{L_WiSe20x21.pdf}
\par
\par\vspace{1.0mm}
\includegraphics[page=2,trim={52.00bp 290.50bp 39.00bp 505.50bp},clip,scale=0.700]{L_WiSe20x21.pdf}
\par
\end{loesung}
\par\vspace{6mm}
\Needspace*{149.5mm}
\marke{\color{A2Farbe}A2 \textendash{} SoSe 2021}
\phantomsection\label{auf:SoSe21:2}
\aufgabenkopf{SoSe 2021 \textendash{} Aufgabe 2}{11 Punkte}{A2Farbe}
\begin{tikzpicture}[inner sep=0]
\node[anchor=north west,inner sep=0] (bi) at (0,0) {\includegraphics[page=3,trim={52.00bp 697.50bp 39.00bp 106.50bp},clip,scale=0.840]{A_SoSe21.pdf}};
\fill[Marker,opacity=0.30,rounded corners=0.6bp] (171.36bp,-2.10bp) rectangle (205.50bp,-12.18bp);
\fill[Marker,opacity=0.30,rounded corners=0.6bp] (246.12bp,-2.10bp) rectangle (264.60bp,-12.18bp);
\fill[Marker,opacity=0.30,rounded corners=0.6bp] (61.32bp,-13.02bp) rectangle (96.60bp,-23.10bp);
\fill[Marker,opacity=0.30,rounded corners=0.6bp] (122.64bp,-13.02bp) rectangle (157.20bp,-23.10bp);
\fill[Marker,opacity=0.30,rounded corners=0.6bp] (231.84bp,-13.02bp) rectangle (263.18bp,-23.10bp);
\fill[Marker,opacity=0.30,rounded corners=0.6bp] (369.60bp,-13.02bp) rectangle (414.20bp,-23.10bp);
\end{tikzpicture}
\par\vspace{1.2mm}
\includegraphics[page=3,trim={52.00bp 479.50bp 39.00bp 181.50bp},clip,scale=0.620]{A_SoSe21.pdf}
\par\vspace{1.2mm}
\phantomsection\label{tf:SoSe21:2:1}
\par\vspace{1.6mm}
\Needspace*{27mm}
\begin{tikzpicture}[inner sep=0]
\node[anchor=north west,inner sep=0] (bi) at (0,0) {\includegraphics[page=3,trim={52.00bp 379.50bp 39.00bp 436.50bp},clip,scale=0.840]{A_SoSe21.pdf}};
\fill[MarkerG,opacity=0.30,rounded corners=0.6bp] (118.23bp,-1.26bp) rectangle (150.22bp,-11.34bp);
\end{tikzpicture}
\par\nopagebreak
\vspace{0.8mm}\nopagebreak
\begin{loesung}{A2Farbe}
\vspace{1.2mm}
\includegraphics[page=2,trim={52.00bp 675.50bp 39.00bp 126.50bp},clip,scale=0.700]{L_SoSe21.pdf}
\par
\end{loesung}
\phantomsection\label{tf:SoSe21:2:2}
\par\vspace{1.6mm}
\Needspace*{52mm}
\begin{tikzpicture}[inner sep=0]
\node[anchor=north west,inner sep=0] (bi) at (0,0) {\includegraphics[page=3,trim={52.00bp 315.50bp 39.00bp 500.50bp},clip,scale=0.840]{A_SoSe21.pdf}};
\fill[MarkerG,opacity=0.30,rounded corners=0.6bp] (55.12bp,-1.26bp) rectangle (97.44bp,-11.34bp);
\fill[MarkerG,opacity=0.30,rounded corners=0.6bp] (95.76bp,-1.26bp) rectangle (102.48bp,-11.34bp);
\end{tikzpicture}
\par\nopagebreak
\vspace{0.8mm}\nopagebreak
\begin{loesung}{A2Farbe}
\vspace{1.2mm}
\includegraphics[page=2,trim={52.00bp 457.50bp 39.00bp 240.50bp},clip,scale=0.700]{L_SoSe21.pdf}
\par
\end{loesung}
\par\vspace{6mm}
\Needspace*{215.6mm}
\marke{\color{A2Farbe}A2 \textendash{} WiSe 2021/22}
\phantomsection\label{auf:WiSe21x22:2}
\aufgabenkopf{WiSe 2021/22 \textendash{} Aufgabe 2}{12 Punkte}{A2Farbe}
\begin{tikzpicture}[inner sep=0]
\node[anchor=north west,inner sep=0] (bi) at (0,0) {\includegraphics[page=3,trim={52.00bp 631.50bp 39.00bp 106.50bp},clip,scale=0.840]{A_WiSe21x22.pdf}};
\fill[Marker,opacity=0.30,rounded corners=0.6bp] (19.32bp,-55.02bp) rectangle (51.15bp,-65.10bp);
\fill[Marker,opacity=0.30,rounded corners=0.6bp] (144.22bp,-55.02bp) rectangle (177.45bp,-65.10bp);
\fill[Marker,opacity=0.30,rounded corners=0.6bp] (13.85bp,-65.94bp) rectangle (41.93bp,-76.02bp);
\fill[Marker,opacity=0.30,rounded corners=0.6bp] (178.92bp,-65.94bp) rectangle (211.40bp,-76.02bp);
\fill[Marker,opacity=0.30,rounded corners=0.6bp] (24.50bp,-76.86bp) rectangle (43.68bp,-86.94bp);
\end{tikzpicture}
\par\vspace{1.2mm}
\includegraphics[page=3,trim={52.00bp 400.50bp 39.00bp 236.50bp},clip,scale=0.620]{A_WiSe21x22.pdf}
\par\vspace{1.2mm}
\phantomsection\label{tf:WiSe21x22:2:1}
\par\vspace{1.6mm}
\Needspace*{36mm}
\begin{tikzpicture}[inner sep=0]
\node[anchor=north west,inner sep=0] (bi) at (0,0) {\includegraphics[page=3,trim={52.00bp 302.50bp 39.00bp 513.50bp},clip,scale=0.840]{A_WiSe21x22.pdf}};
\fill[MarkerG,opacity=0.30,rounded corners=0.6bp] (65.76bp,-1.26bp) rectangle (125.04bp,-11.34bp);
\end{tikzpicture}
\par\nopagebreak
\vspace{0.8mm}\nopagebreak
\begin{loesung}{A2Farbe}
\vspace{1.2mm}
\includegraphics[page=2,trim={52.00bp 658.50bp 39.00bp 120.50bp},clip,scale=0.700]{L_WiSe21x22.pdf}
\par
\par\vspace{1.0mm}
\includegraphics[page=2,trim={52.00bp 625.50bp 39.00bp 201.50bp},clip,scale=0.700]{L_WiSe21x22.pdf}
\par
\end{loesung}
\phantomsection\label{tf:WiSe21x22:2:2}
\par\vspace{1.6mm}
\Needspace*{80mm}
\begin{tikzpicture}[inner sep=0]
\node[anchor=north west,inner sep=0] (bi) at (0,0) {\includegraphics[page=3,trim={52.00bp 201.50bp 39.00bp 576.50bp},clip,scale=0.840]{A_WiSe21x22.pdf}};
\fill[Marker,opacity=0.30,rounded corners=0.6bp] (275.29bp,-1.26bp) rectangle (304.77bp,-11.34bp);
\fill[MarkerG,opacity=0.30,rounded corners=0.6bp] (55.12bp,-1.26bp) rectangle (97.44bp,-11.34bp);
\fill[MarkerG,opacity=0.30,rounded corners=0.6bp] (95.76bp,-1.26bp) rectangle (102.48bp,-11.34bp);
\end{tikzpicture}
\par\nopagebreak
\vspace{0.8mm}\nopagebreak
\begin{loesung}{A2Farbe}
\vspace{1.2mm}
\includegraphics[page=2,trim={52.00bp 536.50bp 39.00bp 270.50bp},clip,scale=0.700]{L_WiSe21x22.pdf}
\par
\par\vspace{1.0mm}
\includegraphics[page=2,trim={52.00bp 492.50bp 39.00bp 335.50bp},clip,scale=0.700]{L_WiSe21x22.pdf}
\par
\par\vspace{1.0mm}
\includegraphics[page=2,trim={52.00bp 306.50bp 39.00bp 375.50bp},clip,scale=0.700]{L_WiSe21x22.pdf}
\par
\end{loesung}
\par\vspace{6mm}
\Needspace*{220.6mm}
\marke{\color{A2Farbe}A2 \textendash{} SoSe 2022}
\phantomsection\label{auf:SoSe22:2}
\aufgabenkopf{SoSe 2022 \textendash{} Aufgabe 2}{20 Punkte}{A2Farbe}
\begin{tikzpicture}[inner sep=0]
\node[anchor=north west,inner sep=0] (bi) at (0,0) {\includegraphics[page=3,trim={52.00bp 695.50bp 39.00bp 105.50bp},clip,scale=0.840]{A_SoSe22.pdf}};
\fill[Marker,opacity=0.30,rounded corners=0.6bp] (298.20bp,-2.10bp) rectangle (332.64bp,-12.18bp);
\fill[Marker,opacity=0.30,rounded corners=0.6bp] (368.40bp,-2.10bp) rectangle (393.12bp,-12.18bp);
\fill[Marker,opacity=0.30,rounded corners=0.6bp] (121.45bp,-13.02bp) rectangle (153.91bp,-23.10bp);
\fill[Marker,opacity=0.30,rounded corners=0.6bp] (24.50bp,-23.94bp) rectangle (43.68bp,-34.02bp);
\end{tikzpicture}
\par\vspace{1.2mm}
\includegraphics[page=3,trim={52.00bp 479.50bp 39.00bp 181.50bp},clip,scale=0.620]{A_SoSe22.pdf}
\par\vspace{1.2mm}
\phantomsection\label{tf:SoSe22:2:1}
\Needspace*{18mm}
\verweis{A2Farbe}{1.}{\gs{Frequenzbereich}, in dem genau ein Wellentyp ausbreitungsfähig ist}{gleiche Frage, nur andere Zahlenwerte \textendash{} siehe \textbf{WiSe 2020/21, Aufgabe 2.1}, \textbf{Seite \pageref{tf:WiSe20x21:2:1}}}
\phantomsection\label{tf:SoSe22:2:2}
\par\vspace{1.6mm}
\Needspace*{39mm}
\begin{tikzpicture}[inner sep=0]
\node[anchor=north west,inner sep=0] (bi) at (0,0) {\includegraphics[page=3,trim={52.00bp 315.50bp 39.00bp 499.50bp},clip,scale=0.840]{A_SoSe22.pdf}};
\fill[Marker,opacity=0.30,rounded corners=0.6bp] (55.69bp,-12.18bp) rectangle (77.48bp,-22.26bp);
\fill[MarkerG,opacity=0.30,rounded corners=0.6bp] (87.86bp,-1.26bp) rectangle (121.71bp,-11.34bp);
\end{tikzpicture}
\par\nopagebreak
\vspace{0.8mm}\nopagebreak
\begin{loesung}{A2Farbe}
\vspace{1.2mm}
\includegraphics[page=2,trim={52.00bp 531.50bp 39.00bp 221.50bp},clip,scale=0.700]{L_SoSe22.pdf}
\par
\end{loesung}
\phantomsection\label{tf:SoSe22:2:3}
\Needspace*{18mm}
\verweis{A2Farbe}{3.}{Maximale transversale magnetische Feldstärke \gs{$\hat{H}_t$}}{gleiche Frage, nur andere Zahlenwerte \textendash{} siehe \textbf{WiSe 2020/21, Aufgabe 2.3}, \textbf{Seite \pageref{tf:WiSe20x21:2:3}}}
\phantomsection\label{tf:SoSe22:2:4}
\par\vspace{1.6mm}
\Needspace*{81mm}
\begin{tikzpicture}[inner sep=0]
\node[anchor=north west,inner sep=0] (bi) at (0,0) {\includegraphics[page=3,trim={52.00bp 176.50bp 39.00bp 638.50bp},clip,scale=0.840]{A_SoSe22.pdf}};
\fill[Marker,opacity=0.30,rounded corners=0.6bp] (401.52bp,-1.26bp) rectangle (415.80bp,-11.34bp);
\fill[MarkerG,opacity=0.30,rounded corners=0.6bp] (51.05bp,-1.26bp) rectangle (72.40bp,-11.34bp);
\end{tikzpicture}
\par\nopagebreak
\vspace{0.8mm}\nopagebreak
\begin{loesung}{A2Farbe}
\vspace{1.2mm}
\includegraphics[page=2,trim={52.00bp 171.50bp 39.00bp 468.50bp},clip,scale=0.700]{L_SoSe22.pdf}
\par
\par\vspace{1.0mm}
\includegraphics[page=2,trim={52.00bp 133.50bp 39.00bp 680.50bp},clip,scale=0.700]{L_SoSe22.pdf}
\par
\par\vspace{1.0mm}
\includegraphics[page=2,trim={52.00bp 92.50bp 39.00bp 718.50bp},clip,scale=0.700]{L_SoSe22.pdf}
\par
\end{loesung}
\par\vspace{6mm}
\Needspace*{256.7mm}
\marke{\color{A2Farbe}A2 \textendash{} WiSe 2022/23}
\phantomsection\label{auf:WiSe22x23:2}
\aufgabenkopf{WiSe 2022/23 \textendash{} Aufgabe 2}{19 Punkte}{A2Farbe}
\begin{tikzpicture}[inner sep=0]
\node[anchor=north west,inner sep=0] (bi) at (0,0) {\includegraphics[page=3,trim={52.00bp 695.50bp 39.00bp 105.50bp},clip,scale=0.840]{A_WiSe22x23.pdf}};
\fill[Marker,opacity=0.30,rounded corners=0.6bp] (332.64bp,-2.10bp) rectangle (367.92bp,-12.18bp);
\fill[Marker,opacity=0.30,rounded corners=0.6bp] (378.00bp,-2.10bp) rectangle (414.96bp,-12.18bp);
\fill[Marker,opacity=0.30,rounded corners=0.6bp] (63.00bp,-13.02bp) rectangle (89.88bp,-23.10bp);
\fill[Marker,opacity=0.30,rounded corners=0.6bp] (218.40bp,-13.02bp) rectangle (260.59bp,-23.10bp);
\fill[Marker,opacity=0.30,rounded corners=0.6bp] (24.50bp,-23.94bp) rectangle (43.68bp,-34.02bp);
\end{tikzpicture}
\par\vspace{1.2mm}
\includegraphics[page=3,trim={52.00bp 479.50bp 39.00bp 181.50bp},clip,scale=0.620]{A_WiSe22x23.pdf}
\par\vspace{1.2mm}
\phantomsection\label{tf:WiSe22x23:2:1}
\Needspace*{18mm}
\verweis{A2Farbe}{1.}{\gs{Frequenzbereich}, in dem genau ein Wellentyp ausbreitungsfähig ist}{gleiche Frage, nur andere Zahlenwerte \textendash{} siehe \textbf{WiSe 2020/21, Aufgabe 2.1}, \textbf{Seite \pageref{tf:WiSe20x21:2:1}}}
\phantomsection\label{tf:WiSe22x23:2:2}
\par\vspace{1.6mm}
\Needspace*{72mm}
\begin{tikzpicture}[inner sep=0]
\node[anchor=north west,inner sep=0] (bi) at (0,0) {\includegraphics[page=3,trim={52.00bp 328.50bp 39.00bp 499.50bp},clip,scale=0.840]{A_WiSe22x23.pdf}};
\fill[MarkerG,opacity=0.30,rounded corners=0.6bp] (50.83bp,-1.26bp) rectangle (114.92bp,-11.34bp);
\end{tikzpicture}
\par\nopagebreak
\vspace{0.8mm}\nopagebreak
\begin{loesung}{A2Farbe}
\vspace{1.2mm}
\includegraphics[page=2,trim={52.00bp 383.50bp 39.00bp 220.50bp},clip,scale=0.700]{L_WiSe22x23.pdf}
\par
\end{loesung}
\par\vspace{1.6mm}
\begin{tikzpicture}[inner sep=0]
\node[anchor=north west,inner sep=0] (bi) at (0,0) {\includegraphics[page=3,trim={52.00bp 264.50bp 39.00bp 563.50bp},clip,scale=0.840]{A_WiSe22x23.pdf}};
\fill[Marker,opacity=0.30,rounded corners=0.6bp] (210.00bp,-1.26bp) rectangle (237.44bp,-11.34bp);
\end{tikzpicture}
\par
\phantomsection\label{tf:WiSe22x23:2:3}
\par\vspace{1.6mm}
\Needspace*{23mm}
\begin{tikzpicture}[inner sep=0]
\node[anchor=north west,inner sep=0] (bi) at (0,0) {\includegraphics[page=3,trim={52.00bp 214.50bp 39.00bp 614.50bp},clip,scale=0.840]{A_WiSe22x23.pdf}};
\fill[MarkerG,opacity=0.30,rounded corners=0.6bp] (66.56bp,-1.26bp) rectangle (99.53bp,-11.34bp);
\end{tikzpicture}
\par\nopagebreak
\vspace{0.8mm}\nopagebreak
\begin{loesung}{A2Farbe}
\vspace{1.2mm}
\includegraphics[page=2,trim={52.00bp 312.50bp 39.00bp 487.50bp},clip,scale=0.700]{L_WiSe22x23.pdf}
\par
\end{loesung}
\phantomsection\label{tf:WiSe22x23:2:4}
\par\vspace{1.6mm}
\Needspace*{66mm}
\begin{tikzpicture}[inner sep=0]
\node[anchor=north west,inner sep=0] (bi) at (0,0) {\includegraphics[page=3,trim={52.00bp 149.50bp 39.00bp 665.50bp},clip,scale=0.840]{A_WiSe22x23.pdf}};
\fill[Marker,opacity=0.30,rounded corners=0.6bp] (206.79bp,-12.18bp) rectangle (227.56bp,-22.26bp);
\fill[MarkerG,opacity=0.30,rounded corners=0.6bp] (90.91bp,-1.26bp) rectangle (141.26bp,-11.34bp);
\end{tikzpicture}
\par\nopagebreak
\vspace{0.8mm}\nopagebreak
\begin{loesung}{A2Farbe}
\vspace{1.2mm}
\includegraphics[page=2,trim={52.00bp 79.50bp 39.00bp 563.50bp},clip,scale=0.700]{L_WiSe22x23.pdf}
\par
\end{loesung}
\par\vspace{6mm}
\Needspace*{238.2mm}
\marke{\color{A2Farbe}A2 \textendash{} SoSe 2023}
\phantomsection\label{auf:SoSe23:2}
\aufgabenkopf{SoSe 2023 \textendash{} Aufgabe 2}{18 Punkte}{A2Farbe}
\begin{tikzpicture}[inner sep=0]
\node[anchor=north west,inner sep=0] (bi) at (0,0) {\includegraphics[page=3,trim={52.00bp 708.50bp 39.00bp 106.50bp},clip,scale=0.840]{A_SoSe23.pdf}};
\fill[Marker,opacity=0.30,rounded corners=0.6bp] (222.60bp,-1.26bp) rectangle (257.04bp,-11.34bp);
\fill[Marker,opacity=0.30,rounded corners=0.6bp] (267.12bp,-1.26bp) rectangle (306.13bp,-11.34bp);
\fill[Marker,opacity=0.30,rounded corners=0.6bp] (370.44bp,-1.26bp) rectangle (401.52bp,-11.34bp);
\fill[Marker,opacity=0.30,rounded corners=0.6bp] (193.20bp,-12.18bp) rectangle (235.60bp,-22.26bp);
\end{tikzpicture}
\par\vspace{1.2mm}
\includegraphics[page=3,trim={52.00bp 451.50bp 39.00bp 210.50bp},clip,scale=0.620]{A_SoSe23.pdf}
\par\vspace{1.2mm}
\phantomsection\label{tf:SoSe23:2:1}
\par\vspace{1.6mm}
\Needspace*{42mm}
\begin{tikzpicture}[inner sep=0]
\node[anchor=north west,inner sep=0] (bi) at (0,0) {\includegraphics[page=3,trim={52.00bp 353.50bp 39.00bp 461.50bp},clip,scale=0.840]{A_SoSe23.pdf}};
\fill[MarkerG,opacity=0.30,rounded corners=0.6bp] (65.36bp,-1.26bp) rectangle (93.67bp,-11.34bp);
\end{tikzpicture}
\par\nopagebreak
\vspace{0.8mm}\nopagebreak
\begin{loesung}{A2Farbe}
\vspace{1.2mm}
\includegraphics[page=2,trim={52.00bp 664.50bp 39.00bp 96.50bp},clip,scale=0.700]{L_SoSe23.pdf}
\par
\par\vspace{1.0mm}
\includegraphics[page=2,trim={52.00bp 635.50bp 39.00bp 187.50bp},clip,scale=0.700]{L_SoSe23.pdf}
\par
\end{loesung}
\phantomsection\label{tf:SoSe23:2:2}
\par\vspace{1.6mm}
\Needspace*{59mm}
\begin{tikzpicture}[inner sep=0]
\node[anchor=north west,inner sep=0] (bi) at (0,0) {\includegraphics[page=3,trim={52.00bp 277.50bp 39.00bp 538.50bp},clip,scale=0.840]{A_SoSe23.pdf}};
\fill[Marker,opacity=0.30,rounded corners=0.6bp] (230.19bp,-1.26bp) rectangle (257.27bp,-11.34bp);
\fill[MarkerG,opacity=0.30,rounded corners=0.6bp] (377.08bp,-1.26bp) rectangle (424.58bp,-11.34bp);
\end{tikzpicture}
\par\nopagebreak
\vspace{0.8mm}\nopagebreak
\begin{loesung}{A2Farbe}
\vspace{1.2mm}
\includegraphics[page=2,trim={52.00bp 412.50bp 39.00bp 256.50bp},clip,scale=0.700]{L_SoSe23.pdf}
\par
\end{loesung}
\phantomsection\label{tf:SoSe23:2:3}
\par\vspace{1.6mm}
\Needspace*{62mm}
\begin{tikzpicture}[inner sep=0]
\node[anchor=north west,inner sep=0] (bi) at (0,0) {\includegraphics[page=3,trim={52.00bp 174.50bp 39.00bp 601.50bp},clip,scale=0.840]{A_SoSe23.pdf}};
\fill[Marker,opacity=0.30,rounded corners=0.6bp] (170.52bp,-1.26bp) rectangle (211.14bp,-11.34bp);
\fill[Marker,opacity=0.30,rounded corners=0.6bp] (205.80bp,-44.94bp) rectangle (234.36bp,-55.02bp);
\fill[MarkerG,opacity=0.30,rounded corners=0.6bp] (116.52bp,-23.10bp) rectangle (166.56bp,-33.18bp);
\end{tikzpicture}
\par\nopagebreak
\vspace{0.8mm}\nopagebreak
\begin{loesung}{A2Farbe}
\vspace{1.2mm}
\includegraphics[page=2,trim={52.00bp 328.50bp 39.00bp 483.50bp},clip,scale=0.700]{L_SoSe23.pdf}
\par
\par\vspace{1.0mm}
\includegraphics[page=2,trim={52.00bp 211.50bp 39.00bp 523.50bp},clip,scale=0.700]{L_SoSe23.pdf}
\par
\end{loesung}
\par\vspace{6mm}
\Needspace*{157.9mm}
\marke{\color{A2Farbe}A2 \textendash{} WiSe 2023/24}
\phantomsection\label{auf:WiSe23x24:2}
\aufgabenkopf{WiSe 2023/24 \textendash{} Aufgabe 2}{18 Punkte}{A2Farbe}
\begin{tikzpicture}[inner sep=0]
\node[anchor=north west,inner sep=0] (bi) at (0,0) {\includegraphics[page=3,trim={52.00bp 695.50bp 39.00bp 105.50bp},clip,scale=0.840]{A_WiSe23x24.pdf}};
\fill[Marker,opacity=0.30,rounded corners=0.6bp] (337.68bp,-2.10bp) rectangle (371.95bp,-12.18bp);
\fill[Marker,opacity=0.30,rounded corners=0.6bp] (14.01bp,-13.02bp) rectangle (43.21bp,-23.10bp);
\fill[Marker,opacity=0.30,rounded corners=0.6bp] (107.52bp,-13.02bp) rectangle (134.88bp,-23.10bp);
\fill[Marker,opacity=0.30,rounded corners=0.6bp] (269.64bp,-13.02bp) rectangle (305.07bp,-23.10bp);
\fill[Marker,opacity=0.30,rounded corners=0.6bp] (128.88bp,-23.94bp) rectangle (152.88bp,-34.02bp);
\end{tikzpicture}
\par\vspace{1.2mm}
\includegraphics[page=3,trim={52.00bp 479.50bp 39.00bp 181.50bp},clip,scale=0.620]{A_WiSe23x24.pdf}
\par\vspace{1.2mm}
\phantomsection\label{tf:WiSe23x24:2:1}
\Needspace*{18mm}
\verweis{A2Farbe}{1.}{\gs{Frequenzbereich}, in dem genau ein Wellentyp ausbreitungsfähig ist}{gleiche Frage, nur andere Zahlenwerte \textendash{} siehe \textbf{WiSe 2020/21, Aufgabe 2.1}, \textbf{Seite \pageref{tf:WiSe20x21:2:1}}}
\phantomsection\label{tf:WiSe23x24:2:2}
\Needspace*{18mm}
\verweis{A2Farbe}{2.}{Maximal übertragbare \gs{Leistung} bei vorgegebener Feldstärke}{gleiche Frage, nur andere Zahlenwerte \textendash{} siehe \textbf{SoSe 2022, Aufgabe 2.2}, \textbf{Seite \pageref{tf:SoSe22:2:2}}}
\phantomsection\label{tf:WiSe23x24:2:3}
\Needspace*{18mm}
\verweis{A2Farbe}{3.}{\gs{Dämpfung in dB} durch die \emph{Wandverluste} beim luftgefüllten Hohlleiter}{gleiche Frage, nur andere Zahlenwerte \textendash{} siehe \textbf{WiSe 2022/23, Aufgabe 2.2}, \textbf{Seite \pageref{tf:WiSe22x23:2:2}}}
\phantomsection\label{tf:WiSe23x24:2:4}
\par\vspace{1.6mm}
\Needspace*{52mm}
\begin{tikzpicture}[inner sep=0]
\node[anchor=north west,inner sep=0] (bi) at (0,0) {\includegraphics[page=3,trim={52.00bp 201.50bp 39.00bp 601.50bp},clip,scale=0.840]{A_WiSe23x24.pdf}};
\fill[Marker,opacity=0.30,rounded corners=0.6bp] (221.76bp,-1.26bp) rectangle (255.10bp,-11.34bp);
\end{tikzpicture}
\par\nopagebreak
\vspace{0.8mm}\nopagebreak
\begin{loesung}{A2Farbe}
\vspace{1.2mm}
\includegraphics[page=2,trim={52.00bp 166.50bp 39.00bp 581.50bp},clip,scale=0.700]{L_WiSe23x24.pdf}
\par
\par\vspace{1.0mm}
\includegraphics[page=2,trim={52.00bp 113.50bp 39.00bp 696.50bp},clip,scale=0.700]{L_WiSe23x24.pdf}
\par
\end{loesung}
\par\vspace{6mm}
\Needspace*{191.2mm}
\marke{\color{A2Farbe}A2 \textendash{} SoSe 2025}
\phantomsection\label{auf:SoSe25:2}
\aufgabenkopf{SoSe 2025 \textendash{} Aufgabe 2}{18 Punkte}{A2Farbe}
\begin{tikzpicture}[inner sep=0]
\node[anchor=north west,inner sep=0] (bi) at (0,0) {\includegraphics[page=3,trim={52.00bp 696.50bp 39.00bp 106.50bp},clip,scale=0.840]{A_SoSe25.pdf}};
\fill[Marker,opacity=0.30,rounded corners=0.6bp] (260.11bp,-1.26bp) rectangle (293.10bp,-11.34bp);
\fill[Marker,opacity=0.30,rounded corners=0.6bp] (81.90bp,-12.18bp) rectangle (117.65bp,-22.26bp);
\end{tikzpicture}
\par\vspace{1.2mm}
\includegraphics[page=3,trim={52.00bp 573.50bp 39.00bp 173.50bp},clip,scale=0.620]{A_SoSe25.pdf}
\par\vspace{1.2mm}
\phantomsection\label{tf:SoSe25:2:1}
\par\vspace{1.6mm}
\Needspace*{29mm}
\begin{tikzpicture}[inner sep=0]
\node[anchor=north west,inner sep=0] (bi) at (0,0) {\includegraphics[page=3,trim={52.00bp 519.50bp 39.00bp 309.50bp},clip,scale=0.840]{A_SoSe25.pdf}};
\fill[MarkerG,opacity=0.30,rounded corners=0.6bp] (53.66bp,-1.26bp) rectangle (128.52bp,-11.34bp);
\fill[MarkerG,opacity=0.30,rounded corners=0.6bp] (126.84bp,-1.26bp) rectangle (133.56bp,-11.34bp);
\end{tikzpicture}
\par\nopagebreak
\vspace{0.8mm}\nopagebreak
\begin{loesung}{A2Farbe}
\vspace{1.2mm}
\includegraphics[page=2,trim={52.00bp 722.50bp 39.00bp 85.50bp},clip,scale=0.700]{L_SoSe25.pdf}
\par
\par\vspace{1.0mm}
\includegraphics[page=2,trim={52.00bp 679.50bp 39.00bp 129.50bp},clip,scale=0.700]{L_SoSe25.pdf}
\par
\end{loesung}
\par\vspace{1.6mm}
\begin{tikzpicture}[inner sep=0]
\node[anchor=north west,inner sep=0] (bi) at (0,0) {\includegraphics[page=3,trim={52.00bp 468.50bp 39.00bp 359.50bp},clip,scale=0.840]{A_SoSe25.pdf}};
\fill[Marker,opacity=0.30,rounded corners=0.6bp] (256.20bp,-1.26bp) rectangle (291.90bp,-11.34bp);
\fill[Marker,opacity=0.30,rounded corners=0.6bp] (317.52bp,-1.26bp) rectangle (352.17bp,-11.34bp);
\end{tikzpicture}
\par
\phantomsection\label{tf:SoSe25:2:2}
\Needspace*{18mm}
\verweis{A2Farbe}{2.}{Umkehrung: welche \gs{Länge} darf der Hohlleiter bei vorgegebener Dämpfung haben?}{gleiche Frage, nur andere Zahlenwerte \textendash{} siehe \textbf{SoSe 2022, Aufgabe 2.4}, \textbf{Seite \pageref{tf:SoSe22:2:4}}}
\phantomsection\label{tf:SoSe25:2:3}
\par\vspace{1.6mm}
\Needspace*{81mm}
\begin{tikzpicture}[inner sep=0]
\node[anchor=north west,inner sep=0] (bi) at (0,0) {\includegraphics[page=3,trim={52.00bp 291.50bp 39.00bp 499.50bp},clip,scale=0.840]{A_SoSe25.pdf}};
\fill[Marker,opacity=0.30,rounded corners=0.6bp] (218.40bp,-11.34bp) rectangle (260.52bp,-21.42bp);
\fill[MarkerG,opacity=0.30,rounded corners=0.6bp] (176.80bp,-22.26bp) rectangle (197.64bp,-32.34bp);
\end{tikzpicture}
\par\nopagebreak
\vspace{0.8mm}\nopagebreak
\begin{loesung}{A2Farbe}
\vspace{1.2mm}
\includegraphics[page=2,trim={52.00bp 396.50bp 39.00bp 431.50bp},clip,scale=0.700]{L_SoSe25.pdf}
\par
\par\vspace{1.0mm}
\includegraphics[page=2,trim={52.00bp 351.50bp 39.00bp 457.50bp},clip,scale=0.700]{L_SoSe25.pdf}
\par
\par\vspace{1.0mm}
\includegraphics[page=2,trim={52.00bp 287.50bp 39.00bp 500.50bp},clip,scale=0.700]{L_SoSe25.pdf}
\par
\par\vspace{1.0mm}
\includegraphics[page=2,trim={52.00bp 263.50bp 39.00bp 564.50bp},clip,scale=0.700]{L_SoSe25.pdf}
\par
\par\vspace{1.0mm}
\includegraphics[page=2,trim={52.00bp 227.50bp 39.00bp 603.50bp},clip,scale=0.700]{L_SoSe25.pdf}
\par
\par\vspace{1.0mm}
\includegraphics[page=2,trim={52.00bp 109.50bp 39.00bp 627.50bp},clip,scale=0.700]{L_SoSe25.pdf}
\par
\end{loesung}
\par\vspace{6mm}
\clearpage
\marke{\color{A3Farbe}A3 \textendash{} Grenzflächen / Optik}
\phantomsection\label{typ:3}
\par\vspace*{-3mm}{\color{A3Farbe}\rule{\textwidth}{2.5pt}}\par\vspace{1.5mm}
{\sffamily\bfseries\huge\color{A3Farbe} A3 \textendash{} Grenzflächen / Optik}\par\vspace{1mm}
{\sffamily\color{Grau} alle 8 Klausuren \textendash{} 21 verschiedene Teilaufgaben ausgeschrieben, 3 Wiederholungen als Verweis}\par\vspace{2mm}
{\color{A3Farbe}\rule{\textwidth}{0.8pt}}\par\vspace{4mm}
{\sffamily\bfseries Register: Wenn gefragt wird \ldots}\hfill{\footnotesize\color{Grau}\colorbox{MarkerG!45}{grün} = das ist gesucht \textbf{·} rechts steht, wo die Aufgabe ausgeschrieben ist}\par\vspace{1.5mm}
\renewcommand{\arraystretch}{1.12}\footnotesize
\begin{tabularx}{\textwidth}{@{}>{\RaggedRight}X>{\RaggedRight\scriptsize}p{27mm}@{}}
\hline\rule{0pt}{3.5mm}
\gs{Beide Grenzwinkel} einer Anordnung: Totalreflexion ($P_d=0$) \emph{und} Brewster ($P_r=0$) & \textbf{S.\,\pageref{tf:WiSe20x21:3:1}} \textcolor{Grau}{WiSe 20/21, 1.} \\
\gs{Leistungen $P_r$ und $P_d$} statt Leistungs\emph{dichten} -- Achtung, der Strahlquerschnitt ändert sich! & \textbf{S.\,\pageref{tf:WiSe20x21:3:2}} \textcolor{Grau}{WiSe 20/21, 2.} \\
Fresnel an \emph{einer} Grenzfläche: \gs{reflektierte und durchgehende Leistungsdichte} \newline \hspace*{3mm}\textcolor{Grau}{\itshape \textendash{} Grundfall} & \textbf{S.\,\pageref{tf:SoSe21:3:1}} \textcolor{Grau}{SoSe 21, 1.} \\
Fresnel an \emph{einer} Grenzfläche: \gs{reflektierte und durchgehende Leistungsdichte} \newline \hspace*{3mm}\textcolor{Grau}{\itshape \textendash{} hier als Leistung, und die Polarisation steht senkrecht zur Brewster-Polarisation} & \textbf{S.\,\pageref{tf:SoSe25:3:3}} \textcolor{Grau}{SoSe 25, 3.} \\
Zweite Grenzfläche im Prisma: erst der \gs{Winkel} über die Geometrie, dann die \gs{Leistungsdichten} & \textbf{S.\,\pageref{tf:SoSe21:3:2}} \textcolor{Grau}{SoSe 21, 2.} \\
\gs{Einfallswinkel}, bei dem reflektierter und durchgehender Strahl aufeinander senkrecht stehen (= Brewster in Verkleidung) & \textbf{S.\,\pageref{tf:WiSe21x22:3:1}} \textcolor{Grau}{WiSe 21/22, 1.} \\
\gs{Maximale und minimale reflektierte Leistungsdichte} je nach Polarisationsrichtung & \textbf{S.\,\pageref{tf:WiSe21x22:3:2}} \textcolor{Grau}{WiSe 21/22, 2.} \\
\gs{$\varepsilon_r$} aus dem Brewsterwinkel (\glqq reflektiert rein linear polarisiert\grqq{} bzw. \glqq ohne Reflexionsverluste\grqq{}) \newline \hspace*{3mm}\textcolor{Grau}{\itshape \textendash{} Grundfall} & \textbf{S.\,\pageref{tf:SoSe22:3:1}} \textcolor{Grau}{SoSe 22, 1.} \\
\gs{$\varepsilon_r$} aus dem Brewsterwinkel (\glqq reflektiert rein linear polarisiert\grqq{} bzw. \glqq ohne Reflexionsverluste\grqq{}) \newline \hspace*{3mm}\textcolor{Grau}{\itshape \textendash{} hier ist der Brewsterwinkel direkt genannt} & \textbf{S.\,\pageref{tf:WiSe22x23:3:1}} \textcolor{Grau}{WiSe 22/23, 1.} \\
\gs{$\varepsilon_r$} aus dem Brewsterwinkel (\glqq reflektiert rein linear polarisiert\grqq{} bzw. \glqq ohne Reflexionsverluste\grqq{}) \newline \hspace*{3mm}\textcolor{Grau}{\itshape \textendash{} hier zusätzlich die E-Feld-Richtung einzeichnen} & \textbf{S.\,\pageref{tf:SoSe23:3:1}} \textcolor{Grau}{SoSe 23, 1.} \\
Unpolarisiertes Licht: \gs{Leistungsdichte} über 50/50 auf beide Polarisationen -- \emph{beide} Fresnel-Zeilen rechnen & \textbf{S.\,\pageref{tf:SoSe22:3:2}} \textcolor{Grau}{SoSe 22, 2.} \\
\gs{Prismenwinkel} so wählen, dass an der nächsten Grenzfläche Totalreflexion eintritt & \textbf{S.\,\pageref{tf:SoSe22:3:3}} \textcolor{Grau}{SoSe 22, 3.} \\
\gs{Gibt es einen Winkel}, bei dem der reflektierte Anteil verschwindet? (Begründung) & \textbf{S.\,\pageref{tf:SoSe22:3:4}} \textcolor{Grau}{SoSe 22, 4.} \\
\gs{Maximale einfallende Leistungsdichte} aus der zulässigen Feldstärke \emph{im} Dielektrikum & \textbf{S.\,\pageref{tf:WiSe22x23:3:2}} \textcolor{Grau}{WiSe 22/23, 2.} \\
\gs{Prismenwinkel} für reflexionsfreien \emph{Austritt} (Brewster von innen nach außen) & \textbf{S.\,\pageref{tf:WiSe22x23:3:4}} \textcolor{Grau}{WiSe 22/23, 4.} \\
Dritte Grenzfläche: \gs{Leistungsdichte} am Ende der ganzen Kette & \textbf{S.\,\pageref{tf:SoSe23:3:3}} \textcolor{Grau}{SoSe 23, 3.} \\
Welche \gs{Polarisation} hat der \emph{reflektierte} Anteil beim Brewsterwinkel? (Begründung, einzeichnen) & \textbf{S.\,\pageref{tf:WiSe23x24:3:1}} \textcolor{Grau}{WiSe 23/24, 1.} \\
Senkrechter Einfall neben Brewster-Einfall: \gs{reflektierte Leistungsdichten} beider Fälle & \textbf{S.\,\pageref{tf:WiSe23x24:3:2}} \textcolor{Grau}{WiSe 23/24, 2.} \\
Dieselben zwei Fälle, jetzt die \gs{durchgehenden Leistungsdichten} & \textbf{S.\,\pageref{tf:WiSe23x24:3:3}} \textcolor{Grau}{WiSe 23/24, 3.} \\
\gs{Einfallswinkel} \emph{und} \gs{Polarisationsrichtung} für reflexionsfreien Durchgang \newline \hspace*{3mm}\textcolor{Grau}{\itshape \textendash{} hier ist nur die Zeichnung gefragt} & \textbf{S.\,\pageref{tf:SoSe25:3:1}} \textcolor{Grau}{SoSe 25, 1.} \\
Für welche \gs{$\varepsilon_r$} werden 100\,\% reflektiert? (Totalreflexion) & \textbf{S.\,\pageref{tf:SoSe25:3:4}} \textcolor{Grau}{SoSe 25, 4.} \\
\hline\end{tabularx}\normalsize
\par\vspace{5mm}
\begin{tcolorbox}[colback=A3Farbe!4, colframe=A3Farbe!55, boxrule=0.5pt, arc=1mm, left=3mm, right=2mm, top=2mm, bottom=2mm, breakable, title={\sffamily\bfseries\small Was diesmal neu drankommen könnte \normalfont\itshape\footnotesize -- eigene Prognose (keine offizielle Ankündigung): Umkehrungen, Kombinationen und Transfers zwischen den Aufgabentypen}, coltitle=white, colbacktitle=A3Farbe!70]
\sffamily\small
\begin{itemize}\setlength{\itemsep}{2.2mm}\setlength{\leftmargini}{4mm}
\item \emph{Senkrechter} Einfall ($\alpha=0$): welcher Anteil wird reflektiert? Kam in keiner Altklausur vor, die Formel steht aber in der FS\\[0.6mm] {\footnotesize\color{A3Farbe}$\rightarrow$ $r=\frac{\sqrt{\varepsilon_{r1}}-\sqrt{\varepsilon_{r2}}}{\sqrt{\varepsilon_{r1}}+\sqrt{\varepsilon_{r2}}}$, $S_r=r^2S_h$ -- polarisationsunabhängig, weil es keine Einfallsebene gibt}
\item Winkel $\alpha_g$ der Totalreflexion direkt gefragt (bisher immer über den Umweg \glqq für welche $\varepsilon_r$\grqq)\\[0.6mm] {\footnotesize\color{A3Farbe}$\rightarrow$ $\alpha_g=\arcsin\sqrt{\varepsilon_{r2}/\varepsilon_{r1}}$, nur vom dichten ins dünne Medium. FS 6.8}
\item Prisma mit drei Grenzflächen \emph{und} unpolarisiertem Licht: Gesamtbilanz, wieviel \% der Leistung kommt hinten an?\\[0.6mm] {\footnotesize\color{A3Farbe}$\rightarrow$ je 50\,\% auf $\perp$ und $\parallel$ aufteilen, jede Grenzfläche einzeln mit ihrem $r$ rechnen, Transmissionsgrade multiplizieren}
\item Umkehrung: aus gemessenem $S_r/S_h$ auf den Einfallswinkel oder auf $\varepsilon_r$ zurückrechnen\\[0.6mm] {\footnotesize\color{A3Farbe}$\rightarrow$ aus $S_r/S_h=r^2$ den Reflexionsfaktor ziehen, Fresnel-Formel nach $q$ bzw.\ $\varepsilon_r$ auflösen}
\item Tritt der Strahl an der dritten Fläche überhaupt noch aus? (Totalreflexion prüfen, Ja/Nein mit Begründung)\\[0.6mm] {\footnotesize\color{A3Farbe}$\rightarrow$ Winkel über die Winkelsumme im Dreieck bestimmen, dann mit $\alpha_g$ vergleichen -- Antwort ist ja/nein plus Begründung}
\item Effektivwerte von $E$ und $H$ im Dielektrikum aus der \emph{durchgehenden} Leistungsdichte \textcolor{Grau}{\itshape(verwandt: WiSe 22/23 A3.2)}\\[0.6mm] {\footnotesize\color{A3Farbe}$\rightarrow$ $S=E_{eff}^2/Z$ mit $Z=Z_{F0}/\sqrt{\varepsilon_r}$ \emph{im Medium}, $H=E/Z$. K3}
\end{itemize}\end{tcolorbox}
\clearpage
\Needspace*{199.8mm}
\marke{\color{A3Farbe}A3 \textendash{} WiSe 2020/21}
\phantomsection\label{auf:WiSe20x21:3}
\aufgabenkopf{WiSe 2020/21 \textendash{} Aufgabe 3}{14 Punkte}{A3Farbe}
\begin{tikzpicture}[inner sep=0]
\node[anchor=north west,inner sep=0] (bi) at (0,0) {\includegraphics[page=4,trim={52.00bp 683.50bp 39.00bp 107.50bp},clip,scale=0.840]{A_WiSe20x21.pdf}};
\fill[Marker,opacity=0.30,rounded corners=0.6bp] (225.47bp,-1.26bp) rectangle (245.71bp,-11.34bp);
\fill[Marker,opacity=0.30,rounded corners=0.6bp] (393.01bp,-1.26bp) rectangle (410.97bp,-11.34bp);
\fill[Marker,opacity=0.30,rounded corners=0.6bp] (364.56bp,-12.18bp) rectangle (383.04bp,-22.26bp);
\end{tikzpicture}
\par\vspace{1.2mm}
\includegraphics[page=4,trim={52.00bp 414.50bp 39.00bp 209.50bp},clip,scale=0.620]{A_WiSe20x21.pdf}
\par\vspace{1.2mm}
\phantomsection\label{tf:WiSe20x21:3:1}
\par\vspace{1.6mm}
\Needspace*{42mm}
\begin{tikzpicture}[inner sep=0]
\node[anchor=north west,inner sep=0] (bi) at (0,0) {\includegraphics[page=4,trim={52.00bp 326.50bp 39.00bp 463.50bp},clip,scale=0.840]{A_WiSe20x21.pdf}};
\fill[Marker,opacity=0.30,rounded corners=0.6bp] (194.04bp,-17.22bp) rectangle (219.24bp,-27.30bp);
\fill[Marker,opacity=0.30,rounded corners=0.6bp] (183.12bp,-33.18bp) rectangle (208.04bp,-43.26bp);
\fill[MarkerG,opacity=0.30,rounded corners=0.6bp] (69.15bp,-1.26bp) rectangle (93.64bp,-11.34bp);
\end{tikzpicture}
\par\nopagebreak
\vspace{0.8mm}\nopagebreak
\begin{loesung}{A3Farbe}
\vspace{1.2mm}
\includegraphics[page=3,trim={52.00bp 696.50bp 39.00bp 108.50bp},clip,scale=0.700]{L_WiSe20x21.pdf}
\par
\par\vspace{1.0mm}
\includegraphics[page=3,trim={52.00bp 645.50bp 39.00bp 160.50bp},clip,scale=0.700]{L_WiSe20x21.pdf}
\par
\end{loesung}
\phantomsection\label{tf:WiSe20x21:3:2}
\par\vspace{1.6mm}
\Needspace*{74mm}
\begin{tikzpicture}[inner sep=0]
\node[anchor=north west,inner sep=0] (bi) at (0,0) {\includegraphics[page=4,trim={52.00bp 237.50bp 39.00bp 578.50bp},clip,scale=0.840]{A_WiSe20x21.pdf}};
\fill[Marker,opacity=0.30,rounded corners=0.6bp] (94.08bp,-1.26bp) rectangle (120.71bp,-11.34bp);
\fill[MarkerG,opacity=0.30,rounded corners=0.6bp] (144.54bp,-1.26bp) rectangle (182.68bp,-11.34bp);
\end{tikzpicture}
\par\nopagebreak
\vspace{0.8mm}\nopagebreak
\begin{loesung}{A3Farbe}
\vspace{1.2mm}
\includegraphics[page=3,trim={52.00bp 364.50bp 39.00bp 247.50bp},clip,scale=0.700]{L_WiSe20x21.pdf}
\par
\end{loesung}
\par\vspace{6mm}
\Needspace*{244.3mm}
\marke{\color{A3Farbe}A3 \textendash{} SoSe 2021}
\phantomsection\label{auf:SoSe21:3}
\aufgabenkopf{SoSe 2021 \textendash{} Aufgabe 3}{15 Punkte}{A3Farbe}
\begin{tikzpicture}[inner sep=0]
\node[anchor=north west,inner sep=0] (bi) at (0,0) {\includegraphics[page=4,trim={52.00bp 695.50bp 39.00bp 106.50bp},clip,scale=0.840]{A_SoSe21.pdf}};
\fill[Marker,opacity=0.30,rounded corners=0.6bp] (5.04bp,-13.02bp) rectangle (43.47bp,-23.10bp);
\fill[Marker,opacity=0.30,rounded corners=0.6bp] (194.04bp,-13.02bp) rectangle (218.61bp,-23.10bp);
\fill[Marker,opacity=0.30,rounded corners=0.6bp] (49.56bp,-23.94bp) rectangle (77.84bp,-34.02bp);
\end{tikzpicture}
\par\vspace{1.2mm}
\includegraphics[page=4,trim={52.00bp 385.50bp 39.00bp 169.50bp},clip,scale=0.620]{A_SoSe21.pdf}
\par\vspace{1.2mm}
\phantomsection\label{tf:SoSe21:3:1}
\par\vspace{1.6mm}
\Needspace*{91mm}
\begin{tikzpicture}[inner sep=0]
\node[anchor=north west,inner sep=0] (bi) at (0,0) {\includegraphics[page=4,trim={52.00bp 302.50bp 39.00bp 513.50bp},clip,scale=0.840]{A_SoSe21.pdf}};
\fill[MarkerG,opacity=0.30,rounded corners=0.6bp] (93.78bp,-1.26bp) rectangle (157.92bp,-11.34bp);
\fill[MarkerG,opacity=0.30,rounded corners=0.6bp] (156.24bp,-1.26bp) rectangle (162.96bp,-11.34bp);
\end{tikzpicture}
\par\nopagebreak
\vspace{0.8mm}\nopagebreak
\begin{loesung}{A3Farbe}
\vspace{1.2mm}
\includegraphics[page=3,trim={52.00bp 699.50bp 39.00bp 109.50bp},clip,scale=0.700]{L_SoSe21.pdf}
\par
\par\vspace{1.0mm}
\includegraphics[page=3,trim={52.00bp 610.50bp 39.00bp 152.50bp},clip,scale=0.700]{L_SoSe21.pdf}
\par
\par\vspace{1.0mm}
\includegraphics[page=3,trim={52.00bp 557.50bp 39.00bp 246.50bp},clip,scale=0.700]{L_SoSe21.pdf}
\par
\par\vspace{1.0mm}
\includegraphics[page=3,trim={52.00bp 515.50bp 39.00bp 294.50bp},clip,scale=0.700]{L_SoSe21.pdf}
\par
\par\vspace{1.0mm}
\includegraphics[page=3,trim={52.00bp 380.50bp 39.00bp 342.50bp},clip,scale=0.700]{L_SoSe21.pdf}
\par
\end{loesung}
\phantomsection\label{tf:SoSe21:3:2}
\par\vspace{1.6mm}
\Needspace*{52mm}
\begin{tikzpicture}[inner sep=0]
\node[anchor=north west,inner sep=0] (bi) at (0,0) {\includegraphics[page=4,trim={52.00bp 225.50bp 39.00bp 589.50bp},clip,scale=0.840]{A_SoSe21.pdf}};
\fill[Marker,opacity=0.30,rounded corners=0.6bp] (219.03bp,-12.18bp) rectangle (244.65bp,-22.26bp);
\fill[MarkerG,opacity=0.30,rounded corners=0.6bp] (57.12bp,-1.26bp) rectangle (84.00bp,-11.34bp);
\end{tikzpicture}
\par\nopagebreak
\vspace{0.8mm}\nopagebreak
\begin{loesung}{A3Farbe}
\vspace{1.2mm}
\includegraphics[page=3,trim={52.00bp 279.50bp 39.00bp 526.50bp},clip,scale=0.700]{L_SoSe21.pdf}
\par
\par\vspace{1.0mm}
\includegraphics[page=3,trim={52.00bp 164.50bp 39.00bp 572.50bp},clip,scale=0.700]{L_SoSe21.pdf}
\par
\end{loesung}
\par\vspace{6mm}
\Needspace*{225.0mm}
\marke{\color{A3Farbe}A3 \textendash{} WiSe 2021/22}
\phantomsection\label{auf:WiSe21x22:3}
\aufgabenkopf{WiSe 2021/22 \textendash{} Aufgabe 3}{13 Punkte}{A3Farbe}
\begin{tikzpicture}[inner sep=0]
\node[anchor=north west,inner sep=0] (bi) at (0,0) {\includegraphics[page=4,trim={52.00bp 708.50bp 39.00bp 105.50bp},clip,scale=0.840]{A_WiSe21x22.pdf}};
\fill[Marker,opacity=0.30,rounded corners=0.6bp] (329.28bp,-2.10bp) rectangle (370.07bp,-12.18bp);
\fill[Marker,opacity=0.30,rounded corners=0.6bp] (315.84bp,-13.86bp) rectangle (347.17bp,-23.94bp);
\end{tikzpicture}
\par\vspace{1.2mm}
\includegraphics[page=4,trim={52.00bp 448.50bp 39.00bp 187.50bp},clip,scale=0.620]{A_WiSe21x22.pdf}
\par\vspace{1.2mm}
\phantomsection\label{tf:WiSe21x22:3:1}
\par\vspace{1.6mm}
\Needspace*{65mm}
\begin{tikzpicture}[inner sep=0]
\node[anchor=north west,inner sep=0] (bi) at (0,0) {\includegraphics[page=4,trim={52.00bp 314.50bp 39.00bp 475.50bp},clip,scale=0.840]{A_WiSe21x22.pdf}};
\fill[MarkerG,opacity=0.30,rounded corners=0.6bp] (67.57bp,-1.26bp) rectangle (120.65bp,-11.34bp);
\end{tikzpicture}
\par\nopagebreak
\vspace{0.8mm}\nopagebreak
\begin{loesung}{A3Farbe}
\vspace{1.2mm}
\includegraphics[page=3,trim={52.00bp 582.50bp 39.00bp 119.50bp},clip,scale=0.700]{L_WiSe21x22.pdf}
\par
\par\vspace{1.0mm}
\includegraphics[page=3,trim={52.00bp 524.50bp 39.00bp 291.50bp},clip,scale=0.700]{L_WiSe21x22.pdf}
\par
\end{loesung}
\phantomsection\label{tf:WiSe21x22:3:2}
\par\vspace{1.6mm}
\Needspace*{82mm}
\begin{tikzpicture}[inner sep=0]
\node[anchor=north west,inner sep=0] (bi) at (0,0) {\includegraphics[page=4,trim={52.00bp 225.50bp 39.00bp 578.50bp},clip,scale=0.840]{A_WiSe21x22.pdf}};
\fill[Marker,opacity=0.30,rounded corners=0.6bp] (123.48bp,-0.42bp) rectangle (149.20bp,-10.50bp);
\fill[MarkerG,opacity=0.30,rounded corners=0.6bp] (161.64bp,-11.34bp) rectangle (220.86bp,-21.42bp);
\end{tikzpicture}
\par\nopagebreak
\vspace{0.8mm}\nopagebreak
\begin{loesung}{A3Farbe}
\vspace{1.2mm}
\includegraphics[page=3,trim={52.00bp 444.50bp 39.00bp 374.50bp},clip,scale=0.700]{L_WiSe21x22.pdf}
\par
\par\vspace{1.0mm}
\includegraphics[page=3,trim={52.00bp 364.50bp 39.00bp 430.50bp},clip,scale=0.700]{L_WiSe21x22.pdf}
\par
\par\vspace{1.0mm}
\includegraphics[page=3,trim={52.00bp 150.50bp 39.00bp 510.50bp},clip,scale=0.700]{L_WiSe21x22.pdf}
\par
\end{loesung}
\par\vspace{6mm}
\Needspace*{273.0mm}
\marke{\color{A3Farbe}A3 \textendash{} SoSe 2022}
\phantomsection\label{auf:SoSe22:3}
\aufgabenkopf{SoSe 2022 \textendash{} Aufgabe 3}{24 Punkte}{A3Farbe}
\begin{tikzpicture}[inner sep=0]
\node[anchor=north west,inner sep=0] (bi) at (0,0) {\includegraphics[page=4,trim={52.00bp 695.50bp 39.00bp 105.50bp},clip,scale=0.840]{A_SoSe22.pdf}};
\fill[Marker,opacity=0.30,rounded corners=0.6bp] (310.80bp,-2.10bp) rectangle (349.23bp,-12.18bp);
\fill[Marker,opacity=0.30,rounded corners=0.6bp] (81.48bp,-13.86bp) rectangle (111.71bp,-23.94bp);
\end{tikzpicture}
\par\vspace{1.2mm}
\includegraphics[page=4,trim={52.00bp 358.50bp 39.00bp 177.50bp},clip,scale=0.620]{A_SoSe22.pdf}
\par\vspace{1.2mm}
\phantomsection\label{tf:SoSe22:3:1}
\par\vspace{1.6mm}
\Needspace*{27mm}
\begin{tikzpicture}[inner sep=0]
\node[anchor=north west,inner sep=0] (bi) at (0,0) {\includegraphics[page=4,trim={52.00bp 315.50bp 39.00bp 513.50bp},clip,scale=0.840]{A_SoSe22.pdf}};
\fill[MarkerG,opacity=0.30,rounded corners=0.6bp] (122.44bp,-1.26bp) rectangle (193.74bp,-11.34bp);
\end{tikzpicture}
\par\nopagebreak
\vspace{0.8mm}\nopagebreak
\begin{loesung}{A3Farbe}
\vspace{1.2mm}
\includegraphics[page=3,trim={52.00bp 661.50bp 39.00bp 121.50bp},clip,scale=0.700]{L_SoSe22.pdf}
\par
\end{loesung}
\par\vspace{1.6mm}
\begin{tikzpicture}[inner sep=0]
\node[anchor=north west,inner sep=0] (bi) at (0,0) {\includegraphics[page=4,trim={52.00bp 263.50bp 39.00bp 564.50bp},clip,scale=0.840]{A_SoSe22.pdf}};
\fill[Marker,opacity=0.30,rounded corners=0.6bp] (87.25bp,-1.26bp) rectangle (105.21bp,-11.34bp);
\end{tikzpicture}
\par
\phantomsection\label{tf:SoSe22:3:2}
\par\vspace{1.6mm}
\Needspace*{75mm}
\begin{tikzpicture}[inner sep=0]
\node[anchor=north west,inner sep=0] (bi) at (0,0) {\includegraphics[page=4,trim={52.00bp 238.50bp 39.00bp 589.50bp},clip,scale=0.840]{A_SoSe22.pdf}};
\fill[MarkerG,opacity=0.30,rounded corners=0.6bp] (51.39bp,-1.26bp) rectangle (116.76bp,-11.34bp);
\fill[MarkerG,opacity=0.30,rounded corners=0.6bp] (115.08bp,-1.26bp) rectangle (121.80bp,-11.34bp);
\end{tikzpicture}
\par\nopagebreak
\vspace{0.8mm}\nopagebreak
\begin{loesung}{A3Farbe}
\vspace{1.2mm}
\includegraphics[page=3,trim={52.00bp 631.50bp 39.00bp 197.50bp},clip,scale=0.700]{L_SoSe22.pdf}
\par
\par\vspace{1.0mm}
\includegraphics[page=3,trim={52.00bp 382.50bp 39.00bp 220.50bp},clip,scale=0.700]{L_SoSe22.pdf}
\par
\end{loesung}
\phantomsection\label{tf:SoSe22:3:3}
\par\vspace{1.6mm}
\Needspace*{35mm}
\begin{tikzpicture}[inner sep=0]
\node[anchor=north west,inner sep=0] (bi) at (0,0) {\includegraphics[page=4,trim={52.00bp 187.50bp 39.00bp 641.50bp},clip,scale=0.840]{A_SoSe22.pdf}};
\fill[MarkerG,opacity=0.30,rounded corners=0.6bp] (69.15bp,-1.26bp) rectangle (93.64bp,-11.34bp);
\end{tikzpicture}
\par\nopagebreak
\vspace{0.8mm}\nopagebreak
\begin{loesung}{A3Farbe}
\vspace{1.2mm}
\includegraphics[page=3,trim={52.00bp 247.50bp 39.00bp 503.50bp},clip,scale=0.700]{L_SoSe22.pdf}
\par
\end{loesung}
\phantomsection\label{tf:SoSe22:3:4}
\par\vspace{1.6mm}
\Needspace*{29mm}
\begin{tikzpicture}[inner sep=0]
\node[anchor=north west,inner sep=0] (bi) at (0,0) {\includegraphics[page=4,trim={52.00bp 110.50bp 39.00bp 693.50bp},clip,scale=0.840]{A_SoSe22.pdf}};
\fill[MarkerG,opacity=0.30,rounded corners=0.6bp] (77.28bp,-0.42bp) rectangle (102.00bp,-10.50bp);
\end{tikzpicture}
\par\nopagebreak
\vspace{0.8mm}\nopagebreak
\begin{loesung}{A3Farbe}
\vspace{1.2mm}
\includegraphics[page=3,trim={52.00bp 177.50bp 39.00bp 627.50bp},clip,scale=0.700]{L_SoSe22.pdf}
\par
\end{loesung}
\par\vspace{6mm}
\Needspace*{252.4mm}
\marke{\color{A3Farbe}A3 \textendash{} WiSe 2022/23}
\phantomsection\label{auf:WiSe22x23:3}
\aufgabenkopf{WiSe 2022/23 \textendash{} Aufgabe 3}{26 Punkte}{A3Farbe}
\begin{tikzpicture}[inner sep=0]
\node[anchor=north west,inner sep=0] (bi) at (0,0) {\includegraphics[page=4,trim={52.00bp 707.50bp 39.00bp 106.50bp},clip,scale=0.840]{A_WiSe22x23.pdf}};
\fill[Marker,opacity=0.30,rounded corners=0.6bp] (10.08bp,-13.02bp) rectangle (35.70bp,-23.10bp);
\fill[Marker,opacity=0.30,rounded corners=0.6bp] (332.30bp,-13.02bp) rectangle (349.44bp,-23.10bp);
\end{tikzpicture}
\par\vspace{1.2mm}
\includegraphics[page=4,trim={52.00bp 361.50bp 39.00bp 161.50bp},clip,scale=0.620]{A_WiSe22x23.pdf}
\par\vspace{1.2mm}
\phantomsection\label{tf:WiSe22x23:3:1}
\par\vspace{1.6mm}
\Needspace*{22mm}
\begin{tikzpicture}[inner sep=0]
\node[anchor=north west,inner sep=0] (bi) at (0,0) {\includegraphics[page=4,trim={52.00bp 327.50bp 39.00bp 500.50bp},clip,scale=0.840]{A_WiSe22x23.pdf}};
\fill[MarkerG,opacity=0.30,rounded corners=0.6bp] (122.44bp,-1.26bp) rectangle (193.74bp,-11.34bp);
\end{tikzpicture}
\par\nopagebreak
\vspace{0.8mm}\nopagebreak
\begin{loesung}{A3Farbe}
\vspace{1.2mm}
\includegraphics[page=3,trim={52.00bp 686.50bp 39.00bp 121.50bp},clip,scale=0.700]{L_WiSe22x23.pdf}
\par
\end{loesung}
\par\vspace{1.6mm}
\begin{tikzpicture}[inner sep=0]
\node[anchor=north west,inner sep=0] (bi) at (0,0) {\includegraphics[page=4,trim={52.00bp 276.50bp 39.00bp 551.50bp},clip,scale=0.840]{A_WiSe22x23.pdf}};
\fill[Marker,opacity=0.30,rounded corners=0.6bp] (73.08bp,-1.26bp) rectangle (105.84bp,-11.34bp);
\end{tikzpicture}
\par
\phantomsection\label{tf:WiSe22x23:3:2}
\par\vspace{1.6mm}
\Needspace*{69mm}
\begin{tikzpicture}[inner sep=0]
\node[anchor=north west,inner sep=0] (bi) at (0,0) {\includegraphics[page=4,trim={52.00bp 238.50bp 39.00bp 577.50bp},clip,scale=0.840]{A_WiSe22x23.pdf}};
\fill[Marker,opacity=0.30,rounded corners=0.6bp] (188.66bp,-11.34bp) rectangle (217.32bp,-21.42bp);
\fill[MarkerG,opacity=0.30,rounded corners=0.6bp] (88.62bp,-1.26bp) rectangle (155.40bp,-11.34bp);
\fill[MarkerG,opacity=0.30,rounded corners=0.6bp] (153.72bp,-1.26bp) rectangle (160.44bp,-11.34bp);
\end{tikzpicture}
\par\nopagebreak
\vspace{0.8mm}\nopagebreak
\begin{loesung}{A3Farbe}
\vspace{1.2mm}
\includegraphics[page=3,trim={52.00bp 632.50bp 39.00bp 197.50bp},clip,scale=0.700]{L_WiSe22x23.pdf}
\par
\par\vspace{1.0mm}
\includegraphics[page=3,trim={52.00bp 458.50bp 39.00bp 220.50bp},clip,scale=0.700]{L_WiSe22x23.pdf}
\par
\par\vspace{1.0mm}
\includegraphics[page=3,trim={52.00bp 411.50bp 39.00bp 393.50bp},clip,scale=0.700]{L_WiSe22x23.pdf}
\par
\end{loesung}
\phantomsection\label{tf:WiSe22x23:3:3}
\Needspace*{18mm}
\verweis{A3Farbe}{3.}{Zweite Grenzfläche im Prisma: erst der \gs{Winkel} über die Geometrie, dann die \gs{Leistungsdichten}}{gleiche Frage, nur andere Zahlenwerte \textendash{} siehe \textbf{SoSe 2021, Aufgabe 3.2}, \textbf{Seite \pageref{tf:SoSe21:3:2}}}
\phantomsection\label{tf:WiSe22x23:3:4}
\par\vspace{1.6mm}
\Needspace*{36mm}
\begin{tikzpicture}[inner sep=0]
\node[anchor=north west,inner sep=0] (bi) at (0,0) {\includegraphics[page=4,trim={52.00bp 98.50bp 39.00bp 717.50bp},clip,scale=0.840]{A_WiSe22x23.pdf}};
\fill[MarkerG,opacity=0.30,rounded corners=0.6bp] (69.15bp,-1.26bp) rectangle (93.64bp,-11.34bp);
\end{tikzpicture}
\par\nopagebreak
\vspace{0.8mm}\nopagebreak
\begin{loesung}{A3Farbe}
\vspace{1.2mm}
\includegraphics[page=3,trim={52.00bp 177.50bp 39.00bp 653.50bp},clip,scale=0.700]{L_WiSe22x23.pdf}
\par
\par\vspace{1.0mm}
\includegraphics[page=3,trim={52.00bp 97.50bp 39.00bp 675.50bp},clip,scale=0.700]{L_WiSe22x23.pdf}
\par
\end{loesung}
\par\vspace{6mm}
\Needspace*{243.3mm}
\marke{\color{A3Farbe}A3 \textendash{} SoSe 2023}
\phantomsection\label{auf:SoSe23:3}
\aufgabenkopf{SoSe 2023 \textendash{} Aufgabe 3}{25 Punkte}{A3Farbe}
\begin{tikzpicture}[inner sep=0]
\node[anchor=north west,inner sep=0] (bi) at (0,0) {\includegraphics[page=4,trim={52.00bp 709.50bp 39.00bp 106.50bp},clip,scale=0.840]{A_SoSe23.pdf}};
\fill[Marker,opacity=0.30,rounded corners=0.6bp] (309.96bp,-1.26bp) rectangle (344.82bp,-11.34bp);
\fill[Marker,opacity=0.30,rounded corners=0.6bp] (85.97bp,-13.02bp) rectangle (102.23bp,-23.10bp);
\end{tikzpicture}
\par\vspace{1.2mm}
\includegraphics[page=4,trim={52.00bp 460.50bp 39.00bp 191.50bp},clip,scale=0.620]{A_SoSe23.pdf}
\par\vspace{1.2mm}
\phantomsection\label{tf:SoSe23:3:1}
\par\vspace{1.6mm}
\Needspace*{45mm}
\begin{tikzpicture}[inner sep=0]
\node[anchor=north west,inner sep=0] (bi) at (0,0) {\includegraphics[page=4,trim={52.00bp 302.50bp 39.00bp 499.50bp},clip,scale=0.840]{A_SoSe23.pdf}};
\fill[Marker,opacity=0.30,rounded corners=0.6bp] (233.28bp,-1.26bp) rectangle (244.32bp,-11.34bp);
\fill[MarkerG,opacity=0.30,rounded corners=0.6bp] (61.90bp,-12.18bp) rectangle (94.29bp,-22.26bp);
\end{tikzpicture}
\par\nopagebreak
\vspace{0.8mm}\nopagebreak
\begin{loesung}{A3Farbe}
\vspace{1.2mm}
\includegraphics[page=3,trim={52.00bp 655.50bp 39.00bp 89.50bp},clip,scale=0.700]{L_SoSe23.pdf}
\par
\end{loesung}
\par\vspace{1.6mm}
\begin{tikzpicture}[inner sep=0]
\node[anchor=north west,inner sep=0] (bi) at (0,0) {\includegraphics[page=4,trim={52.00bp 238.50bp 39.00bp 589.50bp},clip,scale=0.840]{A_SoSe23.pdf}};
\fill[Marker,opacity=0.30,rounded corners=0.6bp] (73.08bp,-1.26bp) rectangle (101.85bp,-11.34bp);
\end{tikzpicture}
\par
\phantomsection\label{tf:SoSe23:3:2}
\Needspace*{18mm}
\verweis{A3Farbe}{2.}{\gs{Prismenwinkel} so wählen, dass an der nächsten Grenzfläche Totalreflexion eintritt}{gleiche Frage, nur andere Zahlenwerte \textendash{} siehe \textbf{SoSe 2022, Aufgabe 3.3}, \textbf{Seite \pageref{tf:SoSe22:3:3}}}
\phantomsection\label{tf:SoSe23:3:3}
\par\vspace{1.6mm}
\Needspace*{105mm}
\begin{tikzpicture}[inner sep=0]
\node[anchor=north west,inner sep=0] (bi) at (0,0) {\includegraphics[page=4,trim={52.00bp 124.50bp 39.00bp 692.50bp},clip,scale=0.840]{A_SoSe23.pdf}};
\fill[MarkerG,opacity=0.30,rounded corners=0.6bp] (118.02bp,-0.42bp) rectangle (134.82bp,-10.50bp);
\end{tikzpicture}
\par\nopagebreak
\vspace{0.8mm}\nopagebreak
\begin{loesung}{A3Farbe}
\vspace{1.2mm}
\includegraphics[page=3,trim={52.00bp 342.50bp 39.00bp 368.50bp},clip,scale=0.700]{L_SoSe23.pdf}
\par
\par\vspace{1.0mm}
\includegraphics[page=3,trim={52.00bp 238.50bp 39.00bp 509.50bp},clip,scale=0.700]{L_SoSe23.pdf}
\par
\par\vspace{1.0mm}
\includegraphics[page=3,trim={52.00bp 176.50bp 39.00bp 615.50bp},clip,scale=0.700]{L_SoSe23.pdf}
\par
\par\vspace{1.0mm}
\includegraphics[page=3,trim={52.00bp 116.50bp 39.00bp 675.50bp},clip,scale=0.700]{L_SoSe23.pdf}
\par
\par\vspace{1.0mm}
\includegraphics[page=3,trim={52.00bp 73.50bp 39.00bp 735.50bp},clip,scale=0.700]{L_SoSe23.pdf}
\par
\end{loesung}
\par\vspace{6mm}
\Needspace*{273.0mm}
\marke{\color{A3Farbe}A3 \textendash{} WiSe 2023/24}
\phantomsection\label{auf:WiSe23x24:3}
\aufgabenkopf{WiSe 2023/24 \textendash{} Aufgabe 3}{27 Punkte}{A3Farbe}
\begin{tikzpicture}[inner sep=0]
\node[anchor=north west,inner sep=0] (bi) at (0,0) {\includegraphics[page=4,trim={52.00bp 708.50bp 39.00bp 105.50bp},clip,scale=0.840]{A_WiSe23x24.pdf}};
\fill[Marker,opacity=0.30,rounded corners=0.6bp] (232.68bp,-2.10bp) rectangle (276.36bp,-12.18bp);
\fill[Marker,opacity=0.30,rounded corners=0.6bp] (100.80bp,-13.02bp) rectangle (131.44bp,-23.10bp);
\end{tikzpicture}
\par\vspace{1.2mm}
\includegraphics[page=4,trim={52.00bp 480.50bp 39.00bp 165.50bp},clip,scale=0.620]{A_WiSe23x24.pdf}
\par\vspace{1.2mm}
\includegraphics[page=4,trim={52.00bp 430.50bp 39.00bp 390.50bp},clip,scale=0.840]{A_WiSe23x24.pdf}
\par\vspace{1.2mm}
\includegraphics[page=4,trim={52.00bp 366.50bp 39.00bp 461.50bp},clip,scale=0.840]{A_WiSe23x24.pdf}
\par\vspace{1.2mm}
\phantomsection\label{tf:WiSe23x24:3:1}
\par\vspace{1.6mm}
\Needspace*{58mm}
\begin{tikzpicture}[inner sep=0]
\node[anchor=north west,inner sep=0] (bi) at (0,0) {\includegraphics[page=4,trim={52.00bp 303.50bp 39.00bp 512.50bp},clip,scale=0.840]{A_WiSe23x24.pdf}};
\fill[MarkerG,opacity=0.30,rounded corners=0.6bp] (48.96bp,-0.42bp) rectangle (94.25bp,-10.50bp);
\end{tikzpicture}
\par\nopagebreak
\vspace{0.8mm}\nopagebreak
\begin{loesung}{A3Farbe}
\vspace{1.2mm}
\includegraphics[page=3,trim={52.00bp 588.50bp 39.00bp 87.50bp},clip,scale=0.700]{L_WiSe23x24.pdf}
\par
\end{loesung}
\phantomsection\label{tf:WiSe23x24:3:2}
\par\vspace{1.6mm}
\Needspace*{107mm}
\begin{tikzpicture}[inner sep=0]
\node[anchor=north west,inner sep=0] (bi) at (0,0) {\includegraphics[page=4,trim={52.00bp 252.50bp 39.00bp 575.50bp},clip,scale=0.840]{A_WiSe23x24.pdf}};
\fill[MarkerG,opacity=0.30,rounded corners=0.6bp] (51.45bp,-1.26bp) rectangle (120.96bp,-11.34bp);
\fill[MarkerG,opacity=0.30,rounded corners=0.6bp] (119.28bp,-1.26bp) rectangle (126.00bp,-11.34bp);
\end{tikzpicture}
\par\nopagebreak
\vspace{0.8mm}\nopagebreak
\begin{loesung}{A3Farbe}
\vspace{1.2mm}
\includegraphics[page=3,trim={52.00bp 444.50bp 39.00bp 316.50bp},clip,scale=0.700]{L_WiSe23x24.pdf}
\par
\par\vspace{1.0mm}
\includegraphics[page=3,trim={52.00bp 386.50bp 39.00bp 429.50bp},clip,scale=0.700]{L_WiSe23x24.pdf}
\par
\par\vspace{1.0mm}
\includegraphics[page=3,trim={52.00bp 90.50bp 39.00bp 476.50bp},clip,scale=0.700]{L_WiSe23x24.pdf}
\par
\end{loesung}
\phantomsection\label{tf:WiSe23x24:3:3}
\par\vspace{1.6mm}
\Needspace*{145mm}
\begin{tikzpicture}[inner sep=0]
\node[anchor=north west,inner sep=0] (bi) at (0,0) {\includegraphics[page=4,trim={52.00bp 189.50bp 39.00bp 625.50bp},clip,scale=0.840]{A_WiSe23x24.pdf}};
\fill[MarkerG,opacity=0.30,rounded corners=0.6bp] (94.37bp,-1.26bp) rectangle (162.96bp,-11.34bp);
\fill[MarkerG,opacity=0.30,rounded corners=0.6bp] (161.28bp,-1.26bp) rectangle (168.00bp,-11.34bp);
\end{tikzpicture}
\par\nopagebreak
\vspace{0.8mm}\nopagebreak
\begin{loesung}{A3Farbe}
\vspace{1.2mm}
\includegraphics[page=4,trim={52.00bp 760.50bp 39.00bp 48.50bp},clip,scale=0.700]{L_WiSe23x24.pdf}
\par
\par\vspace{1.0mm}
\includegraphics[page=4,trim={52.00bp 738.50bp 39.00bp 91.50bp},clip,scale=0.700]{L_WiSe23x24.pdf}
\par
\par\vspace{1.0mm}
\includegraphics[page=4,trim={52.00bp 676.50bp 39.00bp 115.50bp},clip,scale=0.700]{L_WiSe23x24.pdf}
\par
\par\vspace{1.0mm}
\includegraphics[page=4,trim={52.00bp 634.50bp 39.00bp 175.50bp},clip,scale=0.700]{L_WiSe23x24.pdf}
\par
\par\vspace{1.0mm}
\includegraphics[page=4,trim={52.00bp 301.50bp 39.00bp 227.50bp},clip,scale=0.700]{L_WiSe23x24.pdf}
\par
\par\vspace{1.0mm}
\includegraphics[page=4,trim={52.00bp 265.50bp 39.00bp 562.50bp},clip,scale=0.700]{L_WiSe23x24.pdf}
\par
\par\vspace{1.0mm}
\includegraphics[page=4,trim={52.00bp 220.50bp 39.00bp 590.50bp},clip,scale=0.700]{L_WiSe23x24.pdf}
\par
\par\vspace{1.0mm}
\includegraphics[page=4,trim={52.00bp 182.50bp 39.00bp 631.50bp},clip,scale=0.700]{L_WiSe23x24.pdf}
\par
\par\vspace{1.0mm}
\includegraphics[page=4,trim={52.00bp 139.50bp 39.00bp 669.50bp},clip,scale=0.700]{L_WiSe23x24.pdf}
\par
\end{loesung}
\par\vspace{6mm}
\Needspace*{249.6mm}
\marke{\color{A3Farbe}A3 \textendash{} SoSe 2025}
\phantomsection\label{auf:SoSe25:3}
\aufgabenkopf{SoSe 2025 \textendash{} Aufgabe 3}{16 Punkte}{A3Farbe}
\begin{tikzpicture}[inner sep=0]
\node[anchor=north west,inner sep=0] (bi) at (0,0) {\includegraphics[page=4,trim={52.00bp 707.50bp 39.00bp 106.50bp},clip,scale=0.840]{A_SoSe25.pdf}};
\fill[Marker,opacity=0.30,rounded corners=0.6bp] (127.68bp,-1.26bp) rectangle (164.95bp,-11.34bp);
\fill[Marker,opacity=0.30,rounded corners=0.6bp] (201.60bp,-13.02bp) rectangle (228.23bp,-23.10bp);
\end{tikzpicture}
\par\vspace{1.2mm}
\includegraphics[page=4,trim={52.00bp 461.50bp 39.00bp 177.50bp},clip,scale=0.620]{A_SoSe25.pdf}
\par\vspace{1.2mm}
\phantomsection\label{tf:SoSe25:3:1}
\par\vspace{1.6mm}
\Needspace*{42mm}
\begin{tikzpicture}[inner sep=0]
\node[anchor=north west,inner sep=0] (bi) at (0,0) {\includegraphics[page=4,trim={52.00bp 404.50bp 39.00bp 411.50bp},clip,scale=0.840]{A_SoSe25.pdf}};
\fill[MarkerG,opacity=0.30,rounded corners=0.6bp] (119.12bp,-1.26bp) rectangle (135.51bp,-11.34bp);
\end{tikzpicture}
\par\nopagebreak
\vspace{0.8mm}\nopagebreak
\begin{loesung}{A3Farbe}
\vspace{1.2mm}
\includegraphics[page=3,trim={52.00bp 642.50bp 39.00bp 98.50bp},clip,scale=0.700]{L_SoSe25.pdf}
\par
\end{loesung}
\phantomsection\label{tf:SoSe25:3:2}
\Needspace*{18mm}
\verweis{A3Farbe}{2.}{\gs{$\varepsilon_r$} aus dem Brewsterwinkel (\glqq reflektiert rein linear polarisiert\grqq{} bzw. \glqq ohne Reflexionsverluste\grqq{})}{gleiche Frage, nur andere Zahlenwerte \textendash{} siehe \textbf{SoSe 2022, Aufgabe 3.1}, \textbf{Seite \pageref{tf:SoSe22:3:1}}}
\phantomsection\label{tf:SoSe25:3:3}
\par\vspace{1.6mm}
\Needspace*{60mm}
\begin{tikzpicture}[inner sep=0]
\node[anchor=north west,inner sep=0] (bi) at (0,0) {\includegraphics[page=4,trim={52.00bp 287.50bp 39.00bp 500.50bp},clip,scale=0.840]{A_SoSe25.pdf}};
\fill[Marker,opacity=0.30,rounded corners=0.6bp] (175.56bp,-12.18bp) rectangle (202.02bp,-22.26bp);
\fill[Marker,opacity=0.30,rounded corners=0.6bp] (256.20bp,-34.86bp) rectangle (285.60bp,-44.94bp);
\fill[MarkerG,opacity=0.30,rounded corners=0.6bp] (52.18bp,-23.10bp) rectangle (90.72bp,-33.18bp);
\fill[MarkerG,opacity=0.30,rounded corners=0.6bp] (89.04bp,-23.10bp) rectangle (96.60bp,-33.18bp);
\end{tikzpicture}
\par\nopagebreak
\vspace{0.8mm}\nopagebreak
\begin{loesung}{A3Farbe}
\vspace{1.2mm}
\includegraphics[page=3,trim={52.00bp 455.50bp 39.00bp 244.50bp},clip,scale=0.700]{L_SoSe25.pdf}
\par
\end{loesung}
\par\vspace{1.6mm}
\begin{tikzpicture}[inner sep=0]
\node[anchor=north west,inner sep=0] (bi) at (0,0) {\includegraphics[page=4,trim={52.00bp 236.50bp 39.00bp 578.50bp},clip,scale=0.840]{A_SoSe25.pdf}};
\fill[Marker,opacity=0.30,rounded corners=0.6bp] (95.76bp,-13.02bp) rectangle (121.38bp,-23.10bp);
\end{tikzpicture}
\par
\phantomsection\label{tf:SoSe25:3:4}
\par\vspace{1.6mm}
\Needspace*{45mm}
\begin{tikzpicture}[inner sep=0]
\node[anchor=north west,inner sep=0] (bi) at (0,0) {\includegraphics[page=4,trim={52.00bp 197.50bp 39.00bp 618.50bp},clip,scale=0.840]{A_SoSe25.pdf}};
\fill[MarkerG,opacity=0.30,rounded corners=0.6bp] (67.96bp,-1.26bp) rectangle (89.84bp,-11.34bp);
\end{tikzpicture}
\par\nopagebreak
\vspace{0.8mm}\nopagebreak
\begin{loesung}{A3Farbe}
\vspace{1.2mm}
\includegraphics[page=3,trim={52.00bp 337.50bp 39.00bp 447.50bp},clip,scale=0.700]{L_SoSe25.pdf}
\par
\par\vspace{1.0mm}
\includegraphics[page=3,trim={52.00bp 296.50bp 39.00bp 514.50bp},clip,scale=0.700]{L_SoSe25.pdf}
\par
\par\vspace{1.0mm}
\includegraphics[page=3,trim={52.00bp 255.50bp 39.00bp 559.50bp},clip,scale=0.700]{L_SoSe25.pdf}
\par
\end{loesung}
\par\vspace{6mm}
\clearpage
\marke{\color{A4Farbe}A4 \textendash{} Leitungstransformation}
\phantomsection\label{typ:4}
\par\vspace*{-3mm}{\color{A4Farbe}\rule{\textwidth}{2.5pt}}\par\vspace{1.5mm}
{\sffamily\bfseries\huge\color{A4Farbe} A4 \textendash{} Leitungstransformation}\par\vspace{1mm}
{\sffamily\color{Grau} alle 8 Klausuren \textendash{} 17 verschiedene Teilaufgaben ausgeschrieben, 5 Wiederholungen als Verweis}\par\vspace{2mm}
{\color{A4Farbe}\rule{\textwidth}{0.8pt}}\par\vspace{4mm}
{\sffamily\bfseries Register: Wenn gefragt wird \ldots}\hfill{\footnotesize\color{Grau}\colorbox{MarkerG!45}{grün} = das ist gesucht \textbf{·} rechts steht, wo die Aufgabe ausgeschrieben ist}\par\vspace{1.5mm}
\renewcommand{\arraystretch}{1.12}\footnotesize
\begin{tabularx}{\textwidth}{@{}>{\RaggedRight}X>{\RaggedRight\scriptsize}p{27mm}@{}}
\hline\rule{0pt}{3.5mm}
\gs{Längen} und \gs{Wellenwiderstand} so wählen, dass der Eingang reell wird (zwei parallel gespeiste Antennen) & \textbf{S.\,\pageref{tf:WiSe20x21:4:1}} \textcolor{Grau}{WiSe 20/21, 1.} \\
\gs{Eingangsimpedanz}, wenn einer der beiden Zweige kurzgeschlossen wird & \textbf{S.\,\pageref{tf:WiSe20x21:4:2}} \textcolor{Grau}{WiSe 20/21, 2.} \\
\gs{Kompensationsbauteil} ($C$ oder $L$, seriell oder parallel) so wählen, dass die Impedanz \emph{rein reell} wird \newline \hspace*{3mm}\textcolor{Grau}{\itshape \textendash{} Grundfall} & \textbf{S.\,\pageref{tf:SoSe21:4:1}} \textcolor{Grau}{SoSe 21, 1.} \\
\gs{Kompensationsbauteil} ($C$ oder $L$, seriell oder parallel) so wählen, dass die Impedanz \emph{rein reell} wird \newline \hspace*{3mm}\textcolor{Grau}{\itshape \textendash{} hier eine Induktivität statt einer Kapazität} & \textbf{S.\,\pageref{tf:SoSe23:4:2}} \textcolor{Grau}{SoSe 23, 2.} \\
$\lambda/4$-Transformator über das \emph{Dielektrikum}: \gs{$Z_L$} und \gs{$\varepsilon_r$} wählen & \textbf{S.\,\pageref{tf:SoSe21:4:2}} \textcolor{Grau}{SoSe 21, 2.} \\
\gs{Leitungslänge} für einen geforderten Phasenunterschied zwischen zwei Antennen & \textbf{S.\,\pageref{tf:WiSe21x22:4:1}} \textcolor{Grau}{WiSe 21/22, 1.} \\
\gs{Kürzeste Länge}, für die der Eingangswiderstand einen vorgegebenen Wert annimmt & \textbf{S.\,\pageref{tf:WiSe21x22:4:2}} \textcolor{Grau}{WiSe 21/22, 2.} \\
\gs{Impedanz} über eine Leitung nach vorn transformieren (Leitungsgleichung) & \textbf{S.\,\pageref{tf:SoSe22:4:1}} \textcolor{Grau}{SoSe 22, 1.} \\
$\lambda/4$-Transformator entwerfen: \gs{Länge $l$} und \gs{Wellenwiderstand $Z_L$} für ein gefordertes $Z$ \newline \hspace*{3mm}\textcolor{Grau}{\itshape \textendash{} mit vorgegebenem Längenfenster 1,5--2,0\,m} & \textbf{S.\,\pageref{tf:SoSe22:4:3}} \textcolor{Grau}{SoSe 22, 3.} \\
$\lambda/4$-Transformator entwerfen: \gs{Länge $l$} und \gs{Wellenwiderstand $Z_L$} für ein gefordertes $Z$ \newline \hspace*{3mm}\textcolor{Grau}{\itshape \textendash{} Grundfall} & \textbf{S.\,\pageref{tf:WiSe22x23:4:3}} \textcolor{Grau}{WiSe 22/23, 3.} \\
\gs{Reflexionsfaktor-Betrag} am Eingang nach einem Kurzschluss (Begründung ohne Rechnung) & \textbf{S.\,\pageref{tf:SoSe22:4:4}} \textcolor{Grau}{SoSe 22, 4.} \\
\gs{Kürzeste Länge}, damit der Reflexionsfaktor \emph{reell} wird & \textbf{S.\,\pageref{tf:WiSe22x23:4:2}} \textcolor{Grau}{WiSe 22/23, 2.} \\
Stichleitung: \gs{Länge} so wählen, dass die Impedanz rein reell wird & \textbf{S.\,\pageref{tf:SoSe23:4:1}} \textcolor{Grau}{SoSe 23, 1.} \\
Rückwärts durch die Kette: welches \gs{$Z$} muss man an dieser Stelle sehen? & \textbf{S.\,\pageref{tf:WiSe23x24:4:2}} \textcolor{Grau}{WiSe 23/24, 2.} \\
Rückwärts plus Bauteilwert: welches \gs{$Z$} und welches \gs{$C$} sind nötig? & \textbf{S.\,\pageref{tf:WiSe23x24:4:3}} \textcolor{Grau}{WiSe 23/24, 3.} \\
\gs{Impedanz} bei in Serie \emph{eingeschleifter}, kurzgeschlossener Stichleitung & \textbf{S.\,\pageref{tf:SoSe25:4:1}} \textcolor{Grau}{SoSe 25, 1.} \\
\gs{Bauteilwerte $R$ und $L$} aus der transformierten Impedanz zurückrechnen ($\lambda/4$-Inverter) & \textbf{S.\,\pageref{tf:SoSe25:4:2}} \textcolor{Grau}{SoSe 25, 2.} \\
\hline\end{tabularx}\normalsize
\par\vspace{5mm}
\begin{tcolorbox}[colback=A4Farbe!4, colframe=A4Farbe!55, boxrule=0.5pt, arc=1mm, left=3mm, right=2mm, top=2mm, bottom=2mm, breakable, title={\sffamily\bfseries\small Was diesmal neu drankommen könnte \normalfont\itshape\footnotesize -- eigene Prognose (keine offizielle Ankündigung): Umkehrungen, Kombinationen und Transfers zwischen den Aufgabentypen}, coltitle=white, colbacktitle=A4Farbe!70]
\sffamily\small
\begin{itemize}\setlength{\itemsep}{2.2mm}\setlength{\leftmargini}{4mm}
\item Stichleitungsanpassung komplett: \emph{Position und} Länge der Stichleitung bestimmen -- bisher war die Position immer vorgegeben\\[0.6mm] {\footnotesize\color{A4Farbe}$\rightarrow$ mit Leitwerten rechnen (Stichleitung liegt parallel): Position dort, wo $\mathrm{Re}\{Y\}=1/Z_L$, Stichleitung kompensiert dann $\mathrm{Im}\{Y\}$}
\item $\lambda/2$-Leitung als Sonderfall (Impedanz unverändert) bzw.\ Umschlag Leerlauf $\leftrightarrow$ Kurzschluss bei $\lambda/4$\\[0.6mm] {\footnotesize\color{A4Farbe}$\rightarrow$ $\lambda/2$: $Z$ bleibt unverändert. $\lambda/4$: $Z_{ein}=Z_L^2/Z_{ab}$, Leerlauf $\leftrightarrow$ Kurzschluss. K4}
\item Reflexionsfaktor am Eingang wirklich \emph{ausrechnen} (SoSe 22 A4.4 verlangte nur eine Begründung)\\[0.6mm] {\footnotesize\color{A4Farbe}$\rightarrow$ $r=\frac{Z-Z_L}{Z+Z_L}$, daraus $|r|$ und ggf.\ $s=\frac{1+|r|}{1-|r|}$}
\item Leistungsbilanz: welche Leistung kommt an der Last an, welche wird reflektiert?\\[0.6mm] {\footnotesize\color{A4Farbe}$\rightarrow$ $P_{ab}=P_{hin}(1-|r|^2)$ -- der Rest läuft zurück}
\item Koaxialleitung statt Zweidrahtleitung: Durchmesserverhältnis $D/d$ aus $Z_L$ -- Brücke zu Aufgabe 1\\[0.6mm] {\footnotesize\color{A4Farbe}$\rightarrow$ $Z_L=\frac{60\,\Omega}{\sqrt{\varepsilon_r}}\ln\frac{D}{d}$ nach $D/d$ auflösen (statt der Zweidraht-Formel). K1}
\end{itemize}\end{tcolorbox}
\clearpage
\Needspace*{182.0mm}
\marke{\color{A4Farbe}A4 \textendash{} WiSe 2020/21}
\phantomsection\label{auf:WiSe20x21:4}
\aufgabenkopf{WiSe 2020/21 \textendash{} Aufgabe 4}{12 Punkte}{A4Farbe}
\begin{tikzpicture}[inner sep=0]
\node[anchor=north west,inner sep=0] (bi) at (0,0) {\includegraphics[page=5,trim={52.00bp 694.50bp 39.00bp 107.50bp},clip,scale=0.840]{A_WiSe20x21.pdf}};
\fill[Marker,opacity=0.30,rounded corners=0.6bp] (197.06bp,-1.26bp) rectangle (214.20bp,-11.34bp);
\fill[Marker,opacity=0.30,rounded corners=0.6bp] (376.24bp,-12.18bp) rectangle (404.65bp,-22.26bp);
\fill[Marker,opacity=0.30,rounded corners=0.6bp] (173.54bp,-23.10bp) rectangle (190.68bp,-33.18bp);
\end{tikzpicture}
\par\vspace{1.2mm}
\includegraphics[page=5,trim={52.00bp 479.50bp 39.00bp 165.50bp},clip,scale=0.620]{A_WiSe20x21.pdf}
\par\vspace{1.2mm}
\phantomsection\label{tf:WiSe20x21:4:1}
\par\vspace{1.6mm}
\Needspace*{46mm}
\begin{tikzpicture}[inner sep=0]
\node[anchor=north west,inner sep=0] (bi) at (0,0) {\includegraphics[page=5,trim={52.00bp 377.50bp 39.00bp 437.50bp},clip,scale=0.840]{A_WiSe20x21.pdf}};
\fill[Marker,opacity=0.30,rounded corners=0.6bp] (61.82bp,-12.18bp) rectangle (78.96bp,-22.26bp);
\fill[MarkerG,opacity=0.30,rounded corners=0.6bp] (85.95bp,-1.26bp) rectangle (116.76bp,-11.34bp);
\fill[MarkerG,opacity=0.30,rounded corners=0.6bp] (115.08bp,-1.26bp) rectangle (119.28bp,-11.34bp);
\end{tikzpicture}
\par\nopagebreak
\vspace{0.8mm}\nopagebreak
\begin{loesung}{A4Farbe}
\vspace{1.2mm}
\includegraphics[page=4,trim={52.00bp 693.50bp 39.00bp 94.50bp},clip,scale=0.700]{L_WiSe20x21.pdf}
\par
\par\vspace{1.0mm}
\includegraphics[page=4,trim={52.00bp 617.50bp 39.00bp 159.50bp},clip,scale=0.700]{L_WiSe20x21.pdf}
\par
\end{loesung}
\phantomsection\label{tf:WiSe20x21:4:2}
\par\vspace{1.6mm}
\Needspace*{55mm}
\begin{tikzpicture}[inner sep=0]
\node[anchor=north west,inner sep=0] (bi) at (0,0) {\includegraphics[page=5,trim={52.00bp 287.50bp 39.00bp 502.50bp},clip,scale=0.840]{A_WiSe20x21.pdf}};
\fill[Marker,opacity=0.30,rounded corners=0.6bp] (66.17bp,-1.26bp) rectangle (103.90bp,-11.34bp);
\fill[Marker,opacity=0.30,rounded corners=0.6bp] (113.40bp,-1.26bp) rectangle (157.75bp,-11.34bp);
\fill[Marker,opacity=0.30,rounded corners=0.6bp] (196.22bp,-1.26bp) rectangle (213.36bp,-11.34bp);
\fill[MarkerG,opacity=0.30,rounded corners=0.6bp] (54.02bp,-22.26bp) rectangle (94.92bp,-32.34bp);
\fill[MarkerG,opacity=0.30,rounded corners=0.6bp] (92.40bp,-22.26bp) rectangle (99.12bp,-32.34bp);
\end{tikzpicture}
\par\nopagebreak
\vspace{0.8mm}\nopagebreak
\begin{loesung}{A4Farbe}
\vspace{1.2mm}
\includegraphics[page=4,trim={52.00bp 502.50bp 39.00bp 286.50bp},clip,scale=0.700]{L_WiSe20x21.pdf}
\par
\par\vspace{1.0mm}
\includegraphics[page=4,trim={52.00bp 463.50bp 39.00bp 351.50bp},clip,scale=0.700]{L_WiSe20x21.pdf}
\par
\par\vspace{1.0mm}
\includegraphics[page=4,trim={52.00bp 439.50bp 39.00bp 390.50bp},clip,scale=0.700]{L_WiSe20x21.pdf}
\par
\par\vspace{1.0mm}
\includegraphics[page=4,trim={52.00bp 395.50bp 39.00bp 416.50bp},clip,scale=0.700]{L_WiSe20x21.pdf}
\par
\end{loesung}
\par\vspace{6mm}
\Needspace*{174.7mm}
\marke{\color{A4Farbe}A4 \textendash{} SoSe 2021}
\phantomsection\label{auf:SoSe21:4}
\aufgabenkopf{SoSe 2021 \textendash{} Aufgabe 4}{14 Punkte}{A4Farbe}
\begin{tikzpicture}[inner sep=0]
\node[anchor=north west,inner sep=0] (bi) at (0,0) {\includegraphics[page=5,trim={52.00bp 695.50bp 39.00bp 107.50bp},clip,scale=0.840]{A_SoSe21.pdf}};
\fill[Marker,opacity=0.30,rounded corners=0.6bp] (111.72bp,-12.18bp) rectangle (160.09bp,-22.26bp);
\fill[Marker,opacity=0.30,rounded corners=0.6bp] (156.07bp,-22.26bp) rectangle (181.94bp,-32.34bp);
\end{tikzpicture}
\par\vspace{1.2mm}
\includegraphics[page=5,trim={52.00bp 510.50bp 39.00bp 192.50bp},clip,scale=0.620]{A_SoSe21.pdf}
\par\vspace{1.2mm}
\phantomsection\label{tf:SoSe21:4:1}
\par\vspace{1.6mm}
\Needspace*{56mm}
\begin{tikzpicture}[inner sep=0]
\node[anchor=north west,inner sep=0] (bi) at (0,0) {\includegraphics[page=5,trim={52.00bp 417.50bp 39.00bp 411.50bp},clip,scale=0.840]{A_SoSe21.pdf}};
\fill[MarkerG,opacity=0.30,rounded corners=0.6bp] (94.26bp,-1.26bp) rectangle (135.24bp,-11.34bp);
\fill[MarkerG,opacity=0.30,rounded corners=0.6bp] (133.56bp,-1.26bp) rectangle (141.12bp,-11.34bp);
\end{tikzpicture}
\par\nopagebreak
\vspace{0.8mm}\nopagebreak
\begin{loesung}{A4Farbe}
\vspace{1.2mm}
\includegraphics[page=4,trim={52.00bp 644.50bp 39.00bp 132.50bp},clip,scale=0.700]{L_SoSe21.pdf}
\par
\par\vspace{1.0mm}
\includegraphics[page=4,trim={52.00bp 570.50bp 39.00bp 207.50bp},clip,scale=0.700]{L_SoSe21.pdf}
\par
\par\vspace{1.0mm}
\includegraphics[page=4,trim={52.00bp 544.50bp 39.00bp 284.50bp},clip,scale=0.700]{L_SoSe21.pdf}
\par
\par\vspace{1.0mm}
\includegraphics[page=4,trim={52.00bp 501.50bp 39.00bp 307.50bp},clip,scale=0.700]{L_SoSe21.pdf}
\par
\end{loesung}
\phantomsection\label{tf:SoSe21:4:2}
\par\vspace{1.6mm}
\Needspace*{49mm}
\begin{tikzpicture}[inner sep=0]
\node[anchor=north west,inner sep=0] (bi) at (0,0) {\includegraphics[page=5,trim={52.00bp 314.50bp 39.00bp 488.50bp},clip,scale=0.840]{A_SoSe21.pdf}};
\fill[Marker,opacity=0.30,rounded corners=0.6bp] (180.26bp,-0.42bp) rectangle (197.40bp,-10.50bp);
\fill[Marker,opacity=0.30,rounded corners=0.6bp] (248.64bp,-22.26bp) rectangle (278.04bp,-32.34bp);
\fill[MarkerG,opacity=0.30,rounded corners=0.6bp] (373.27bp,-0.42bp) rectangle (397.40bp,-10.50bp);
\end{tikzpicture}
\par\nopagebreak
\vspace{0.8mm}\nopagebreak
\begin{loesung}{A4Farbe}
\vspace{1.2mm}
\includegraphics[page=4,trim={52.00bp 411.50bp 39.00bp 417.50bp},clip,scale=0.700]{L_SoSe21.pdf}
\par
\par\vspace{1.0mm}
\includegraphics[page=4,trim={52.00bp 368.50bp 39.00bp 440.50bp},clip,scale=0.700]{L_SoSe21.pdf}
\par
\par\vspace{1.0mm}
\includegraphics[page=4,trim={52.00bp 287.50bp 39.00bp 483.50bp},clip,scale=0.700]{L_SoSe21.pdf}
\par
\end{loesung}
\par\vspace{6mm}
\Needspace*{218.9mm}
\marke{\color{A4Farbe}A4 \textendash{} WiSe 2021/22}
\phantomsection\label{auf:WiSe21x22:4}
\aufgabenkopf{WiSe 2021/22 \textendash{} Aufgabe 4}{17 Punkte}{A4Farbe}
\begin{tikzpicture}[inner sep=0]
\node[anchor=north west,inner sep=0] (bi) at (0,0) {\includegraphics[page=5,trim={52.00bp 682.50bp 39.00bp 107.50bp},clip,scale=0.840]{A_WiSe21x22.pdf}};
\fill[Marker,opacity=0.30,rounded corners=0.6bp] (362.54bp,-1.26bp) rectangle (379.68bp,-11.34bp);
\fill[Marker,opacity=0.30,rounded corners=0.6bp] (237.72bp,-12.18bp) rectangle (279.93bp,-22.26bp);
\fill[Marker,opacity=0.30,rounded corners=0.6bp] (89.54bp,-23.10bp) rectangle (106.68bp,-33.18bp);
\fill[Marker,opacity=0.30,rounded corners=0.6bp] (205.80bp,-23.10bp) rectangle (227.92bp,-33.18bp);
\fill[Marker,opacity=0.30,rounded corners=0.6bp] (210.00bp,-34.02bp) rectangle (246.43bp,-44.10bp);
\end{tikzpicture}
\par\vspace{1.2mm}
\includegraphics[page=5,trim={52.00bp 447.50bp 39.00bp 199.50bp},clip,scale=0.620]{A_WiSe21x22.pdf}
\par\vspace{1.2mm}
\phantomsection\label{tf:WiSe21x22:4:1}
\par\vspace{1.6mm}
\Needspace*{32mm}
\begin{tikzpicture}[inner sep=0]
\node[anchor=north west,inner sep=0] (bi) at (0,0) {\includegraphics[page=5,trim={52.00bp 353.50bp 39.00bp 462.50bp},clip,scale=0.840]{A_WiSe21x22.pdf}};
\fill[MarkerG,opacity=0.30,rounded corners=0.6bp] (68.08bp,-1.26bp) rectangle (95.76bp,-11.34bp);
\fill[MarkerG,opacity=0.30,rounded corners=0.6bp] (93.24bp,-1.26bp) rectangle (97.44bp,-11.34bp);
\end{tikzpicture}
\par\nopagebreak
\vspace{0.8mm}\nopagebreak
\begin{loesung}{A4Farbe}
\vspace{1.2mm}
\includegraphics[page=4,trim={52.00bp 682.50bp 39.00bp 126.50bp},clip,scale=0.700]{L_WiSe21x22.pdf}
\par
\par\vspace{1.0mm}
\includegraphics[page=4,trim={52.00bp 643.50bp 39.00bp 169.50bp},clip,scale=0.700]{L_WiSe21x22.pdf}
\par
\end{loesung}
\phantomsection\label{tf:WiSe21x22:4:2}
\par\vspace{1.6mm}
\Needspace*{102mm}
\begin{tikzpicture}[inner sep=0]
\node[anchor=north west,inner sep=0] (bi) at (0,0) {\includegraphics[page=5,trim={52.00bp 276.50bp 39.00bp 551.50bp},clip,scale=0.840]{A_WiSe21x22.pdf}};
\fill[Marker,opacity=0.30,rounded corners=0.6bp] (270.98bp,-1.26bp) rectangle (288.12bp,-11.34bp);
\fill[MarkerG,opacity=0.30,rounded corners=0.6bp] (102.35bp,-1.26bp) rectangle (127.68bp,-11.34bp);
\fill[MarkerG,opacity=0.30,rounded corners=0.6bp] (125.16bp,-1.26bp) rectangle (129.36bp,-11.34bp);
\end{tikzpicture}
\par\nopagebreak
\vspace{0.8mm}\nopagebreak
\begin{loesung}{A4Farbe}
\vspace{1.2mm}
\includegraphics[page=4,trim={52.00bp 523.50bp 39.00bp 266.50bp},clip,scale=0.700]{L_WiSe21x22.pdf}
\par
\par\vspace{1.0mm}
\includegraphics[page=4,trim={52.00bp 279.50bp 39.00bp 328.50bp},clip,scale=0.700]{L_WiSe21x22.pdf}
\par
\par\vspace{1.0mm}
\includegraphics[page=4,trim={52.00bp 238.50bp 39.00bp 572.50bp},clip,scale=0.700]{L_WiSe21x22.pdf}
\par
\par\vspace{1.0mm}
\includegraphics[page=4,trim={52.00bp 214.50bp 39.00bp 613.50bp},clip,scale=0.700]{L_WiSe21x22.pdf}
\par
\par\vspace{1.0mm}
\includegraphics[page=4,trim={52.00bp 176.50bp 39.00bp 637.50bp},clip,scale=0.700]{L_WiSe21x22.pdf}
\par
\end{loesung}
\par\vspace{6mm}
\Needspace*{222.9mm}
\marke{\color{A4Farbe}A4 \textendash{} SoSe 2022}
\phantomsection\label{auf:SoSe22:4}
\aufgabenkopf{SoSe 2022 \textendash{} Aufgabe 4}{23 Punkte}{A4Farbe}
\begin{tikzpicture}[inner sep=0]
\node[anchor=north west,inner sep=0] (bi) at (0,0) {\includegraphics[page=5,trim={52.00bp 708.50bp 39.00bp 106.50bp},clip,scale=0.840]{A_SoSe22.pdf}};
\fill[Marker,opacity=0.30,rounded corners=0.6bp] (220.08bp,-12.18bp) rectangle (260.68bp,-22.26bp);
\end{tikzpicture}
\par\vspace{1.2mm}
\includegraphics[page=5,trim={52.00bp 417.50bp 39.00bp 164.50bp},clip,scale=0.620]{A_SoSe22.pdf}
\par\vspace{1.2mm}
\phantomsection\label{tf:SoSe22:4:1}
\par\vspace{1.6mm}
\Needspace*{51mm}
\begin{tikzpicture}[inner sep=0]
\node[anchor=north west,inner sep=0] (bi) at (0,0) {\includegraphics[page=5,trim={52.00bp 367.50bp 39.00bp 461.50bp},clip,scale=0.840]{A_SoSe22.pdf}};
\fill[MarkerG,opacity=0.30,rounded corners=0.6bp] (58.18bp,-0.42bp) rectangle (77.19bp,-10.50bp);
\end{tikzpicture}
\par\nopagebreak
\vspace{0.8mm}\nopagebreak
\begin{loesung}{A4Farbe}
\vspace{1.2mm}
\includegraphics[page=4,trim={52.00bp 656.50bp 39.00bp 94.50bp},clip,scale=0.700]{L_SoSe22.pdf}
\par
\par\vspace{1.0mm}
\includegraphics[page=4,trim={52.00bp 583.50bp 39.00bp 195.50bp},clip,scale=0.700]{L_SoSe22.pdf}
\par
\end{loesung}
\phantomsection\label{tf:SoSe22:4:2}
\Needspace*{18mm}
\verweis{A4Farbe}{2.}{\gs{Kompensationsbauteil} ($C$ oder $L$, seriell oder parallel) so wählen, dass die Impedanz \emph{rein reell} wird}{gleiche Frage, nur andere Zahlenwerte \textendash{} siehe \textbf{SoSe 2021, Aufgabe 4.1}, \textbf{Seite \pageref{tf:SoSe21:4:1}}}
\phantomsection\label{tf:SoSe22:4:3}
\par\vspace{1.6mm}
\Needspace*{38mm}
\begin{tikzpicture}[inner sep=0]
\node[anchor=north west,inner sep=0] (bi) at (0,0) {\includegraphics[page=5,trim={52.00bp 213.50bp 39.00bp 589.50bp},clip,scale=0.840]{A_SoSe22.pdf}};
\fill[Marker,opacity=0.30,rounded corners=0.6bp] (361.70bp,-0.42bp) rectangle (378.84bp,-10.50bp);
\fill[Marker,opacity=0.30,rounded corners=0.6bp] (193.83bp,-11.34bp) rectangle (210.99bp,-21.42bp);
\fill[Marker,opacity=0.30,rounded corners=0.6bp] (237.72bp,-11.34bp) rectangle (253.96bp,-21.42bp);
\fill[Marker,opacity=0.30,rounded corners=0.6bp] (248.64bp,-23.10bp) rectangle (278.04bp,-33.18bp);
\fill[MarkerG,opacity=0.30,rounded corners=0.6bp] (74.67bp,-0.42bp) rectangle (99.96bp,-10.50bp);
\fill[MarkerG,opacity=0.30,rounded corners=0.6bp] (98.28bp,-0.42bp) rectangle (102.48bp,-10.50bp);
\end{tikzpicture}
\par\nopagebreak
\vspace{0.8mm}\nopagebreak
\begin{loesung}{A4Farbe}
\vspace{1.2mm}
\includegraphics[page=4,trim={52.00bp 298.50bp 39.00bp 531.50bp},clip,scale=0.700]{L_SoSe22.pdf}
\par
\par\vspace{1.0mm}
\includegraphics[page=4,trim={52.00bp 218.50bp 39.00bp 563.50bp},clip,scale=0.700]{L_SoSe22.pdf}
\par
\end{loesung}
\phantomsection\label{tf:SoSe22:4:4}
\par\vspace{1.6mm}
\Needspace*{29mm}
\begin{tikzpicture}[inner sep=0]
\node[anchor=north west,inner sep=0] (bi) at (0,0) {\includegraphics[page=5,trim={52.00bp 137.50bp 39.00bp 665.50bp},clip,scale=0.840]{A_SoSe22.pdf}};
\fill[MarkerG,opacity=0.30,rounded corners=0.6bp] (55.78bp,-12.18bp) rectangle (126.00bp,-22.26bp);
\end{tikzpicture}
\par\nopagebreak
\vspace{0.8mm}\nopagebreak
\begin{loesung}{A4Farbe}
\vspace{1.2mm}
\includegraphics[page=4,trim={52.00bp 143.50bp 39.00bp 665.50bp},clip,scale=0.700]{L_SoSe22.pdf}
\par
\end{loesung}
\par\vspace{6mm}
\Needspace*{185.1mm}
\marke{\color{A4Farbe}A4 \textendash{} WiSe 2022/23}
\phantomsection\label{auf:WiSe22x23:4}
\aufgabenkopf{WiSe 2022/23 \textendash{} Aufgabe 4}{19 Punkte}{A4Farbe}
\begin{tikzpicture}[inner sep=0]
\node[anchor=north west,inner sep=0] (bi) at (0,0) {\includegraphics[page=5,trim={52.00bp 707.50bp 39.00bp 106.50bp},clip,scale=0.840]{A_WiSe22x23.pdf}};
\fill[Marker,opacity=0.30,rounded corners=0.6bp] (209.16bp,-13.02bp) rectangle (248.54bp,-23.10bp);
\fill[Marker,opacity=0.30,rounded corners=0.6bp] (358.34bp,-13.02bp) rectangle (375.48bp,-23.10bp);
\end{tikzpicture}
\par\vspace{1.2mm}
\includegraphics[page=5,trim={52.00bp 501.50bp 39.00bp 190.50bp},clip,scale=0.620]{A_WiSe22x23.pdf}
\par\vspace{1.2mm}
\phantomsection\label{tf:WiSe22x23:4:1}
\Needspace*{18mm}
\verweis{A4Farbe}{1.}{\gs{Impedanz} über eine Leitung nach vorn transformieren (Leitungsgleichung)}{gleiche Frage, nur andere Zahlenwerte \textendash{} siehe \textbf{SoSe 2022, Aufgabe 4.1}, \textbf{Seite \pageref{tf:SoSe22:4:1}}}
\phantomsection\label{tf:WiSe22x23:4:2}
\par\vspace{1.6mm}
\Needspace*{70mm}
\begin{tikzpicture}[inner sep=0]
\node[anchor=north west,inner sep=0] (bi) at (0,0) {\includegraphics[page=5,trim={52.00bp 353.50bp 39.00bp 449.50bp},clip,scale=0.840]{A_WiSe22x23.pdf}};
\fill[MarkerG,opacity=0.30,rounded corners=0.6bp] (68.49bp,-0.42bp) rectangle (122.64bp,-10.50bp);
\fill[MarkerG,opacity=0.30,rounded corners=0.6bp] (120.12bp,-0.42bp) rectangle (124.32bp,-10.50bp);
\end{tikzpicture}
\par\nopagebreak
\vspace{0.8mm}\nopagebreak
\begin{loesung}{A4Farbe}
\vspace{1.2mm}
\includegraphics[page=4,trim={52.00bp 416.50bp 39.00bp 341.50bp},clip,scale=0.700]{L_WiSe22x23.pdf}
\par
\par\vspace{1.0mm}
\includegraphics[page=4,trim={52.00bp 289.50bp 39.00bp 435.50bp},clip,scale=0.700]{L_WiSe22x23.pdf}
\par
\end{loesung}
\phantomsection\label{tf:WiSe22x23:4:3}
\par\vspace{1.6mm}
\Needspace*{39mm}
\begin{tikzpicture}[inner sep=0]
\node[anchor=north west,inner sep=0] (bi) at (0,0) {\includegraphics[page=5,trim={52.00bp 275.50bp 39.00bp 526.50bp},clip,scale=0.840]{A_WiSe22x23.pdf}};
\fill[Marker,opacity=0.30,rounded corners=0.6bp] (79.80bp,-12.18bp) rectangle (109.20bp,-22.26bp);
\fill[Marker,opacity=0.30,rounded corners=0.6bp] (248.64bp,-23.10bp) rectangle (282.24bp,-33.18bp);
\fill[MarkerG,opacity=0.30,rounded corners=0.6bp] (53.20bp,-0.42bp) rectangle (114.24bp,-10.50bp);
\fill[MarkerG,opacity=0.30,rounded corners=0.6bp] (112.56bp,-0.42bp) rectangle (116.76bp,-10.50bp);
\end{tikzpicture}
\par\nopagebreak
\vspace{0.8mm}\nopagebreak
\begin{loesung}{A4Farbe}
\vspace{1.2mm}
\includegraphics[page=4,trim={52.00bp 241.50bp 39.00bp 588.50bp},clip,scale=0.700]{L_WiSe22x23.pdf}
\par
\par\vspace{1.0mm}
\includegraphics[page=4,trim={52.00bp 161.50bp 39.00bp 620.50bp},clip,scale=0.700]{L_WiSe22x23.pdf}
\par
\end{loesung}
\par\vspace{6mm}
\Needspace*{225.5mm}
\marke{\color{A4Farbe}A4 \textendash{} SoSe 2023}
\phantomsection\label{auf:SoSe23:4}
\aufgabenkopf{SoSe 2023 \textendash{} Aufgabe 4}{17 Punkte}{A4Farbe}
\begin{tikzpicture}[inner sep=0]
\node[anchor=north west,inner sep=0] (bi) at (0,0) {\includegraphics[page=5,trim={52.00bp 684.50bp 39.00bp 106.50bp},clip,scale=0.840]{A_SoSe23.pdf}};
\fill[Marker,opacity=0.30,rounded corners=0.6bp] (43.68bp,-12.18bp) rectangle (84.28bp,-22.26bp);
\fill[Marker,opacity=0.30,rounded corners=0.6bp] (80.30bp,-23.10bp) rectangle (97.44bp,-33.18bp);
\fill[Marker,opacity=0.30,rounded corners=0.6bp] (83.33bp,-34.02bp) rectangle (107.18bp,-44.10bp);
\end{tikzpicture}
\par\vspace{1.2mm}
\includegraphics[page=5,trim={52.00bp 415.50bp 39.00bp 213.50bp},clip,scale=0.620]{A_SoSe23.pdf}
\par\vspace{1.2mm}
\phantomsection\label{tf:SoSe23:4:1}
\par\vspace{1.6mm}
\Needspace*{85mm}
\begin{tikzpicture}[inner sep=0]
\node[anchor=north west,inner sep=0] (bi) at (0,0) {\includegraphics[page=5,trim={52.00bp 341.50bp 39.00bp 474.50bp},clip,scale=0.840]{A_SoSe23.pdf}};
\fill[MarkerG,opacity=0.30,rounded corners=0.6bp] (96.60bp,-1.26bp) rectangle (122.64bp,-11.34bp);
\fill[MarkerG,opacity=0.30,rounded corners=0.6bp] (120.96bp,-1.26bp) rectangle (125.16bp,-11.34bp);
\end{tikzpicture}
\par\nopagebreak
\vspace{0.8mm}\nopagebreak
\begin{loesung}{A4Farbe}
\vspace{1.2mm}
\includegraphics[page=4,trim={52.00bp 717.50bp 39.00bp 93.50bp},clip,scale=0.700]{L_SoSe23.pdf}
\par
\par\vspace{1.0mm}
\includegraphics[page=4,trim={52.00bp 463.50bp 39.00bp 134.50bp},clip,scale=0.700]{L_SoSe23.pdf}
\par
\end{loesung}
\par\vspace{1.6mm}
\begin{tikzpicture}[inner sep=0]
\node[anchor=north west,inner sep=0] (bi) at (0,0) {\includegraphics[page=5,trim={52.00bp 289.50bp 39.00bp 538.50bp},clip,scale=0.840]{A_SoSe23.pdf}};
\fill[Marker,opacity=0.30,rounded corners=0.6bp] (86.18bp,-1.26bp) rectangle (103.32bp,-11.34bp);
\end{tikzpicture}
\par
\phantomsection\label{tf:SoSe23:4:2}
\par\vspace{1.6mm}
\Needspace*{39mm}
\begin{tikzpicture}[inner sep=0]
\node[anchor=north west,inner sep=0] (bi) at (0,0) {\includegraphics[page=5,trim={52.00bp 251.50bp 39.00bp 576.50bp},clip,scale=0.840]{A_SoSe23.pdf}};
\fill[MarkerG,opacity=0.30,rounded corners=0.6bp] (65.86bp,-1.26bp) rectangle (113.40bp,-11.34bp);
\fill[MarkerG,opacity=0.30,rounded corners=0.6bp] (111.72bp,-1.26bp) rectangle (118.44bp,-11.34bp);
\end{tikzpicture}
\par\nopagebreak
\vspace{0.8mm}\nopagebreak
\begin{loesung}{A4Farbe}
\vspace{1.2mm}
\includegraphics[page=4,trim={52.00bp 301.50bp 39.00bp 436.50bp},clip,scale=0.700]{L_SoSe23.pdf}
\par
\end{loesung}
\phantomsection\label{tf:SoSe23:4:3}
\Needspace*{18mm}
\verweis{A4Farbe}{3.}{$\lambda/4$-Transformator entwerfen: \gs{Länge $l$} und \gs{Wellenwiderstand $Z_L$} für ein gefordertes $Z$}{gleiche Frage, nur andere Zahlenwerte \textendash{} siehe \textbf{WiSe 2022/23, Aufgabe 4.3}, \textbf{Seite \pageref{tf:WiSe22x23:4:3}}}
\par\vspace{6mm}
\Needspace*{225.0mm}
\marke{\color{A4Farbe}A4 \textendash{} WiSe 2023/24}
\phantomsection\label{auf:WiSe23x24:4}
\aufgabenkopf{WiSe 2023/24 \textendash{} Aufgabe 4}{19 Punkte}{A4Farbe}
\begin{tikzpicture}[inner sep=0]
\node[anchor=north west,inner sep=0] (bi) at (0,0) {\includegraphics[page=5,trim={52.00bp 707.50bp 39.00bp 107.50bp},clip,scale=0.840]{A_WiSe23x24.pdf}};
\fill[Marker,opacity=0.30,rounded corners=0.6bp] (333.48bp,-1.26bp) rectangle (375.85bp,-11.34bp);
\end{tikzpicture}
\par\vspace{1.2mm}
\includegraphics[page=5,trim={52.00bp 492.50bp 39.00bp 177.50bp},clip,scale=0.620]{A_WiSe23x24.pdf}
\par\vspace{1.2mm}
\phantomsection\label{tf:WiSe23x24:4:1}
\Needspace*{18mm}
\verweis{A4Farbe}{1.}{$\lambda/4$-Transformator entwerfen: \gs{Länge $l$} und \gs{Wellenwiderstand $Z_L$} für ein gefordertes $Z$}{gleiche Frage, nur andere Zahlenwerte \textendash{} siehe \textbf{WiSe 2022/23, Aufgabe 4.3}, \textbf{Seite \pageref{tf:WiSe22x23:4:3}}}
\phantomsection\label{tf:WiSe23x24:4:2}
\par\vspace{1.6mm}
\Needspace*{42mm}
\begin{tikzpicture}[inner sep=0]
\node[anchor=north west,inner sep=0] (bi) at (0,0) {\includegraphics[page=5,trim={52.00bp 351.50bp 39.00bp 464.50bp},clip,scale=0.840]{A_WiSe23x24.pdf}};
\fill[Marker,opacity=0.30,rounded corners=0.6bp] (208.32bp,-0.42bp) rectangle (241.92bp,-10.50bp);
\fill[MarkerG,opacity=0.30,rounded corners=0.6bp] (53.66bp,-0.42bp) rectangle (94.08bp,-10.50bp);
\fill[MarkerG,opacity=0.30,rounded corners=0.6bp] (92.40bp,-0.42bp) rectangle (99.12bp,-10.50bp);
\end{tikzpicture}
\par\nopagebreak
\vspace{0.8mm}\nopagebreak
\begin{loesung}{A4Farbe}
\vspace{1.2mm}
\includegraphics[page=5,trim={52.00bp 512.50bp 39.00bp 228.50bp},clip,scale=0.700]{L_WiSe23x24.pdf}
\par
\end{loesung}
\phantomsection\label{tf:WiSe23x24:4:3}
\par\vspace{1.6mm}
\Needspace*{102mm}
\begin{tikzpicture}[inner sep=0]
\node[anchor=north west,inner sep=0] (bi) at (0,0) {\includegraphics[page=5,trim={52.00bp 287.50bp 39.00bp 527.50bp},clip,scale=0.840]{A_WiSe23x24.pdf}};
\fill[MarkerG,opacity=0.30,rounded corners=0.6bp] (53.66bp,-1.26bp) rectangle (94.08bp,-11.34bp);
\fill[MarkerG,opacity=0.30,rounded corners=0.6bp] (92.40bp,-1.26bp) rectangle (99.12bp,-11.34bp);
\end{tikzpicture}
\par\nopagebreak
\vspace{0.8mm}\nopagebreak
\begin{loesung}{A4Farbe}
\vspace{1.2mm}
\includegraphics[page=5,trim={52.00bp 357.50bp 39.00bp 378.50bp},clip,scale=0.700]{L_WiSe23x24.pdf}
\par
\par\vspace{1.0mm}
\includegraphics[page=5,trim={52.00bp 167.50bp 39.00bp 494.50bp},clip,scale=0.700]{L_WiSe23x24.pdf}
\par
\par\vspace{1.0mm}
\includegraphics[page=5,trim={52.00bp 75.50bp 39.00bp 706.50bp},clip,scale=0.700]{L_WiSe23x24.pdf}
\par
\end{loesung}
\par\vspace{6mm}
\Needspace*{232.8mm}
\marke{\color{A4Farbe}A4 \textendash{} SoSe 2025}
\phantomsection\label{auf:SoSe25:4}
\aufgabenkopf{SoSe 2025 \textendash{} Aufgabe 4}{22 Punkte}{A4Farbe}
\begin{tikzpicture}[inner sep=0]
\node[anchor=north west,inner sep=0] (bi) at (0,0) {\includegraphics[page=5,trim={52.00bp 683.50bp 39.00bp 119.50bp},clip,scale=0.840]{A_SoSe25.pdf}};
\fill[Marker,opacity=0.30,rounded corners=0.6bp] (5.88bp,-22.26bp) rectangle (32.08bp,-32.34bp);
\end{tikzpicture}
\par\vspace{1.2mm}
\includegraphics[page=5,trim={52.00bp 323.50bp 39.00bp 169.50bp},clip,scale=0.620]{A_SoSe25.pdf}
\par\vspace{1.2mm}
\phantomsection\label{tf:SoSe25:4:1}
\par\vspace{1.6mm}
\Needspace*{59mm}
\begin{tikzpicture}[inner sep=0]
\node[anchor=north west,inner sep=0] (bi) at (0,0) {\includegraphics[page=5,trim={52.00bp 263.50bp 39.00bp 564.50bp},clip,scale=0.840]{A_SoSe25.pdf}};
\fill[MarkerG,opacity=0.30,rounded corners=0.6bp] (56.33bp,-1.26bp) rectangle (74.41bp,-11.34bp);
\end{tikzpicture}
\par\nopagebreak
\vspace{0.8mm}\nopagebreak
\begin{loesung}{A4Farbe}
\vspace{1.2mm}
\includegraphics[page=4,trim={52.00bp 652.50bp 39.00bp 107.50bp},clip,scale=0.700]{L_SoSe25.pdf}
\par
\par\vspace{1.0mm}
\includegraphics[page=4,trim={52.00bp 537.50bp 39.00bp 199.50bp},clip,scale=0.700]{L_SoSe25.pdf}
\par
\end{loesung}
\par\vspace{1.6mm}
\includegraphics[page=5,trim={52.00bp 213.50bp 39.00bp 615.50bp},clip,scale=0.840]{A_SoSe25.pdf}
\par
\phantomsection\label{tf:SoSe25:4:2}
\par\vspace{1.6mm}
\Needspace*{45mm}
\begin{tikzpicture}[inner sep=0]
\node[anchor=north west,inner sep=0] (bi) at (0,0) {\includegraphics[page=5,trim={52.00bp 175.50bp 39.00bp 653.50bp},clip,scale=0.840]{A_SoSe25.pdf}};
\fill[MarkerG,opacity=0.30,rounded corners=0.6bp] (92.82bp,-1.26bp) rectangle (144.55bp,-11.34bp);
\end{tikzpicture}
\par\nopagebreak
\vspace{0.8mm}\nopagebreak
\begin{loesung}{A4Farbe}
\vspace{1.2mm}
\includegraphics[page=4,trim={52.00bp 381.50bp 39.00bp 346.50bp},clip,scale=0.700]{L_SoSe25.pdf}
\par
\par\vspace{1.0mm}
\includegraphics[page=4,trim={52.00bp 355.50bp 39.00bp 470.50bp},clip,scale=0.700]{L_SoSe25.pdf}
\par
\end{loesung}
\phantomsection\label{tf:SoSe25:4:3}
\Needspace*{18mm}
\verweis{A4Farbe}{3.}{\gs{Kompensationsbauteil} ($C$ oder $L$, seriell oder parallel) so wählen, dass die Impedanz \emph{rein reell} wird}{gleiche Frage, nur andere Zahlenwerte \textendash{} siehe \textbf{SoSe 2021, Aufgabe 4.1}, \textbf{Seite \pageref{tf:SoSe21:4:1}}}
\par\vspace{6mm}
\clearpage
\marke{\color{A5Farbe}A5 \textendash{} Antennen}
\phantomsection\label{typ:5}
\par\vspace*{-3mm}{\color{A5Farbe}\rule{\textwidth}{2.5pt}}\par\vspace{1.5mm}
{\sffamily\bfseries\huge\color{A5Farbe} A5 \textendash{} Antennen}\par\vspace{1mm}
{\sffamily\color{Grau} alle 8 Klausuren \textendash{} 23 verschiedene Teilaufgaben ausgeschrieben, 3 Wiederholungen als Verweis}\par\vspace{2mm}
{\color{A5Farbe}\rule{\textwidth}{0.8pt}}\par\vspace{4mm}
{\sffamily\bfseries Register: Wenn gefragt wird \ldots}\hfill{\footnotesize\color{Grau}\colorbox{MarkerG!45}{grün} = das ist gesucht \textbf{·} rechts steht, wo die Aufgabe ausgeschrieben ist}\par\vspace{1.5mm}
\renewcommand{\arraystretch}{1.12}\footnotesize
\begin{tabularx}{\textwidth}{@{}>{\RaggedRight}X>{\RaggedRight\scriptsize}p{27mm}@{}}
\hline\rule{0pt}{3.5mm}
\gs{Maximale Entfernung} aus der geforderten Mindest-Empfangsleistung (dBm/dBi) \newline \hspace*{3mm}\textcolor{Grau}{\itshape \textendash{} Grundfall} & \textbf{S.\,\pageref{tf:WiSe20x21:5:1}} \textcolor{Grau}{WiSe 20/21, 1.} \\
\gs{Maximale Entfernung} aus der geforderten Mindest-Empfangsleistung (dBm/dBi) \newline \hspace*{3mm}\textcolor{Grau}{\itshape \textendash{} hier \emph{bidirektional} -- beide Richtungen rechnen, die schwächere begrenzt} & \textbf{S.\,\pageref{tf:WiSe21x22:5:1}} \textcolor{Grau}{WiSe 21/22, 1.} \\
\gs{Maximale Entfernung} aus der geforderten Mindest-Empfangsleistung (dBm/dBi) \newline \hspace*{3mm}\textcolor{Grau}{\itshape \textendash{} hier ist statt $P_{E,min}$ direkt $S_{E,min}$ gegeben} & \textbf{S.\,\pageref{tf:WiSe23x24:5:3}} \textcolor{Grau}{WiSe 23/24, 3.} \\
Störsender: bis zu welchem \gs{Abstand} liegt das Signal noch $x$\,dB über der Störung? & \textbf{S.\,\pageref{tf:WiSe20x21:5:2}} \textcolor{Grau}{WiSe 20/21, 2.} \\
Um welchen \gs{Faktor} ändert sich die Reichweite bei anderem Antennengewinn? \newline \hspace*{3mm}\textcolor{Grau}{\itshape \textendash{} getrennt für den ungestörten und den gestörten Fall} & \textbf{S.\,\pageref{tf:WiSe20x21:5:3}} \textcolor{Grau}{WiSe 20/21, 3.} \\
Um welchen \gs{Faktor} ändert sich die Reichweite bei anderem Antennengewinn? \newline \hspace*{3mm}\textcolor{Grau}{\itshape \textendash{} hier werden \emph{beide} Antennen getauscht -- dann fällt die Wurzel weg} & \textbf{S.\,\pageref{tf:WiSe23x24:5:4}} \textcolor{Grau}{WiSe 23/24, 4.} \\
Welche \gs{Sendeleistung} wäre für einen isotropen Kugelstrahler nötig? & \textbf{S.\,\pageref{tf:SoSe21:5:1}} \textcolor{Grau}{SoSe 21, 1.} \\
\gs{Anzahl der Dipole} bzw. Dipol-Linien aus dem geforderten Gewinn & \textbf{S.\,\pageref{tf:SoSe21:5:2}} \textcolor{Grau}{SoSe 21, 2.} \\
Gruppenantenne verkippt: welche \gs{Leistung} wird jetzt noch empfangen? & \textbf{S.\,\pageref{tf:WiSe21x22:5:2}} \textcolor{Grau}{WiSe 21/22, 2.} \\
\gs{Leistungsdichte} $S = G\,P/(4\pi r^2)$ am Empfangsort & \textbf{S.\,\pageref{tf:SoSe22:5:1}} \textcolor{Grau}{SoSe 22, 1.} \\
Welchen \gs{Antennengewinn} muss die \emph{Empfangs}antenne mindestens haben? & \textbf{S.\,\pageref{tf:SoSe22:5:2}} \textcolor{Grau}{SoSe 22, 2.} \\
\gs{Kippwinkel}: um wieviel darf die Empfangsantenne verdreht werden? \newline \hspace*{3mm}\textcolor{Grau}{\itshape \textendash{} Grundfall} & \textbf{S.\,\pageref{tf:SoSe22:5:3}} \textcolor{Grau}{SoSe 22, 3.} \\
\gs{Kippwinkel}: um wieviel darf die Empfangsantenne verdreht werden? \newline \hspace*{3mm}\textcolor{Grau}{\itshape \textendash{} SoSe 22 gibt einen Mindest\emph{gewinn} vor, hier kommt die Winkelgrenze aus $P_{E,min}$ bzw. $S_{E,min}$} & \textbf{S.\,\pageref{tf:WiSe23x24:5:2}} \textcolor{Grau}{WiSe 23/24, 2.} \\
\gs{Kippwinkel}: um wieviel darf die Empfangsantenne verdreht werden? \newline \hspace*{3mm}\textcolor{Grau}{\itshape \textendash{} SoSe 22 gibt einen Mindest\emph{gewinn} vor, hier kommt die Winkelgrenze aus $P_M$ und $S_M$} & \textbf{S.\,\pageref{tf:SoSe25:5:2}} \textcolor{Grau}{SoSe 25, 2.} \\
\gs{Gewinn} einer Dipolzeile / Dipol-Linie in dBi & \textbf{S.\,\pageref{tf:WiSe22x23:5:1}} \textcolor{Grau}{WiSe 22/23, 1.} \\
\gs{Leistungsdichte} einer Gruppenantenne unter einem Winkel (Gruppencharakteristik) \newline \hspace*{3mm}\textcolor{Grau}{\itshape \textendash{} Grundfall} & \textbf{S.\,\pageref{tf:WiSe22x23:5:2}} \textcolor{Grau}{WiSe 22/23, 2.} \\
\gs{Leistungsdichte} einer Gruppenantenne unter einem Winkel (Gruppencharakteristik) \newline \hspace*{3mm}\textcolor{Grau}{\itshape \textendash{} hier genau in Hauptstrahlrichtung} & \textbf{S.\,\pageref{tf:SoSe23:5:2}} \textcolor{Grau}{SoSe 23, 2.} \\
\gs{Leistungsdichte} einer Gruppenantenne unter einem Winkel (Gruppencharakteristik) \newline \hspace*{3mm}\textcolor{Grau}{\itshape \textendash{} hier für zwei Gruppen \emph{neben} der Hauptstrahlrichtung} & \textbf{S.\,\pageref{tf:SoSe23:5:3}} \textcolor{Grau}{SoSe 23, 3.} \\
\gs{Empfangsleistung} aus Leistungsdichte und Gewinn (über die Wirkfläche $A_W$) & \textbf{S.\,\pageref{tf:WiSe22x23:5:3}} \textcolor{Grau}{WiSe 22/23, 3.} \\
\gs{Verhältnis zweier Empfangsleistungen} in dB & \textbf{S.\,\pageref{tf:WiSe22x23:5:4}} \textcolor{Grau}{WiSe 22/23, 4.} \\
\gs{Phasenunterschied} für eine gewünschte Hauptstrahlrichtung & \textbf{S.\,\pageref{tf:SoSe23:5:1}} \textcolor{Grau}{SoSe 23, 1.} \\
\gs{Empfangsleistungen} mehrerer Antennengruppen an demselben Ort & \textbf{S.\,\pageref{tf:SoSe23:5:4}} \textcolor{Grau}{SoSe 23, 4.} \\
Frequenzwechsel: um wieviel \gs{dB} ändert sich die Empfangsleistung? & \textbf{S.\,\pageref{tf:SoSe25:5:4}} \textcolor{Grau}{SoSe 25, 4.} \\
\hline\end{tabularx}\normalsize
\par\vspace{5mm}
\begin{tcolorbox}[colback=A5Farbe!4, colframe=A5Farbe!55, boxrule=0.5pt, arc=1mm, left=3mm, right=2mm, top=2mm, bottom=2mm, breakable, title={\sffamily\bfseries\small Was diesmal neu drankommen könnte \normalfont\itshape\footnotesize -- eigene Prognose (keine offizielle Ankündigung): Umkehrungen, Kombinationen und Transfers zwischen den Aufgabentypen}, coltitle=white, colbacktitle=A5Farbe!70]
\sffamily\small
\begin{itemize}\setlength{\itemsep}{2.2mm}\setlength{\leftmargini}{4mm}
\item Kippwinkel beim \emph{Halbwellendipol} statt beim Elementardipol -- andere Richtcharakteristik $C(\vartheta)$!\\[0.6mm] {\footnotesize\color{A5Farbe}$\rightarrow$ $A_W=g\lambda^2/(4\pi)$ mit $g=1{,}64$ und $C(\vartheta)=\left|\frac{\cos(\frac{\pi}{2}\cos\vartheta)}{\sin\vartheta}\right|$, dann $P=S\,A_W\,C^2$. FS 8.5}
\item Umkehrung der Phasensteuerung: Schwenkwinkel der Hauptstrahlrichtung aus gegebenem Phasenunterschied $\Delta\varphi$\\[0.6mm] {\footnotesize\color{A5Farbe}$\rightarrow$ $u=\delta+\frac{2\pi}{\lambda_0}b\cos\vartheta$; Hauptkeule liegt bei $u=0$, also nach $\vartheta$ auflösen. FS 8.7}
\item Richtung der ersten Nullstelle einer Dipolzeile bzw.\ Nebenkeulendämpfung\\[0.6mm] {\footnotesize\color{A5Farbe}$\rightarrow$ Gruppen-Charakteristik null setzen: $\sin(Nu/2)=0$, d.\,h.\ $Nu/2=\pi$ -- daraus $\vartheta$ der ersten Nullstelle}
\item Mittenabstand $\lambda$ statt $\lambda/2$: was passiert mit der Gruppen-Charakteristik (zusätzliche Hauptkeulen)?\\[0.6mm] {\footnotesize\color{A5Farbe}$\rightarrow$ $u$ überstreicht mehr als $2\pi$ $\Rightarrow$ weitere Hauptkeulen; Grenze ist $b\le\lambda/2$. FS 8.7}
\item Zusätzliche Streckendämpfung in dB (Regen, Wand, Kabel zwischen Sender und Antenne) $\rightarrow$ Reichweite\\[0.6mm] {\footnotesize\color{A5Farbe}$\rightarrow$ Dämpfung in der dBm-Bilanz abziehen; für die Reichweite gilt $S\sim1/r^2$, also $r_{neu}=r_{alt}\cdot10^{-a_{dB}/20}$}
\item Wieviele Halbwellendipole braucht man für den Gewinn einer Parabolantenne? Und wie lang wird die Gruppe?\\[0.6mm] {\footnotesize\color{A5Farbe}$\rightarrow$ $g$ der Parabolantenne aus $A_W$; $N=g_{ges}/g_{Einzel}$ (bzw.\ $+3$\,dB je Verdopplung), Länge $=(N-1)\cdot b$}
\item Polarisationsverlust bzw.\ nicht optimal ausgerichtete Antennen kombiniert mit der dBm-Leistungsbilanz\\[0.6mm] {\footnotesize\color{A5Farbe}$\rightarrow$ Kippung als $C^2(\vartheta)$ ansetzen und als weiteren dB-Term in die Bilanz aufnehmen. K5}
\end{itemize}\end{tcolorbox}
\clearpage
\Needspace*{273.0mm}
\marke{\color{A5Farbe}A5 \textendash{} WiSe 2020/21}
\phantomsection\label{auf:WiSe20x21:5}
\aufgabenkopf{WiSe 2020/21 \textendash{} Aufgabe 5}{15 Punkte}{A5Farbe}
\begin{tikzpicture}[inner sep=0]
\node[anchor=north west,inner sep=0] (bi) at (0,0) {\includegraphics[page=6,trim={52.00bp 652.50bp 39.00bp 106.50bp},clip,scale=0.840]{A_WiSe20x21.pdf}};
\fill[Marker,opacity=0.30,rounded corners=0.6bp] (244.44bp,-1.26bp) rectangle (288.48bp,-11.34bp);
\end{tikzpicture}
\par\vspace{1.2mm}
\includegraphics[page=6,trim={52.00bp 458.50bp 39.00bp 211.50bp},clip,scale=0.620]{A_WiSe20x21.pdf}
\par\vspace{1.2mm}
\begin{tikzpicture}[inner sep=0]
\node[anchor=north west,inner sep=0] (bi) at (0,0) {\includegraphics[page=6,trim={52.00bp 367.50bp 39.00bp 416.50bp},clip,scale=0.840]{A_WiSe20x21.pdf}};
\fill[Marker,opacity=0.30,rounded corners=0.6bp] (123.48bp,-27.30bp) rectangle (169.49bp,-37.38bp);
\fill[Marker,opacity=0.30,rounded corners=0.6bp] (299.88bp,-27.30bp) rectangle (331.94bp,-37.38bp);
\fill[Marker,opacity=0.30,rounded corners=0.6bp] (141.12bp,-39.06bp) rectangle (171.72bp,-49.14bp);
\fill[Marker,opacity=0.30,rounded corners=0.6bp] (277.20bp,-39.06bp) rectangle (316.34bp,-49.14bp);
\end{tikzpicture}
\par\vspace{1.2mm}
\phantomsection\label{tf:WiSe20x21:5:1}
\par\vspace{1.6mm}
\Needspace*{64mm}
\begin{tikzpicture}[inner sep=0]
\node[anchor=north west,inner sep=0] (bi) at (0,0) {\includegraphics[page=6,trim={52.00bp 272.50bp 39.00bp 541.50bp},clip,scale=0.840]{A_WiSe20x21.pdf}};
\fill[MarkerG,opacity=0.30,rounded corners=0.6bp] (101.59bp,-2.94bp) rectangle (141.56bp,-13.02bp);
\end{tikzpicture}
\par\nopagebreak
\vspace{0.8mm}\nopagebreak
\begin{loesung}{A5Farbe}
\vspace{1.2mm}
\includegraphics[page=5,trim={52.00bp 626.50bp 39.00bp 94.50bp},clip,scale=0.700]{L_WiSe20x21.pdf}
\par
\par\vspace{1.0mm}
\includegraphics[page=5,trim={52.00bp 546.50bp 39.00bp 225.50bp},clip,scale=0.700]{L_WiSe20x21.pdf}
\par
\end{loesung}
\phantomsection\label{tf:WiSe20x21:5:2}
\par\vspace{1.6mm}
\Needspace*{59mm}
\begin{tikzpicture}[inner sep=0]
\node[anchor=north west,inner sep=0] (bi) at (0,0) {\includegraphics[page=6,trim={52.00bp 178.50bp 39.00bp 612.50bp},clip,scale=0.840]{A_WiSe20x21.pdf}};
\fill[Marker,opacity=0.30,rounded corners=0.6bp] (73.60bp,-1.26bp) rectangle (98.58bp,-11.34bp);
\fill[Marker,opacity=0.30,rounded corners=0.6bp] (336.84bp,-1.26bp) rectangle (373.57bp,-11.34bp);
\fill[Marker,opacity=0.30,rounded corners=0.6bp] (129.74bp,-11.34bp) rectangle (153.87bp,-21.42bp);
\fill[Marker,opacity=0.30,rounded corners=0.6bp] (184.75bp,-33.18bp) rectangle (201.97bp,-43.26bp);
\fill[MarkerG,opacity=0.30,rounded corners=0.6bp] (75.45bp,-22.26bp) rectangle (104.11bp,-32.34bp);
\end{tikzpicture}
\par\nopagebreak
\vspace{0.8mm}\nopagebreak
\begin{loesung}{A5Farbe}
\vspace{1.2mm}
\includegraphics[page=5,trim={52.00bp 364.50bp 39.00bp 334.50bp},clip,scale=0.700]{L_WiSe20x21.pdf}
\par
\end{loesung}
\phantomsection\label{tf:WiSe20x21:5:3}
\par\vspace{1.6mm}
\Needspace*{56mm}
\begin{tikzpicture}[inner sep=0]
\node[anchor=north west,inner sep=0] (bi) at (0,0) {\includegraphics[page=6,trim={52.00bp 89.50bp 39.00bp 688.50bp},clip,scale=0.840]{A_WiSe20x21.pdf}};
\fill[Marker,opacity=0.30,rounded corners=0.6bp] (97.72bp,-11.34bp) rectangle (128.52bp,-21.42bp);
\fill[MarkerG,opacity=0.30,rounded corners=0.6bp] (60.40bp,-22.26bp) rectangle (85.40bp,-32.34bp);
\end{tikzpicture}
\par\nopagebreak
\vspace{0.8mm}\nopagebreak
\begin{loesung}{A5Farbe}
\vspace{1.2mm}
\includegraphics[page=5,trim={52.00bp 265.50bp 39.00bp 499.50bp},clip,scale=0.700]{L_WiSe20x21.pdf}
\par
\par\vspace{1.0mm}
\includegraphics[page=5,trim={52.00bp 229.50bp 39.00bp 588.50bp},clip,scale=0.700]{L_WiSe20x21.pdf}
\par
\par\vspace{1.0mm}
\includegraphics[page=5,trim={52.00bp 202.50bp 39.00bp 625.50bp},clip,scale=0.700]{L_WiSe20x21.pdf}
\par
\end{loesung}
\par\vspace{6mm}
\Needspace*{188.6mm}
\marke{\color{A5Farbe}A5 \textendash{} SoSe 2021}
\phantomsection\label{auf:SoSe21:5}
\aufgabenkopf{SoSe 2021 \textendash{} Aufgabe 5}{13 Punkte}{A5Farbe}
\begin{tikzpicture}[inner sep=0]
\node[anchor=north west,inner sep=0] (bi) at (0,0) {\includegraphics[page=6,trim={52.00bp 626.50bp 39.00bp 106.50bp},clip,scale=0.840]{A_SoSe21.pdf}};
\fill[Marker,opacity=0.30,rounded corners=0.6bp] (192.83bp,-1.26bp) rectangle (221.31bp,-11.34bp);
\fill[Marker,opacity=0.30,rounded corners=0.6bp] (22.37bp,-38.22bp) rectangle (47.39bp,-48.30bp);
\end{tikzpicture}
\par\vspace{1.2mm}
\includegraphics[page=6,trim={52.00bp 506.50bp 39.00bp 257.50bp},clip,scale=0.840]{A_SoSe21.pdf}
\par\vspace{1.2mm}
\includegraphics[page=6,trim={52.00bp 471.50bp 39.00bp 358.50bp},clip,scale=0.840]{A_SoSe21.pdf}
\par\vspace{1.2mm}
\phantomsection\label{tf:SoSe21:5:1}
\par\vspace{1.6mm}
\Needspace*{62mm}
\begin{tikzpicture}[inner sep=0]
\node[anchor=north west,inner sep=0] (bi) at (0,0) {\includegraphics[page=6,trim={52.00bp 347.50bp 39.00bp 442.50bp},clip,scale=0.840]{A_SoSe21.pdf}};
\fill[Marker,opacity=0.30,rounded corners=0.6bp] (128.52bp,-1.26bp) rectangle (162.93bp,-11.34bp);
\fill[Marker,opacity=0.30,rounded corners=0.6bp] (66.36bp,-12.18bp) rectangle (94.53bp,-22.26bp);
\fill[Marker,opacity=0.30,rounded corners=0.6bp] (262.42bp,-12.18bp) rectangle (289.97bp,-22.26bp);
\fill[MarkerG,opacity=0.30,rounded corners=0.6bp] (52.18bp,-23.10bp) rectangle (90.72bp,-33.18bp);
\fill[MarkerG,opacity=0.30,rounded corners=0.6bp] (89.04bp,-23.10bp) rectangle (96.60bp,-33.18bp);
\end{tikzpicture}
\par\nopagebreak
\vspace{0.8mm}\nopagebreak
\begin{loesung}{A5Farbe}
\vspace{1.2mm}
\includegraphics[page=5,trim={52.00bp 583.50bp 39.00bp 107.50bp},clip,scale=0.700]{L_SoSe21.pdf}
\par
\end{loesung}
\phantomsection\label{tf:SoSe21:5:2}
\par\vspace{1.6mm}
\Needspace*{45mm}
\begin{tikzpicture}[inner sep=0]
\node[anchor=north west,inner sep=0] (bi) at (0,0) {\includegraphics[page=6,trim={52.00bp 259.50bp 39.00bp 544.50bp},clip,scale=0.840]{A_SoSe21.pdf}};
\fill[Marker,opacity=0.30,rounded corners=0.6bp] (162.43bp,-0.42bp) rectangle (186.13bp,-10.50bp);
\fill[Marker,opacity=0.30,rounded corners=0.6bp] (245.26bp,-0.42bp) rectangle (272.02bp,-10.50bp);
\fill[Marker,opacity=0.30,rounded corners=0.6bp] (209.16bp,-22.26bp) rectangle (246.68bp,-32.34bp);
\fill[MarkerG,opacity=0.30,rounded corners=0.6bp] (63.35bp,-0.42bp) rectangle (112.83bp,-10.50bp);
\end{tikzpicture}
\par\nopagebreak
\vspace{0.8mm}\nopagebreak
\begin{loesung}{A5Farbe}
\vspace{1.2mm}
\includegraphics[page=5,trim={52.00bp 505.50bp 39.00bp 322.50bp},clip,scale=0.700]{L_SoSe21.pdf}
\par
\par\vspace{1.0mm}
\includegraphics[page=5,trim={52.00bp 392.50bp 39.00bp 363.50bp},clip,scale=0.700]{L_SoSe21.pdf}
\par
\end{loesung}
\par\vspace{6mm}
\Needspace*{255.2mm}
\marke{\color{A5Farbe}A5 \textendash{} WiSe 2021/22}
\phantomsection\label{auf:WiSe21x22:5}
\aufgabenkopf{WiSe 2021/22 \textendash{} Aufgabe 5}{11 Punkte}{A5Farbe}
\begin{tikzpicture}[inner sep=0]
\node[anchor=north west,inner sep=0] (bi) at (0,0) {\includegraphics[page=6,trim={52.00bp 657.50bp 39.00bp 106.50bp},clip,scale=0.840]{A_WiSe21x22.pdf}};
\fill[Marker,opacity=0.30,rounded corners=0.6bp] (210.73bp,-1.26bp) rectangle (236.96bp,-11.34bp);
\fill[Marker,opacity=0.30,rounded corners=0.6bp] (374.38bp,-1.26bp) rectangle (392.42bp,-11.34bp);
\fill[Marker,opacity=0.30,rounded corners=0.6bp] (54.35bp,-12.18bp) rectangle (91.62bp,-22.26bp);
\fill[Marker,opacity=0.30,rounded corners=0.6bp] (373.80bp,-12.18bp) rectangle (417.06bp,-22.26bp);
\fill[Marker,opacity=0.30,rounded corners=0.6bp] (21.76bp,-22.26bp) rectangle (39.33bp,-32.34bp);
\fill[Marker,opacity=0.30,rounded corners=0.6bp] (48.72bp,-34.02bp) rectangle (75.81bp,-44.10bp);
\end{tikzpicture}
\par\vspace{1.2mm}
\includegraphics[page=6,trim={52.00bp 495.50bp 39.00bp 214.50bp},clip,scale=0.620]{A_WiSe21x22.pdf}
\par\vspace{1.2mm}
\phantomsection\label{tf:WiSe21x22:5:1}
\par\vspace{1.6mm}
\Needspace*{59mm}
\begin{tikzpicture}[inner sep=0]
\node[anchor=north west,inner sep=0] (bi) at (0,0) {\includegraphics[page=6,trim={52.00bp 428.50bp 39.00bp 387.50bp},clip,scale=0.840]{A_WiSe21x22.pdf}};
\fill[MarkerG,opacity=0.30,rounded corners=0.6bp] (123.85bp,-1.26bp) rectangle (169.68bp,-11.34bp);
\fill[MarkerG,opacity=0.30,rounded corners=0.6bp] (168.00bp,-1.26bp) rectangle (174.72bp,-11.34bp);
\end{tikzpicture}
\par\nopagebreak
\vspace{0.8mm}\nopagebreak
\begin{loesung}{A5Farbe}
\vspace{1.2mm}
\includegraphics[page=5,trim={52.00bp 562.50bp 39.00bp 107.50bp},clip,scale=0.700]{L_WiSe21x22.pdf}
\par
\end{loesung}
\phantomsection\label{tf:WiSe21x22:5:2}
\par\vspace{1.6mm}
\Needspace*{72mm}
\begin{tikzpicture}[inner sep=0]
\node[anchor=north west,inner sep=0] (bi) at (0,0) {\includegraphics[page=6,trim={52.00bp 338.50bp 39.00bp 463.50bp},clip,scale=0.840]{A_WiSe21x22.pdf}};
\fill[Marker,opacity=0.30,rounded corners=0.6bp] (225.12bp,-23.10bp) rectangle (254.20bp,-33.18bp);
\end{tikzpicture}
\par\nopagebreak
\vspace{0.8mm}\nopagebreak
\begin{loesung}{A5Farbe}
\vspace{1.2mm}
\includegraphics[page=5,trim={52.00bp 495.50bp 39.00bp 322.50bp},clip,scale=0.700]{L_WiSe21x22.pdf}
\par
\par\vspace{1.0mm}
\includegraphics[page=5,trim={52.00bp 302.50bp 39.00bp 356.50bp},clip,scale=0.700]{L_WiSe21x22.pdf}
\par
\end{loesung}
\par\vspace{1.6mm}
\includegraphics[page=6,trim={52.00bp 159.50bp 39.00bp 528.50bp},clip,scale=0.620]{A_WiSe21x22.pdf}
\par
\par\vspace{1.6mm}
\begin{tikzpicture}[inner sep=0]
\node[anchor=north west,inner sep=0] (bi) at (0,0) {\includegraphics[page=6,trim={52.00bp 91.50bp 39.00bp 711.50bp},clip,scale=0.840]{A_WiSe21x22.pdf}};
\fill[Marker,opacity=0.30,rounded corners=0.6bp] (335.16bp,-1.26bp) rectangle (361.15bp,-11.34bp);
\end{tikzpicture}
\par
\par\vspace{6mm}
\Needspace*{236.4mm}
\marke{\color{A5Farbe}A5 \textendash{} SoSe 2022}
\phantomsection\label{auf:SoSe22:5}
\aufgabenkopf{SoSe 2022 \textendash{} Aufgabe 5}{15 Punkte}{A5Farbe}
\includegraphics[page=6,trim={52.00bp 721.50bp 39.00bp 106.50bp},clip,scale=0.840]{A_SoSe22.pdf}
\par\vspace{1.2mm}
\includegraphics[page=6,trim={52.00bp 508.50bp 39.00bp 131.50bp},clip,scale=0.620]{A_SoSe22.pdf}
\par\vspace{1.2mm}
\phantomsection\label{tf:SoSe22:5:1}
\par\vspace{1.6mm}
\Needspace*{22mm}
\begin{tikzpicture}[inner sep=0]
\node[anchor=north west,inner sep=0] (bi) at (0,0) {\includegraphics[page=6,trim={52.00bp 482.50bp 39.00bp 346.50bp},clip,scale=0.840]{A_SoSe22.pdf}};
\fill[MarkerG,opacity=0.30,rounded corners=0.6bp] (51.39bp,-0.42bp) rectangle (116.76bp,-10.50bp);
\fill[MarkerG,opacity=0.30,rounded corners=0.6bp] (115.08bp,-0.42bp) rectangle (121.80bp,-10.50bp);
\end{tikzpicture}
\par\nopagebreak
\vspace{0.8mm}\nopagebreak
\begin{loesung}{A5Farbe}
\vspace{1.2mm}
\includegraphics[page=5,trim={52.00bp 696.50bp 39.00bp 107.50bp},clip,scale=0.700]{L_SoSe22.pdf}
\par
\end{loesung}
\phantomsection\label{tf:SoSe22:5:2}
\par\vspace{1.6mm}
\Needspace*{53mm}
\begin{tikzpicture}[inner sep=0]
\node[anchor=north west,inner sep=0] (bi) at (0,0) {\includegraphics[page=6,trim={52.00bp 407.50bp 39.00bp 408.50bp},clip,scale=0.840]{A_SoSe22.pdf}};
\fill[Marker,opacity=0.30,rounded corners=0.6bp] (213.30bp,-11.34bp) rectangle (252.28bp,-21.42bp);
\fill[MarkerG,opacity=0.30,rounded corners=0.6bp] (56.75bp,-1.26bp) rectangle (116.57bp,-11.34bp);
\end{tikzpicture}
\par\nopagebreak
\vspace{0.8mm}\nopagebreak
\begin{loesung}{A5Farbe}
\vspace{1.2mm}
\includegraphics[page=5,trim={52.00bp 485.50bp 39.00bp 208.50bp},clip,scale=0.700]{L_SoSe22.pdf}
\par
\end{loesung}
\phantomsection\label{tf:SoSe22:5:3}
\par\vspace{1.6mm}
\Needspace*{49mm}
\begin{tikzpicture}[inner sep=0]
\node[anchor=north west,inner sep=0] (bi) at (0,0) {\includegraphics[page=6,trim={52.00bp 330.50bp 39.00bp 471.50bp},clip,scale=0.840]{A_SoSe22.pdf}};
\fill[Marker,opacity=0.30,rounded corners=0.6bp] (64.07bp,-23.10bp) rectangle (80.50bp,-33.18bp);
\fill[MarkerG,opacity=0.30,rounded corners=0.6bp] (68.61bp,-12.18bp) rectangle (123.58bp,-22.26bp);
\end{tikzpicture}
\par\nopagebreak
\vspace{0.8mm}\nopagebreak
\begin{loesung}{A5Farbe}
\vspace{1.2mm}
\includegraphics[page=5,trim={52.00bp 315.50bp 39.00bp 411.50bp},clip,scale=0.700]{L_SoSe22.pdf}
\par
\end{loesung}
\par\vspace{1.6mm}
\includegraphics[page=6,trim={52.00bp 143.50bp 39.00bp 538.50bp},clip,scale=0.620]{A_SoSe22.pdf}
\par
\par\vspace{6mm}
\Needspace*{273.0mm}
\marke{\color{A5Farbe}A5 \textendash{} WiSe 2022/23}
\phantomsection\label{auf:WiSe22x23:5}
\aufgabenkopf{WiSe 2022/23 \textendash{} Aufgabe 5}{18 Punkte}{A5Farbe}
\begin{tikzpicture}[inner sep=0]
\node[anchor=north west,inner sep=0] (bi) at (0,0) {\includegraphics[page=6,trim={52.00bp 650.50bp 39.00bp 106.50bp},clip,scale=0.840]{A_WiSe22x23.pdf}};
\fill[Marker,opacity=0.30,rounded corners=0.6bp] (241.08bp,-12.18bp) rectangle (276.46bp,-22.26bp);
\fill[Marker,opacity=0.30,rounded corners=0.6bp] (57.96bp,-23.10bp) rectangle (93.14bp,-33.18bp);
\fill[Marker,opacity=0.30,rounded corners=0.6bp] (14.76bp,-34.02bp) rectangle (35.67bp,-44.10bp);
\fill[Marker,opacity=0.30,rounded corners=0.6bp] (120.48bp,-34.02bp) rectangle (151.92bp,-44.10bp);
\fill[Marker,opacity=0.30,rounded corners=0.6bp] (67.20bp,-49.98bp) rectangle (95.89bp,-60.06bp);
\fill[Marker,opacity=0.30,rounded corners=0.6bp] (120.96bp,-49.98bp) rectangle (152.96bp,-60.06bp);
\fill[Marker,opacity=0.30,rounded corners=0.6bp] (340.20bp,-49.98bp) rectangle (367.92bp,-60.06bp);
\fill[Marker,opacity=0.30,rounded corners=0.6bp] (195.72bp,-60.90bp) rectangle (233.71bp,-70.98bp);
\fill[Marker,opacity=0.30,rounded corners=0.6bp] (360.59bp,-60.90bp) rectangle (390.98bp,-70.98bp);
\end{tikzpicture}
\par\vspace{1.2mm}
\includegraphics[page=6,trim={52.00bp 546.50bp 39.00bp 213.50bp},clip,scale=0.620]{A_WiSe22x23.pdf}
\par\vspace{1.2mm}
\includegraphics[page=6,trim={52.00bp 312.50bp 39.00bp 347.50bp},clip,scale=0.620]{A_WiSe22x23.pdf}
\par\vspace{1.2mm}
\phantomsection\label{tf:WiSe22x23:5:1}
\par\vspace{1.6mm}
\Needspace*{21mm}
\begin{tikzpicture}[inner sep=0]
\node[anchor=north west,inner sep=0] (bi) at (0,0) {\includegraphics[page=6,trim={52.00bp 271.50bp 39.00bp 557.50bp},clip,scale=0.840]{A_WiSe22x23.pdf}};
\fill[Marker,opacity=0.30,rounded corners=0.6bp] (234.36bp,-1.26bp) rectangle (262.08bp,-11.34bp);
\fill[MarkerG,opacity=0.30,rounded corners=0.6bp] (53.76bp,-1.26bp) rectangle (78.12bp,-11.34bp);
\end{tikzpicture}
\par\nopagebreak
\vspace{0.8mm}\nopagebreak
\begin{loesung}{A5Farbe}
\vspace{1.2mm}
\includegraphics[page=5,trim={52.00bp 720.50bp 39.00bp 106.50bp},clip,scale=0.700]{L_WiSe22x23.pdf}
\par
\par\vspace{1.0mm}
\includegraphics[page=5,trim={52.00bp 687.50bp 39.00bp 136.50bp},clip,scale=0.700]{L_WiSe22x23.pdf}
\par
\end{loesung}
\phantomsection\label{tf:WiSe22x23:5:2}
\par\vspace{1.6mm}
\Needspace*{62mm}
\begin{tikzpicture}[inner sep=0]
\node[anchor=north west,inner sep=0] (bi) at (0,0) {\includegraphics[page=6,trim={52.00bp 221.50bp 39.00bp 607.50bp},clip,scale=0.840]{A_WiSe22x23.pdf}};
\fill[MarkerG,opacity=0.30,rounded corners=0.6bp] (51.39bp,-1.26bp) rectangle (116.76bp,-11.34bp);
\fill[MarkerG,opacity=0.30,rounded corners=0.6bp] (115.08bp,-1.26bp) rectangle (121.80bp,-11.34bp);
\end{tikzpicture}
\par\nopagebreak
\vspace{0.8mm}\nopagebreak
\begin{loesung}{A5Farbe}
\vspace{1.2mm}
\includegraphics[page=5,trim={52.00bp 432.50bp 39.00bp 208.50bp},clip,scale=0.700]{L_WiSe22x23.pdf}
\par
\end{loesung}
\par\vspace{1.6mm}
\begin{tikzpicture}[inner sep=0]
\node[anchor=north west,inner sep=0] (bi) at (0,0) {\includegraphics[page=6,trim={52.00bp 164.50bp 39.00bp 651.50bp},clip,scale=0.840]{A_WiSe22x23.pdf}};
\fill[Marker,opacity=0.30,rounded corners=0.6bp] (294.84bp,-1.26bp) rectangle (330.27bp,-11.34bp);
\fill[Marker,opacity=0.30,rounded corners=0.6bp] (247.67bp,-11.34bp) rectangle (277.85bp,-21.42bp);
\end{tikzpicture}
\par
\phantomsection\label{tf:WiSe22x23:5:3}
\par\vspace{1.6mm}
\Needspace*{22mm}
\begin{tikzpicture}[inner sep=0]
\node[anchor=north west,inner sep=0] (bi) at (0,0) {\includegraphics[page=6,trim={52.00bp 126.50bp 39.00bp 701.50bp},clip,scale=0.840]{A_WiSe22x23.pdf}};
\fill[MarkerG,opacity=0.30,rounded corners=0.6bp] (51.52bp,-1.26bp) rectangle (85.20bp,-11.34bp);
\end{tikzpicture}
\par\nopagebreak
\vspace{0.8mm}\nopagebreak
\begin{loesung}{A5Farbe}
\vspace{1.2mm}
\includegraphics[page=5,trim={52.00bp 358.50bp 39.00bp 449.50bp},clip,scale=0.700]{L_WiSe22x23.pdf}
\par
\end{loesung}
\phantomsection\label{tf:WiSe22x23:5:4}
\par\vspace{1.6mm}
\Needspace*{54mm}
\begin{tikzpicture}[inner sep=0]
\node[anchor=north west,inner sep=0] (bi) at (0,0) {\includegraphics[page=6,trim={52.00bp 82.50bp 39.00bp 733.50bp},clip,scale=0.840]{A_WiSe22x23.pdf}};
\fill[MarkerG,opacity=0.30,rounded corners=0.6bp] (69.66bp,-0.42bp) rectangle (109.80bp,-10.50bp);
\end{tikzpicture}
\par\nopagebreak
\vspace{0.8mm}\nopagebreak
\begin{loesung}{A5Farbe}
\vspace{1.2mm}
\includegraphics[page=5,trim={52.00bp 167.50bp 39.00bp 525.50bp},clip,scale=0.700]{L_WiSe22x23.pdf}
\par
\end{loesung}
\par\vspace{6mm}
\Needspace*{273.0mm}
\marke{\color{A5Farbe}A5 \textendash{} SoSe 2023}
\phantomsection\label{auf:SoSe23:5}
\aufgabenkopf{SoSe 2023 \textendash{} Aufgabe 5}{24 Punkte}{A5Farbe}
\begin{tikzpicture}[inner sep=0]
\node[anchor=north west,inner sep=0] (bi) at (0,0) {\includegraphics[page=6,trim={52.00bp 650.50bp 39.00bp 106.50bp},clip,scale=0.840]{A_SoSe23.pdf}};
\fill[Marker,opacity=0.30,rounded corners=0.6bp] (220.92bp,-44.94bp) rectangle (246.33bp,-55.02bp);
\fill[Marker,opacity=0.30,rounded corners=0.6bp] (227.64bp,-60.90bp) rectangle (264.91bp,-70.98bp);
\fill[Marker,opacity=0.30,rounded corners=0.6bp] (337.68bp,-60.90bp) rectangle (377.08bp,-70.98bp);
\end{tikzpicture}
\par\vspace{1.2mm}
\includegraphics[page=6,trim={52.00bp 578.50bp 39.00bp 209.50bp},clip,scale=0.840]{A_SoSe23.pdf}
\par\vspace{1.2mm}
\begin{tikzpicture}[inner sep=0]
\node[anchor=north west,inner sep=0] (bi) at (0,0) {\includegraphics[page=6,trim={52.00bp 358.50bp 39.00bp 287.50bp},clip,scale=0.840]{A_SoSe23.pdf}};
\fill[Marker,opacity=0.30,rounded corners=0.6bp] (133.56bp,-40.74bp) rectangle (164.64bp,-54.18bp);
\fill[Marker,opacity=0.30,rounded corners=0.6bp] (30.24bp,-46.62bp) rectangle (48.30bp,-60.06bp);
\fill[Marker,opacity=0.30,rounded corners=0.6bp] (99.12bp,-100.38bp) rectangle (130.20bp,-113.82bp);
\end{tikzpicture}
\par\vspace{1.2mm}
\includegraphics[page=6,trim={52.00bp 315.50bp 39.00bp 501.50bp},clip,scale=0.840]{A_SoSe23.pdf}
\par\vspace{1.2mm}
\phantomsection\label{tf:SoSe23:5:1}
\par\vspace{1.6mm}
\Needspace*{26mm}
\begin{tikzpicture}[inner sep=0]
\node[anchor=north west,inner sep=0] (bi) at (0,0) {\includegraphics[page=6,trim={52.00bp 263.50bp 39.00bp 552.50bp},clip,scale=0.840]{A_SoSe23.pdf}};
\fill[MarkerG,opacity=0.30,rounded corners=0.6bp] (55.83bp,-1.26bp) rectangle (126.16bp,-11.34bp);
\end{tikzpicture}
\par\nopagebreak
\vspace{0.8mm}\nopagebreak
\begin{loesung}{A5Farbe}
\vspace{1.2mm}
\includegraphics[page=5,trim={52.00bp 717.50bp 39.00bp 88.50bp},clip,scale=0.700]{L_SoSe23.pdf}
\par
\end{loesung}
\par\vspace{1.6mm}
\begin{tikzpicture}[inner sep=0]
\node[anchor=north west,inner sep=0] (bi) at (0,0) {\includegraphics[page=6,trim={52.00bp 213.50bp 39.00bp 603.50bp},clip,scale=0.840]{A_SoSe23.pdf}};
\fill[Marker,opacity=0.30,rounded corners=0.6bp] (340.94bp,-1.26bp) rectangle (352.38bp,-11.34bp);
\fill[Marker,opacity=0.30,rounded corners=0.6bp] (52.08bp,-11.34bp) rectangle (88.20bp,-21.42bp);
\end{tikzpicture}
\par
\phantomsection\label{tf:SoSe23:5:2}
\par\vspace{1.6mm}
\Needspace*{35mm}
\begin{tikzpicture}[inner sep=0]
\node[anchor=north west,inner sep=0] (bi) at (0,0) {\includegraphics[page=6,trim={52.00bp 180.50bp 39.00bp 647.50bp},clip,scale=0.840]{A_SoSe23.pdf}};
\fill[MarkerG,opacity=0.30,rounded corners=0.6bp] (50.86bp,-1.26bp) rectangle (111.13bp,-11.34bp);
\end{tikzpicture}
\par\nopagebreak
\vspace{0.8mm}\nopagebreak
\begin{loesung}{A5Farbe}
\vspace{1.2mm}
\includegraphics[page=5,trim={52.00bp 603.50bp 39.00bp 151.50bp},clip,scale=0.700]{L_SoSe23.pdf}
\par
\end{loesung}
\phantomsection\label{tf:SoSe23:5:3}
\par\vspace{1.6mm}
\Needspace*{67mm}
\begin{tikzpicture}[inner sep=0]
\node[anchor=north west,inner sep=0] (bi) at (0,0) {\includegraphics[page=6,trim={52.00bp 155.50bp 39.00bp 672.50bp},clip,scale=0.840]{A_SoSe23.pdf}};
\fill[MarkerG,opacity=0.30,rounded corners=0.6bp] (50.88bp,-1.26bp) rectangle (115.09bp,-11.34bp);
\end{tikzpicture}
\par\nopagebreak
\vspace{0.8mm}\nopagebreak
\begin{loesung}{A5Farbe}
\vspace{1.2mm}
\includegraphics[page=5,trim={52.00bp 358.50bp 39.00bp 265.50bp},clip,scale=0.700]{L_SoSe23.pdf}
\par
\end{loesung}
\phantomsection\label{tf:SoSe23:5:4}
\par\vspace{1.6mm}
\Needspace*{87mm}
\begin{tikzpicture}[inner sep=0]
\node[anchor=north west,inner sep=0] (bi) at (0,0) {\includegraphics[page=6,trim={52.00bp 91.50bp 39.00bp 711.50bp},clip,scale=0.840]{A_SoSe23.pdf}};
\fill[Marker,opacity=0.30,rounded corners=0.6bp] (101.64bp,-1.26bp) rectangle (127.45bp,-11.34bp);
\fill[Marker,opacity=0.30,rounded corners=0.6bp] (160.13bp,-1.26bp) rectangle (177.09bp,-11.34bp);
\fill[Marker,opacity=0.30,rounded corners=0.6bp] (320.88bp,-1.26bp) rectangle (357.84bp,-11.34bp);
\fill[Marker,opacity=0.30,rounded corners=0.6bp] (102.68bp,-11.34bp) rectangle (121.03bp,-21.42bp);
\fill[MarkerG,opacity=0.30,rounded corners=0.6bp] (50.79bp,-22.26bp) rectangle (122.59bp,-32.34bp);
\end{tikzpicture}
\par\nopagebreak
\vspace{0.8mm}\nopagebreak
\begin{loesung}{A5Farbe}
\vspace{1.2mm}
\includegraphics[page=5,trim={52.00bp 54.50bp 39.00bp 519.50bp},clip,scale=0.700]{L_SoSe23.pdf}
\par
\end{loesung}
\par\vspace{6mm}
\Needspace*{263.1mm}
\marke{\color{A5Farbe}A5 \textendash{} WiSe 2023/24}
\phantomsection\label{auf:WiSe23x24:5}
\aufgabenkopf{WiSe 2023/24 \textendash{} Aufgabe 5}{17 Punkte}{A5Farbe}
\begin{tikzpicture}[inner sep=0]
\node[anchor=north west,inner sep=0] (bi) at (0,0) {\includegraphics[page=6,trim={52.00bp 695.50bp 39.00bp 106.50bp},clip,scale=0.840]{A_WiSe23x24.pdf}};
\fill[Marker,opacity=0.30,rounded corners=0.6bp] (238.56bp,-1.26bp) rectangle (280.68bp,-11.34bp);
\fill[Marker,opacity=0.30,rounded corners=0.6bp] (366.87bp,-1.26bp) rectangle (404.46bp,-11.34bp);
\fill[Marker,opacity=0.30,rounded corners=0.6bp] (320.88bp,-12.18bp) rectangle (349.86bp,-22.26bp);
\end{tikzpicture}
\par\vspace{1.2mm}
\includegraphics[page=6,trim={52.00bp 582.50bp 39.00bp 177.50bp},clip,scale=0.620]{A_WiSe23x24.pdf}
\par\vspace{1.2mm}
\includegraphics[page=6,trim={52.00bp 548.50bp 39.00bp 267.50bp},clip,scale=0.840]{A_WiSe23x24.pdf}
\par\vspace{1.2mm}
\phantomsection\label{tf:WiSe23x24:5:1}
\Needspace*{18mm}
\verweis{A5Farbe}{1.}{\gs{Leistungsdichte} $S = G\,P/(4\pi r^2)$ am Empfangsort}{gleiche Frage, nur andere Zahlenwerte \textendash{} siehe \textbf{SoSe 2022, Aufgabe 5.1}, \textbf{Seite \pageref{tf:SoSe22:5:1}}}
\phantomsection\label{tf:WiSe23x24:5:2}
\par\vspace{1.6mm}
\Needspace*{66mm}
\begin{tikzpicture}[inner sep=0]
\node[anchor=north west,inner sep=0] (bi) at (0,0) {\includegraphics[page=6,trim={52.00bp 405.50bp 39.00bp 397.50bp},clip,scale=0.840]{A_WiSe23x24.pdf}};
\fill[Marker,opacity=0.30,rounded corners=0.6bp] (60.48bp,-12.18bp) rectangle (89.36bp,-22.26bp);
\fill[Marker,opacity=0.30,rounded corners=0.6bp] (208.32bp,-23.10bp) rectangle (252.00bp,-33.18bp);
\fill[MarkerG,opacity=0.30,rounded corners=0.6bp] (71.77bp,-1.26bp) rectangle (99.77bp,-11.34bp);
\end{tikzpicture}
\par\nopagebreak
\vspace{0.8mm}\nopagebreak
\begin{loesung}{A5Farbe}
\vspace{1.2mm}
\includegraphics[page=6,trim={52.00bp 487.50bp 39.00bp 170.50bp},clip,scale=0.700]{L_WiSe23x24.pdf}
\par
\end{loesung}
\par\vspace{1.6mm}
\includegraphics[page=6,trim={52.00bp 249.50bp 39.00bp 449.50bp},clip,scale=0.620]{A_WiSe23x24.pdf}
\par
\phantomsection\label{tf:WiSe23x24:5:3}
\par\vspace{1.6mm}
\Needspace*{35mm}
\begin{tikzpicture}[inner sep=0]
\node[anchor=north west,inner sep=0] (bi) at (0,0) {\includegraphics[page=6,trim={52.00bp 155.50bp 39.00bp 660.50bp},clip,scale=0.840]{A_WiSe23x24.pdf}};
\fill[Marker,opacity=0.30,rounded corners=0.6bp] (229.84bp,-11.34bp) rectangle (254.80bp,-21.42bp);
\fill[MarkerG,opacity=0.30,rounded corners=0.6bp] (90.72bp,-1.26bp) rectangle (138.60bp,-11.34bp);
\fill[MarkerG,opacity=0.30,rounded corners=0.6bp] (136.92bp,-1.26bp) rectangle (141.12bp,-11.34bp);
\end{tikzpicture}
\par\nopagebreak
\vspace{0.8mm}\nopagebreak
\begin{loesung}{A5Farbe}
\vspace{1.2mm}
\includegraphics[page=6,trim={52.00bp 357.50bp 39.00bp 410.50bp},clip,scale=0.700]{L_WiSe23x24.pdf}
\par
\end{loesung}
\phantomsection\label{tf:WiSe23x24:5:4}
\par\vspace{1.6mm}
\Needspace*{53mm}
\begin{tikzpicture}[inner sep=0]
\node[anchor=north west,inner sep=0] (bi) at (0,0) {\includegraphics[page=6,trim={52.00bp 92.50bp 39.00bp 722.50bp},clip,scale=0.840]{A_WiSe23x24.pdf}};
\fill[MarkerG,opacity=0.30,rounded corners=0.6bp] (67.12bp,-1.26bp) rectangle (92.58bp,-11.34bp);
\end{tikzpicture}
\par\nopagebreak
\vspace{0.8mm}\nopagebreak
\begin{loesung}{A5Farbe}
\vspace{1.2mm}
\includegraphics[page=6,trim={52.00bp 159.50bp 39.00bp 538.50bp},clip,scale=0.700]{L_WiSe23x24.pdf}
\par
\end{loesung}
\par\vspace{6mm}
\Needspace*{273.0mm}
\marke{\color{A5Farbe}A5 \textendash{} SoSe 2025}
\phantomsection\label{auf:SoSe25:5}
\aufgabenkopf{SoSe 2025 \textendash{} Aufgabe 5}{21 Punkte}{A5Farbe}
\begin{tikzpicture}[inner sep=0]
\node[anchor=north west,inner sep=0] (bi) at (0,0) {\includegraphics[page=6,trim={52.00bp 685.50bp 39.00bp 106.50bp},clip,scale=0.840]{A_SoSe25.pdf}};
\fill[Marker,opacity=0.30,rounded corners=0.6bp] (238.02bp,-1.26bp) rectangle (272.62bp,-11.34bp);
\fill[Marker,opacity=0.30,rounded corners=0.6bp] (16.43bp,-12.18bp) rectangle (48.78bp,-22.26bp);
\end{tikzpicture}
\par\vspace{1.2mm}
\includegraphics[page=6,trim={52.00bp 431.50bp 39.00bp 180.50bp},clip,scale=0.620]{A_SoSe25.pdf}
\par\vspace{1.2mm}
\phantomsection\label{tf:SoSe25:5:1}
\Needspace*{18mm}
\verweis{A5Farbe}{1.}{\gs{Leistungsdichte} $S = G\,P/(4\pi r^2)$ am Empfangsort}{gleiche Frage, nur andere Zahlenwerte \textendash{} siehe \textbf{SoSe 2022, Aufgabe 5.1}, \textbf{Seite \pageref{tf:SoSe22:5:1}}}
\phantomsection\label{tf:SoSe25:5:2}
\par\vspace{1.6mm}
\Needspace*{61mm}
\begin{tikzpicture}[inner sep=0]
\node[anchor=north west,inner sep=0] (bi) at (0,0) {\includegraphics[page=6,trim={52.00bp 282.50bp 39.00bp 521.50bp},clip,scale=0.840]{A_SoSe25.pdf}};
\fill[Marker,opacity=0.30,rounded corners=0.6bp] (242.00bp,-0.42bp) rectangle (269.39bp,-10.50bp);
\fill[Marker,opacity=0.30,rounded corners=0.6bp] (37.10bp,-22.26bp) rectangle (59.78bp,-32.34bp);
\fill[MarkerG,opacity=0.30,rounded corners=0.6bp] (336.25bp,-0.42bp) rectangle (363.64bp,-10.50bp);
\end{tikzpicture}
\par\nopagebreak
\vspace{0.8mm}\nopagebreak
\begin{loesung}{A5Farbe}
\vspace{1.2mm}
\includegraphics[page=5,trim={52.00bp 526.50bp 39.00bp 152.50bp},clip,scale=0.700]{L_SoSe25.pdf}
\par
\end{loesung}
\phantomsection\label{tf:SoSe25:5:3}
\Needspace*{18mm}
\verweis{A5Farbe}{3.}{\gs{Anzahl der Dipole} bzw. Dipol-Linien aus dem geforderten Gewinn}{gleiche Frage, nur andere Zahlenwerte \textendash{} siehe \textbf{SoSe 2021, Aufgabe 5.2}, \textbf{Seite \pageref{tf:SoSe21:5:2}}}
\phantomsection\label{tf:SoSe25:5:4}
\par\vspace{1.6mm}
\Needspace*{98mm}
\begin{tikzpicture}[inner sep=0]
\node[anchor=north west,inner sep=0] (bi) at (0,0) {\includegraphics[page=6,trim={52.00bp 111.50bp 39.00bp 666.50bp},clip,scale=0.840]{A_SoSe25.pdf}};
\fill[Marker,opacity=0.30,rounded corners=0.6bp] (257.64bp,-0.42bp) rectangle (291.00bp,-10.50bp);
\fill[MarkerG,opacity=0.30,rounded corners=0.6bp] (198.44bp,-22.26bp) rectangle (231.92bp,-32.34bp);
\end{tikzpicture}
\par\nopagebreak
\vspace{0.8mm}\nopagebreak
\begin{loesung}{A5Farbe}
\vspace{1.2mm}
\includegraphics[page=5,trim={52.00bp 360.50bp 39.00bp 438.50bp},clip,scale=0.700]{L_SoSe25.pdf}
\par
\par\vspace{1.0mm}
\includegraphics[page=5,trim={52.00bp 316.50bp 39.00bp 491.50bp},clip,scale=0.700]{L_SoSe25.pdf}
\par
\par\vspace{1.0mm}
\includegraphics[page=5,trim={52.00bp 148.50bp 39.00bp 535.50bp},clip,scale=0.700]{L_SoSe25.pdf}
\par
\par\vspace{1.0mm}
\includegraphics[page=5,trim={52.00bp 105.50bp 39.00bp 703.50bp},clip,scale=0.700]{L_SoSe25.pdf}
\par
\par\vspace{1.0mm}
\includegraphics[page=5,trim={52.00bp 80.50bp 39.00bp 747.50bp},clip,scale=0.700]{L_SoSe25.pdf}
\par
\end{loesung}
\par\vspace{6mm}
\end{document}